


\documentclass[sigconf,authordraft]{acmart}

\usepackage{hyperref}
\usepackage[capitalise]{cleveref}   % must be loaded AFTER hyperref
\usepackage{pgfplots}\pgfplotsset{compat=1.18}
\usepgfplotslibrary{groupplots}
\usepackage{siunitx}
\sisetup{
  group-minimum-digits = 5,
  group-separator = {,}
}
\usepackage{tikz}
\usetikzlibrary{arrows.meta, backgrounds, fit, positioning, calc, shapes}
\tikzset{ % UML "extends" arrow
	extends/.style={->, >={Triangle[open, width=0.25cm, length=0.25cm]}}
}
\usepackage{tabularx}
\usepackage{listings}
\usepackage{xcolor}
\usepackage{colortbl}
\usepackage{url}
\usepackage{subcaption}
\let\Bbbk\relax % avoid clash between acmart fonts and amssymb
\usepackage{amssymb}
\usepackage{pifont}
\usepackage[nointegrals]{wasysym}
\usepackage{booktabs}
\usepackage{CJKutf8}   % Japanese (Tokyo/PLATEAU) attribute examples; pdflatex-compatible, does not affect Latin text
\newcommand{\jp}[1]{\begin{CJK}{UTF8}{ipxm}#1\end{CJK}}
\providecommand{\pct}[1]{{\small\textcolor{black!100}{(#1\%)}}}   % coverage-share styling (gray) in the per-city content tables
\usepackage{array}

\lstdefinelanguage{Cypher}{
  keywords={MATCH, RETURN, WHERE, CREATE, DELETE, SET, MERGE, OPTIONAL, WITH, AS, AND, OR, NOT, LIMIT, ORDER, BY},
  sensitive=true,
  morecomment=[l]{//},
  morestring=[b]"
}

\lstset{
  language=Cypher,
  basicstyle=\ttfamily\small,
  keywordstyle=\bfseries,
  commentstyle=\itshape\color{gray},
  stringstyle=\color{teal},
  breaklines=true,
  columns=fullflexible,
  keepspaces=true,
  showstringspaces=false
}

\usepackage{xcolor}
\newcommand{\alert}[1]{\textcolor{red}{\textbf{#1}}}

%% LUKAS SACHEN
\usepackage{booktabs,amsmath}
\newcommand{\TODO}[1]{\textcolor{red}{\textbf{[TODO:} #1\textbf{]}}}
% \newcommand{\num}[1]{\textcolor{black}{\textbf{#1}}}
\newcommand{\figplaceholder}[2]{\setlength{\fboxrule}{1pt}%
\fcolorbox{red}{gray!8}{\parbox[c][#1][c]{0.94\linewidth}{\centering\footnotesize\color{red} #2}}}
\newbool{showcomments}

        % Show comments (Uncomment next line to show)
   \booltrue{showcomments}

        % Hide comments (Uncomment next line to hide)
  %       \boolfalse{showcomments}

    \newcommand{\pinaforecomment}[4]{%
    \ifbool{showcomments}{%
        \colorbox{#1}{\textcolor{#4}{\parbox{.8\linewidth}{#2: #3}}}%
    }{}%
}
\newcommand{\lukascomment}[1]{\pinaforecomment{blue}{Lukas}{#1}{white}}

%%
%% \BibTeX command to typeset BibTeX logo in the docs
\AtBeginDocument{%
  \providecommand\BibTeX{{%
    Bib\TeX}}}

%% Rights management information.
\setcopyright{acmlicensed}
\copyrightyear{2027}
\acmYear{2027}
\acmDOI{XXXXXXX.XXXXXXX}
%% These commands are for a PROCEEDINGS abstract or paper.
\acmConference[KDD '27]{The 33rd ACM SIGKDD Conference on Knowledge Discovery and Data Mining
V.1 (KDD '27), Month Day-Day, 2027, San Jose, CA, USA}{Month Day-Day, 2027}{San Jose, CA, USA}
\acmISBN{978-1-4503-XXXX-X/2027/XX}

% Custom commands
\newcommand{\mytoolname}{\textit{pykci}}
\newcommand{\mytoolsound}{\textit{pixie}}
\newcommand{\benchname}{\textit{Authenti\-City}}

% ---------------------------------------------------------------------------
% Draft-annotation macros (remove before submission)
% \addedcite: citation suggested/added by Claude -> yellow in PDF
\newcommand{\addedcite}[1]{\begingroup\setlength{\fboxsep}{1pt}\colorbox{yellow}{\cite{#1}}\endgroup}
% \verify: fact Claude could not verify -> red in PDF, must be checked
\newcommand{\verify}[1]{\begingroup\setlength{\fboxsep}{1pt}\colorbox{red!25}{#1}\endgroup}
% \dsq: a datasheet question -> flush-left italic lead-in with a uniform gap,
% so every question aligns regardless of position after a heading
\newcommand{\dsq}[1]{\par\smallskip\noindent\emph{#1}\space}
% benchtodo: orange box marking pending benchmark work
\newsavebox{\benchtodobox}
\newenvironment{benchtodo}
  {\par\medskip\noindent
   \begin{lrbox}{\benchtodobox}%
   \begin{minipage}{\dimexpr\linewidth-2\fboxsep-2\fboxrule\relax}%
   \textbf{[BENCHMARK TODO]}\ }
  {\end{minipage}\end{lrbox}%
   \noindent\fcolorbox{orange}{orange!15}{\usebox{\benchtodobox}}\par\medskip}
% ---------------------------------------------------------------------------

\begin{document}

\sloppy

\title{AuthentiCity: A Multi-Source Provenance-Aware Knowledge Graph and Benchmark for 3D City Models}
% Name locked 2026-07-05: AuthentiCity (authenticity <-> provenance + city pun;
% collision-checked, no dataset/benchmark uses it). pykci is credited as the
% tooling in the abstract/methods, deliberately not in the title.

\author{Huynh Duc An Son Nguyen}
\email{son.nguyen@hcu-hamburg.de}
\orcid{https://orcid.org/0000-0001-8711-1587}
\affiliation{%
  \institution{HafenCity University Hamburg, \\ Computational Methods Lab}
  \streetaddress{Henning-Voscherau-Platz 1}
  \city{Hamburg}
  \state{}
  \country{Germany}
  \postcode{20457}
}

\author{Lukas Arzoumanidis}
\email{lukas.arzou@hcu-hamburg.de}
\orcid{https://orcid.org/0000-0001-6668-1695}
\affiliation{%
  \institution{HafenCity University Hamburg, \\ Computational Methods Lab}
  \streetaddress{Henning-Voscherau-Platz 1}
  \city{Hamburg}
  \state{}
  \country{Germany}
  \postcode{20457}
}

\author{Youness Dehbi}
\email{youness.dehbi@hcu-hamburg.de}
\orcid{https://orcid.org/0000-0003-0133-4099}
\affiliation{%
  \institution{HafenCity University Hamburg, \\ Computational Methods Lab}
  \streetaddress{Henning-Voscherau-Platz 1}
  \city{Hamburg}
  \state{}
  \country{Germany}
  \postcode{20457}
}

\renewcommand{\shortauthors}{Nguyen et al.}

\begin{abstract}
Urban digital twins increasingly combine authoritative, crowd-sourced, machine-learned, and reconstructed data with differing reliability, coverage, and semantics. Yet few urban datasets provide a unified representation that supports multi-source integration, provenance tracking, spatial reasoning, and machine learning. As a result, existing benchmarks rarely evaluate reasoning about source origin, confidence, coverage, and agreement. We present \benchname{}, a multi-source, provenance-aware 3D city knowledge graph spanning five cities across three continents (Hamburg, Helsinki, Zurich, New York, and Tokyo) and comprising 180 GiB, 180M nodes, 220M edges, 1.2B properties, and 3.6M buildings. The labeled property graphs integrate authoritative CityGML features with OpenStreetMap data for all cities, adding roof-material predictions and reconstructed LoD3 geometry for Hamburg, under a provenance model in which derived information never replaces authoritative data. Confidence-weighted edges resolve many-to-many cross-source correspondences, constructing canonical urban entities while preserving traceable links to all contributing evidence. \benchname{} is primarily a data contribution. We introduce two benchmark families to demonstrate the tasks enabled by the representation. The first evaluates natural-language-to-query translation with tasks beyond conventional text-to-SQL and text-to-Cypher benchmarks, including 3D spatial reasoning, provenance-aware filtering, cross-source agreement and disagreement, coverage-aware aggregation, and infeasible-query detection. The second evaluates graph representation learning through multi-source attribute prediction, node classification, and cross-source matching prediction, facilitating comparison of provenance-agnostic and provenance-aware embeddings. Even a strong commercial LLM reaches only 54--69\,\% execution accuracy and a 7B open-weight model 6--19\,\%, and the open-weight model never abstains on an unanswerable question. We release the complete artifact under open licenses with an archival DOI: the five enriched property graphs, loaders, the question suite with gold queries and materialized answers, task splits, an evaluation harness, and a datasheet.
%\benchname{} includes benchmark questions, gold-standard queries, and an evaluation framework, providing a reproducible testbed for provenance-aware urban analytics, querying, and graph learning. 
% TL;DR: AuthentiCity is a city-scale, provenance-aware 3D knowledge graph dataset and benchmark that integrates heterogeneous urban data sources and introduces challenging tasks for natural-language querying and graph representation learning.
% TL;DR: We present AuthentiCity, a large-scale knowledge graph and benchmark for 3D city models that integrates multiple urban data sources into a unified, canonical representation while preserving data provenance.
%
% TODO Add graph embedding benchmark results once the final version is fixed.
% TODO Add quantitative summary of benchmark size and coverage (number of questions, number of cities, number of buildings, etc.) once the final version is fixed.
% TODO Empirical findings: Something like "current LLMs handle spatial reasoning reasonably but fail at provenance-aware filtering and infeasible-query detection, and provenance-aware embeddings outperform provenance-agnostic ones on cross-source matching by X"
\end{abstract}

\begin{CCSXML}
<ccs2012>
   <concept>
       <concept_id>10002951.10002952.10002953.10010146</concept_id>
       <concept_desc>Information systems~Graph-based database models</concept_desc>
       <concept_significance>500</concept_significance>
       </concept>
 </ccs2012>
\end{CCSXML}

\ccsdesc[500]{Information systems~Graph-based database models}

% \begin{CCSXML}
% <ccs2012>
%    <concept>
%        <concept_id>10002951.10002952.10002953.10010146</concept_id>
%        <concept_desc>Information systems~Graph-based database models</concept_desc>
%        <concept_significance>500</concept_significance>
%        </concept>
%    <concept>
%        <concept_id>10002951.10003227.10003236.10003237</concept_id>
%        <concept_desc>Information systems~Geographic information systems</concept_desc>
%        <concept_significance>500</concept_significance>
%        </concept>
%    <concept>
%        <concept_id>10010147.10010178.10010187</concept_id>
%        <concept_desc>Computing methodologies~Knowledge representation and reasoning</concept_desc>
%        <concept_significance>300</concept_significance>
%        </concept>
%  </ccs2012>
% \end{CCSXML}
%
% \ccsdesc[500]{Information systems~Graph-based database models}
% \ccsdesc[500]{Information systems~Geographic information systems}
% \ccsdesc[300]{Computing methodologies~Knowledge representation and reasoning}

% \begin{CCSXML}
% <ccs2012>
%    <concept>
%        <concept_id>10002951.10002952.10002953.10010146</concept_id>
%        <concept_desc>Information systems~Graph-based database models</concept_desc>
%        <concept_significance>500</concept_significance>
%        </concept>
%    <concept>
%        <concept_id>10002951.10003227.10003236.10003237</concept_id>
%        <concept_desc>Information systems~Geographic information systems</concept_desc>
%        <concept_significance>500</concept_significance>
%        </concept>
%    <concept>
%        <concept_id>10010147.10010257.10010258.10010259.10010263</concept_id>
%        <concept_desc>Computing methodologies~Learning latent representations</concept_desc>
%        <concept_significance>300</concept_significance>
%        </concept>
%  </ccs2012>
% \end{CCSXML}
%
% \ccsdesc[500]{Information systems~Graph-based database models}
% \ccsdesc[500]{Information systems~Geographic information systems}
% \ccsdesc[300]{Computing methodologies~Learning latent representations}

\keywords{Dataset, Benchmark, CityGML, OSM, Knowledge Graph}
% \keywords{Knowledge Graph, Provenance, CityGML, OpenStreetMap, LLM}
% \keywords{Dataset, Benchmark, Knowledge Graph, Provenance, Multi-Source, Spatial, CityGML, OpenStreetMap, LLM, Representation Learning}

\begin{teaserfigure}
  \includegraphics[width=\textwidth]{figures/Teaser_AuthentiCity_new.png} %Teaser_AuthentiCity_opaque.png %Teaser_AuthentiCity_3DTiles_Bldg_KG_OSM_Roofmats_LoD3.png %Teaser_AuthentiCity_TopView_3DTiles.png
  \caption{
  Hamburg, one of five \benchname{} datasets: a provenance-aware 3D knowledge graph integrating CityGML, OpenStreetMap, roof-material predictions (color-coded buildings), and reconstructed LoD3 geometry (orange subgraph). 
  %Multi-source, provenance-aware 3D knowledge graph of Hamburg, one of the five datasets in \benchname{}, generated from CityGML and enriched with OpenStreetMap features, machine-learned roof materials, and reconstructed LoD3 geometries using \mytoolname{}~\cite{nguyen_pykci_2026}. Buildings are colored by roof material; the reconstructed LoD3 building is highlighted in gold.
  % blue metal, green glass, graphite tar paper, sand concrete, red tiles
  }
  \Description{3D building models colored by predicted roof material, with one reconstructed LoD3 building highlighted in gold, overlaid with a spatially aligned knowledge graph of colored nodes and edges above an OpenStreetMap basemap of the HafenCity district in Hamburg.}
  \label{fig:teaser}
\end{teaserfigure}

\received{D M Y}
\received[revised]{D M Y}
\received[accepted]{D M Y}

\maketitle

\section{Introduction}\label{sec:introduction}

More than half of the world's population lives in cities~\cite{un_urban_2018}, and national mapping agencies increasingly provide authoritative semantic 3D city models at city, regional, and national scales~\cite{nyc_3dbuildingmodel_2016,plateau_3dcitymodel_2021,adv_buildings_2025}. Typically encoded in CityGML~\cite{groeger_citygml2_2012, kolbe_citygml3_2021}, these models provide virtual representations of urban environments and form a key component of urban digital twins (UDTs) for planning, simulation, and decision-making~\cite{ketzler_udtreview_2020, lei_udt_survey_2023}. Yet no single source provides a complete description of a city. Cadastral models offer surveyed geometry and official semantics but often lack use-level detail. Crowd-sourced maps such as OpenStreetMap (OSM) add names, addresses, points of interest, and street networks, but vary in coverage and quality~\cite{biljecki_qualitycrowsourced_2023}. Machine learning can infer attributes such as roof materials from orthophotos~\cite{lukas_roofmaterials, kanna2026semantic}, while photogrammetry provides detailed facade geometry. Integrating these sources therefore requires reasoning over facts that differ in origin, confidence, and coverage.
% More than half of the world's population lives in cities~\cite{un_urban_2018}, and national mapping agencies increasingly provide authoritative semantic 3D city models at city, regional, and national scales~\cite{nyc_3dbuildingmodel_2016,plateau_3dcitymodel_2021,adv_buildings_2025}. Typically encoded in CityGML~\cite{groeger_citygml2_2012, kolbe_citygml3_2021}, these models provide virtual representations of real-world urban environments and form a key component of urban digital twins (UDTs), which support planning, simulation, and decision-making processes~\cite{ketzler_udtreview_2020, lei_udt_survey_2023}. Yet no single source provides a complete or uniformly reliable description of a city. Cadastral models contain surveyed geometry and official semantics but often lack use-level detail. Crowd-sourced maps such as OpenStreetMap (OSM) add names, addresses, points of interest, and street networks, but their coverage and quality vary~\cite{biljecki_qualitycrowsourced_2023}. Machine learning can supply predicted attributes, such as roofing materials inferred from orthophotos~\cite{lukas_roofmaterials, kanna2026semantic}, but these attributes are probabilistic and limited to areas with suitable imagery. Photogrammetric reconstruction can further provide detailed facade geometry, usually for selected buildings. Integrating these sources therefore produces more than a heterogeneous data management problem: systems must reason over facts that differ in origin, confidence, and coverage.

Existing benchmarks do not directly evaluate this capability. Text-to-query benchmarks such as Spider~\cite{yu_spider_2018}, BIRD~\cite{li_bird_2023}, and CypherBench~\cite{feng_cypherbench_2025} focus on relational or encyclopedic data and do not combine spatial predicates, 3D city geometry, and source provenance. Urban benchmarks such as CityBench~\cite{feng_citybench_2025} and UUKG~\cite{ning_uukg_2023} evaluate urban reasoning or spatiotemporal prediction, but neither exposes a queryable semantic 3D city model nor represents differing trust levels across integrated sources. Consequently, current benchmarks cannot test whether systems select authoritative rather than predicted values, qualify uncertain evidence, or recognize unsupported aggregates arising from incomplete coverage. We refer to these capabilities as \emph{provenance-aware reasoning}.
% Existing benchmarks do not directly evaluate this capability. Text-to-query benchmarks such as Spider~\cite{yu_spider_2018}, BIRD~\cite{li_bird_2023}, and CypherBench~\cite{feng_cypherbench_2025} focus on relational or encyclopedic data and do not combine spatial predicates, 3D city geometry, and source provenance. Urban benchmarks such as CityBench~\cite{feng_citybench_2025} and UUKG~\cite{ning_uukg_2023} evaluate urban reasoning or spatiotemporal prediction, but they neither expose a queryable semantic 3D city model nor represent differing trust levels across integrated sources. Consequently, current resources cannot test whether a system selects authoritative rather than predicted values, qualifies answers supported only by uncertain evidence, or detects that an aggregate is unsupported because an attribute covers only part of the study area. We refer to these behaviors collectively as \emph{provenance-aware reasoning}.

We introduce \benchname{}, a multi-city knowledge graph (KG) and benchmark for evaluating provenance-aware querying and representation learning over heterogeneous 3D city data. The KG is designed around traceability: authoritative, crowd-sourced, predicted, and reconstructed information are represented as distinct entities and relations. External objects remain separate nodes, fusion edges record matching confidence, and source, confidence, and coverage become explicit inputs to benchmark evaluation.
% We introduce \benchname{}, a multi-city knowledge graph (KG) and benchmark for evaluating provenance-aware querying and representation learning over heterogeneous 3D city data. The KG is designed around traceability: authoritative, crowd-sourced, predicted, and reconstructed information is represented as traceable entities and relations. Derived values do not overwrite authoritative values, external objects remain separate nodes connected through traversal, and fusion edges record match confidence. This design makes source, confidence, and coverage explicit inputs to benchmark evaluation rather than undocumented properties of the underlying data.
%
\benchname{} is constructed with \mytoolname{}~\cite{nguyen_pykci_2026}, an open-source pipeline for mapping CityGML datasets to a compact labeled property graph (LPG) in Neo4j with an R-tree spatial index. On top of the authoritative layer, \benchname{} integrates OSM through confidence-weighted matching edges, attaches machine-learned roof-material predictions, and incorporates reconstructed LoD3 building models.
% \benchname{} is constructed with \mytoolname{}~\cite{nguyen_pykci_2026}, an open-source pipeline for mapping CityGML datasets to a compact labeled property graph (LPG) in Neo4j with an R-tree spatial index.
%% \footnote{Code is available at \url{https://github.com/hcu-cml/pykci}. The mapping pipeline is described in the companion system paper~\cite{nguyen_pykci_2026}; this paper contributes the multi-source, provenance-aware dataset and benchmark.} 
% On top of the authoritative layer, \benchname{} integrates OSM through explicit, confidence-weighted matching edges that represent 1:1, 1:n, n:1, and n:m correspondences between crowd-sourced and cadastral building footprints. It further attaches machine-learned roof material predictions and incorporates reconstructed Level of Detail 3 (LoD3) building models.

The dataset follows a two-tier design that balances cross-city comparability with source depth. Tier~1 applies the same CityGML-plus-OSM construction to all cities, enabling a comparable evaluation core. Tier~2 provides a deep-fusion instance for Hamburg, where all four sources are available: CityGML, OSM, ML predictions, and LoD3 reconstructions. The dataset also preserves incomplete coverage rather than masking it. For example, roof-material predictions are available for only about half of Hamburg's buildings, enabling evaluation of whether models distinguish unavailable facts from negative facts and avoid unsupported city-wide conclusions.
% The dataset follows a two-tier design that balances cross-city comparability with source depth. Tier~1 applies the same CityGML-plus-OSM construction to all selected cities, enabling a comparable core for evaluation. Tier~2 provides a deep-fusion instance for Hamburg, Germany, where all four source layers are available: CityGML, OSM, ML predictions, and LoD3 reconstructions. The dataset also preserves incomplete coverage rather than masking it. For example, ML-predicted roof materials are available for approximately half of buildings in Hamburg. This partial coverage enables direct evaluation of whether a model distinguishes unavailable facts from negative facts and avoids unsupported city-wide conclusions.

We define two complementary benchmark task families. The first evaluates natural-language-to-query translation with gold Cypher over the property graph, covering spatial predicates, cross-source agreement and disagreement, provenance-filtered retrieval, coverage-aware aggregation, and infeasible-question detection (\Cref{subsec:query_tasks}). The second evaluates graph representation learning through attribute imputation, building-function classification, and matching-link prediction, comparing provenance-agnostic and provenance-weighted embeddings. Together, these tasks evaluate provenance awareness at both the symbolic and representation-learning levels.
% On this dataset we define two complementary benchmark task families. The first evaluates natural-language-to-query translation with gold Cypher over the property graph, enabling provenance-aware categories grounded in fusion structures that a relational model cannot naturally express (\Cref{subsec:query_tasks}): spatial predicates over indexed 3D geometry, cross-source agreement and disagreement, provenance-filtered retrieval, coverage-aware aggregation, and infeasible-question detection. \textcolor{blue}{The second evaluates graph representation learning through attribute imputation against held-out ML predictions, building-function classification, and matching-link prediction, comparing provenance-agnostic and provenance-weighted embeddings.} Together the two families test provenance at both the symbolic and the learned-representation level.

Our contributions are: 
\begin{itemize} 
  \item \textbf{A provenance-preserving, multi-source 3D city dataset.} We release five two-tier city KGs with 180~GiB of data, 180\,million nodes, 220\,million edges, 1.2\,billion properties, and 3.6\,million buildings. The KG integrates authoritative CityGML data, crowd-sourced OSM data, ML-predicted attributes, and reconstructed LoD3 geometry while preserving source provenance through canonical feature nodes and source-specific attachments. The release includes loaders, benchmark splits, a datasheet~\cite{gebru_datasheets_2021}, and an archival DOI. 
  \item \textbf{A benchmark for provenance-aware querying and graph learning.} We provide a text-to-query suite covering spatial, cross-source, coverage-aware, and infeasibility cases, together with a representation-learning suite for attribute imputation, building classification, and matching-link prediction. 
  \item \textbf{Baselines and diagnostic analysis.} 
  We report reference results and a failure decomposition that separates Cypher-generation errors from semantic errors and from failures to abstain.
  %We report reference results and category-level failure analysis, separating errors in spatial reasoning, cross-source reasoning, coverage awareness, and infeasibility detection. 
  For graph learning, we compare provenance-agnostic and provenance-weighted variants to measure the value of explicitly representing source information. 
\end{itemize}
% Our contributions are: 
% \begin{itemize} 
%   \item \textbf{A provenance-preserving, multi-source 3D city dataset.} We release five two-tier city KGs with 180~GiB of data, 180\,million nodes, 220\,million edges, 1.2\,billion properties, and 3.6\,million buildings. The KG integrates authoritative CityGML data, crowd-sourced OSM data, ML-predicted attributes, and reconstructed LoD3 geometry while preserving source identities through canonical feature nodes and source-specific attachments. The release includes loaders, benchmark splits, a \verify{datasheet}~\cite{gebru_datasheets_2021}, and an archival DOI.
%   \item \textbf{A benchmark for provenance-aware querying and graph learning.} We provide a text-to-query suite covering spatial, cross-source, coverage-aware, and infeasibility cases, together with a representation learning suite for \textcolor{blue}{attribute imputation, building classification, and matching-link prediction}. 
%   \item \textbf{Baselines and diagnostic analysis.} We report reference results together with category-level failure analysis, separating errors in spatial reasoning, cross-source reasoning, coverage awareness, and infeasibility detection. \textcolor{blue}{For the graph learning tasks, we compare provenance-agnostic and provenance-weighted variants to measure the effect of explicitly representing source information.}
% \end{itemize}
% TODO Publish loaders, benchmark splits, and datasheet with DOI.

\section{Related Work}\label{sec:related_work}

\subsection{Natural-Language-to-Query Benchmarks}

Text-to-SQL is the most mature setting. Spider~\cite{yu_spider_2018} established cross-domain evaluation (\num{10181} questions over 200 databases) but its schemas are small and carry no spatial types. BIRD~\cite{li_bird_2023} scaled to \num{12751} questions over 95 large, noisy databases and introduced the execution-centric evaluation %and efficiency scoring 
we adopt, yet still contains no geospatial reasoning and treats every value as equally trustworthy. Spider~2.0~\cite{li_spider2_2025} adds enterprise-scale realism but remains relational. NL2SQL-BUGs~\cite{liu_nl2sqlbugs_2025} shifts focus to detecting semantically incorrect SQL, which motivates our infeasibility category, and Dr.Spider~\cite{chang_drspider_2023} stresses robustness under perturbations; neither addresses spatial reasoning, provenance, or source disagreement.

For property graphs, the Neo4j Text2Cypher dataset~\cite{ozsoy_text2cypher_2025} aggregates about \num{44000} instances, though many are not grounded in an executable graph. CypherBench~\cite{feng_cypherbench_2025} provides 11 Wikidata-derived graphs with over \num{10000} questions and the execution-accuracy metric we adopt, but its graphs are encyclopedic, with no geometry. Mind the Query~\cite{chauhan_mindthequery_2025} contributes more than \num{27000} validated text-to-Cypher pairs with a rigorous validation pipeline we take as a quality template. All share two limits relevant here: their schemas contain neither spatial geometry nor multiple sources describing the same entity, so the capabilities \benchname{} targets are outside their scope.

\subsection{Urban Benchmarks and Knowledge Graphs}

CityBench~\cite{feng_citybench_2025} evaluates LLMs and VLMs on eight urban tasks across 13 cities but exposes no queryable semantic 3D city model. CityGPT~\cite{feng_citygpt_2025} embeds urban knowledge into the model itself, which does not generalize across cities or updates and gives no auditable grounding. UUKG~\cite{ning_uukg_2023} releases unified urban KGs for spatiotemporal prediction and UrbanKGent~\cite{ning_urbankgent_2024} automates KG construction with LLM agents, but both operate on POI- and region-level entities rather than 3D building models, and neither provides queryable provenance. Surveys of urban KGs and digital twins~\cite{liu_urbankg_2023, wang_urbankgreview_2024, saifwajid_kgudt_2024, akroyd_kgudt_2021} consistently name data fusion as a primary motivation, yet provenance is rarely a first-class queryable component.

Closest to our sources,~\citet{ding_osmgraph_2025} integrate CityGML and OSM into an RDF KG queryable with GeoSPARQL, but the integration is ontology-mediated and not released as a benchmark with tasks, splits, and baselines. KCityChatBot~\cite{liu_kcitychatbot_2025} pairs a CityGML KG with a multi-agent LLM pipeline but provides no reusable benchmark. On the systems side, 3DCityDB~\cite{yao_3dcitydb_2018, yao_3dcitydb5_2025} is the most widely adopted CityGML platform, and Semantic Web representations have been studied extensively~\cite{chadzynski_sem3dcitydb_2021}, but neither natively supports confidence-weighted cross-source correspondences or fact-level provenance, the capabilities central to our tasks (\Cref{subsec:query_tasks}). \benchname{} is complementary to these systems.

\subsection{Multi-Source Integration and Matching}

Fusing crowd-sourced and authoritative geodata requires entity resolution across polygon datasets. Optimal many-to-many polygon matching under the Jaccard measure is NP-hard~\cite{naumann_nnmatching_2024}, with scalable formulations building on tree-constrained bipartite matching~\cite{canzar_treematching_2011, naumann_nnmatching_2025}. Rather than commit to a single set of hard matches, \benchname{} retains the full weighted overlap graph as first-class confidence-annotated edges, letting queries and learning methods resolve correspondences as needed (\Cref{sec:dataset}); the released correspondences are also a resource for polygon-matching research, where ground-truth labels remain scarce.

\subsection{Positioning}

\Cref{tab:gap_matrix} summarizes the gap in existing datasets and benchmarks. Each row represents a strong benchmark within a well-established research area, yet, to the best of our knowledge, no existing benchmark combines these dimensions. \benchname{} addresses this gap. To facilitate comparison with the closest prior work, we adopt evaluation metrics compatible with CypherBench and BIRD, particularly the execution accuracy.
% and structural match.

% \begin{table*}[h!]
%   \caption{\benchname{} and related datasets and benchmarks. \benchname{}
%   trades question volume for graph scale, executable provenance-grounded
%   gold, and capability coverage.}
%   \label{tab:gap_matrix}
%   \begin{tabular}{l r l c ccccc}
%     \toprule
%     & \multicolumn{3}{c}{Resource} & \multicolumn{5}{c}{Capabilities} \\
%     \cmidrule(lr){2-4}\cmidrule(lr){5-9}
%     Benchmark & \#Q & Target data & Gold & Spatial & Multi & Prov & Infeas & RL \\
%     \midrule
%     Spider~\cite{yu_spider_2018}                & 10{,}181 & 200 DBs, \SI{<1}{\giga\byte}   & \ding{51} & \ding{55} & \ding{55} & \ding{55} & \ding{55} & \ding{55} \\
%     BIRD~\cite{li_bird_2023}                    & 12{,}751 & 95 DBs, \SI{33}{\giga\byte}    & \ding{51} & \ding{55} & \ding{55} & \ding{55} & \ding{55} & \ding{55} \\
%     Spider 2.0~\cite{lei_spider2_2025}          & ??       & ?? enterprise DBs              & \ding{51} & \ding{55} & \ding{55} & \ding{55} & \ding{55} & \ding{55} \\
%     NL2SQL-BUGs~\cite{liu_nl2sqlbugs_2025}      & ??       & ?? DBs                         & \ding{51} & \ding{55} & \ding{55} & (\ding{51}) & \ding{55} & \ding{55} \\
%     Text2Cypher~\cite{ozsoy_text2cypher_2025}   & 44{,}387 & aggregated, largely ungrounded  & (\ding{51}) & \ding{55} & \ding{55} & \ding{55} & \ding{55} & \ding{55} \\
%     CypherBench~\cite{feng_cypherbench_2025}    & 11{,}011 & 11 graphs, ?? nodes            & \ding{51} & \ding{55} & \ding{55} & \ding{55} & \ding{55} & \ding{55} \\
%     Mind the Query~\cite{chauhan_mindthequery_2025} & 27{,}000 & ?? graphs                  & \ding{51} & \ding{55} & \ding{55} & \ding{55} & \ding{55} & \ding{55} \\
%     \midrule
%     CityBench~\cite{feng_citybench_2025}        & ??       & 13 cities, no queryable KG     & ---       & (\ding{51}) & \ding{55} & \ding{55} & \ding{55} & \ding{55} \\
%     UUKG~\cite{ning_uukg_2023}                  & ---      & ?? cities, POI/region KG       & ---       & (\ding{51}) & \ding{55} & \ding{55} & \ding{55} & \ding{51} \\
%     \citet{ding_osmgraph_2025}    & ---      & 1 city, RDF (CityGML+OSM)      & ---       & \ding{51} & \ding{51} & \ding{55} & \ding{55} & \ding{55} \\
%     \midrule
%     \benchname{} & 1{,}394 & 5 city LPGs, 180M nodes & \ding{51} & \ding{51} & \ding{51} & \ding{51} & \ding{51} \\
%     \bottomrule
%   \end{tabular}
  
%   \smallskip\raggedright\footnotesize
%   \#Q: released natural-language questions; --- marks resources that ship no
%   question suite. Target data: number and scale of the queried databases or
%   graphs. Gold: gold queries execute against the released data and are
%   verified ((\ding{51}) partial: many instances are not grounded in an
%   executable graph). Spatial: predicates over 2D/3D geometries. Multi:
%   multiple sources adding to the same entities. Prov: explicit provenance
%   representation and provenance-aware evaluation. Infeas: deliberately
%   unanswerable queries. RL: representation learning tasks. Parenthesized
%   (\ding{51}) marks partial or narrower support. \benchname{} gold queries
%   additionally pass five machine-checked gates (Appendix~\ref{sec:gates})
%   and a \SI{20}{\percent} stratified human review.
% \end{table*}

\begin{table}[t]
  \caption{\benchname{} and related datasets and benchmarks. The comparison is
  on capability coverage; see the note below on suite size.}
  \label{tab:gap_matrix}
  \begin{tabular}{lccccc}
    \toprule
    Resource & Spatial & Multi & Prov & Infeas & RL \\
    \midrule
    \multicolumn{6}{@{}l}{\textit{Text-to-query benchmarks}}\\
    \quad Spider~\cite{yu_spider_2018} & \ding{55} & \ding{55} & \ding{55} & \ding{55} & \ding{55} \\
    \quad BIRD~\cite{li_bird_2023} & \ding{55} & \ding{55} & \ding{55} & \ding{55} & \ding{55} \\
    \quad NL2SQL-BUGs~\cite{liu_nl2sqlbugs_2025} & \ding{55} & \ding{55} & \ding{55} & (\ding{51}) & \ding{55} \\
    \quad Text2Cypher~\cite{ozsoy_text2cypher_2025} & \ding{55} & \ding{55} & \ding{55} & \ding{55} & \ding{55} \\
    \quad CypherBench~\cite{feng_cypherbench_2025} & \ding{55} & \ding{55} & \ding{55} & \ding{55} & \ding{55} \\
    \quad Mind the Query~\cite{chauhan_mindthequery_2025} & \ding{55} & \ding{55} & \ding{55} & \ding{55} & \ding{55} \\
    \addlinespace[3pt]
    \multicolumn{6}{@{}l}{\textit{Urban data resources and benchmarks}}\\
    \quad CityBench~\cite{feng_citybench_2025} & (\ding{51}) & \ding{55} & \ding{55} & \ding{55} & \ding{55} \\
    \quad UUKG~\cite{ning_uukg_2023} & (\ding{51}) & \ding{55} & \ding{55} & \ding{55} & \ding{51} \\
    \quad \citet{ding_osmgraph_2025} & \ding{51} & \ding{51} & \ding{55} & \ding{55} & \ding{55} \\
    \midrule
    \benchname{} (+\mytoolname{}~\cite{nguyen_pykci_2026}) & \ding{51} & \ding{51} & \ding{51} & \ding{51} & \ding{51} \\
    \bottomrule
  \end{tabular}

  \smallskip\raggedright\footnotesize
  Spatial: predicates over 2D/3D geometries. Multi: multiple sources adding to the same entities. Prov: explicit provenance representation and provenance-aware evaluation. Infeas: deliberately unanswerable queries. RL: representation learning tasks. Parenthesized (\ding{51}) marks partial or narrower support. 
  % \citet{ding_osmgraph_2025} integrate CityGML and OpenStreetMap into an RDF knowledge graph but release no task suite, splits, or baselines. 
  The table compares \textbf{capability coverage}, not suite size: \benchname{} releases 1{,}394 executable, gold-verified questions, against 10{,}181 for Spider and more than 10{,}000 for CypherBench over far smaller, non-geometric graphs.
\end{table}

% \begin{table}[h!]
%   \caption{\benchname{} and related datasets and benchmarks.}
%   \label{tab:gap_matrix}
%   \begin{tabular}{lccccc}
%     \toprule
%     Benchmark & Spatial & Multi & Prov & Infeas & RL \\
%     \midrule
%     Spider~\cite{yu_spider_2018} & \ding{55} & \ding{55} & \ding{55} & \ding{55} & \ding{55} \\
%     BIRD~\cite{li_bird_2023} & \ding{55} & \ding{55} & \ding{55} & \ding{55} & \ding{55} \\
%     NL2SQL-BUGs~\cite{liu_nl2sqlbugs_2025} & \ding{55} & \ding{55} & \ding{55} & (\ding{51}) & \ding{55} \\
%     Text2Cypher~\cite{ozsoy_text2cypher_2025} & \ding{55} & \ding{55} & \ding{55} & \ding{55} & \ding{55} \\
%     CypherBench~\cite{feng_cypherbench_2025} & \ding{55} & \ding{55} & \ding{55} & \ding{55} & \ding{55} \\
%     Mind the Query~\cite{chauhan_mindthequery_2025} & \ding{55} & \ding{55} & \ding{55} & \ding{55} & \ding{55} \\
%     CityBench~\cite{feng_citybench_2025} & (\ding{51}) & \ding{55} & \ding{55} & \ding{55} & \ding{55} \\
%     UUKG~\cite{ning_uukg_2023} & (\ding{51}) & \ding{55} & \ding{55} & \ding{55} & \ding{51} \\
%     \midrule
%     \benchname{} & \ding{51} & \ding{51} & \ding{51} & \ding{51} & \ding{51} \\
%     \bottomrule
%   \end{tabular}
  
%   \smallskip\raggedright\footnotesize
%   Spatial: predicates over 2D/3D geometries. Multi: multiple sources adding to the same entities. Prov: explicit provenance representation and provenance-aware evaluation. Infeas: deliberately unanswerable queries. RL: representation learning tasks. Parenthesized (\ding{51}) marks partial or narrower support.
% \end{table}

\section{The AuthentiCity Dataset}\label{sec:dataset}

\begin{center}
\fbox{%
  \begin{minipage}{0.93\columnwidth}
    \small\raggedright
    \textbf{The \benchname{} artifact.}\enspace
    Five city-scale knowledge graphs: 180\,GiB, 180M nodes, 220M edges,
    1.2B properties, and 3.6M buildings.
    \par\smallskip
    \textbf{Archive.}\enspace DOI \texttt{10.5281/zenodo.21547211}
    \par\smallskip
    \textbf{Code.}\enspace \url{https://github.com/hcu-cml/pykci/authenticity}
    \par\smallskip
    \textbf{Contents.}\enspace Neo4j dump and backend-neutral node/edge
    export, loaders, question suite with gold queries and materialized
    answers, task splits, evaluation harness, and datasheet
    (Appendix~\ref{app:datasheet}).
    \par\smallskip
    \textbf{Licenses.}\enspace Per source (Table~\ref{tab:city_sources});
    OpenStreetMap under ODbL~\cite{opendatacommons_odbl_2026}; our annotations and question
    suite under CC\,BY\,4.0.
  \end{minipage}%
}
\end{center}

This section introduces \benchname{}'s five datasets and its approach to multi-source integration and provenance management. An overview is provided in~\Cref{fig:overview,tab:city_summary}.

% ===================== Figure 2: AuthentiCity at a glance =====================
\newsavebox{\overviewfigbox}
\begin{figure*}
\centering
\savebox{\overviewfigbox}{%
% per-source-layer accent colors (used consistently in all three panels)
\colorlet{cAuth}{green!45!black}%
\colorlet{cOsm}{blue!55!black}%
\colorlet{cPred}{orange!75!black}%
\colorlet{cRecon}{violet!65!black}%
% availability symbols (wasysym circles, per source-layer color)
\newcommand{\avF}[1]{\textcolor{#1}{\scriptsize\CIRCLE}}%
\newcommand{\avH}[1]{\textcolor{#1}{\scriptsize\LEFTcircle}}%
\newcommand{\avS}[1]{\textcolor{#1}{\scriptsize\Circle}}%
\newcommand{\avN}{\textcolor{gray!60}{--}}%
% task-family badge
\newcommand{\famb}[1]{{\setlength{\fboxsep}{1.6pt}\colorbox{black!75}{\textcolor{white}{\scriptsize\bfseries #1}}}}%
\begin{tikzpicture}[
  font=\small,
  nohyph/.style={execute at begin node={\hyphenpenalty=10000\relax}},
  layer/.style={nohyph, rounded corners=2pt, draw, line width=0.8pt, align=left, text width=2.75cm, inner sep=4pt},
  gnode/.style={rounded corners=2pt, draw, line width=0.8pt, align=center, inner sep=3.5pt},
  card/.style={nohyph, rounded corners=2pt, draw=gray!50, line width=0.5pt, fill=white, align=left, text width=6.5cm, inner sep=4.5pt},
  task/.style={nohyph, rounded corners=2pt, draw=gray!50, line width=0.5pt, fill=white, align=left, text width=4.15cm, inner sep=4.5pt},
  flow/.style={-{Triangle[width=6pt,length=5.5pt]}, line width=1.4pt, gray!55},
  edgeflow/.style={-{Triangle[width=4pt,length=4pt]}, semithick},
  panel/.style={rounded corners=4pt, draw=gray!45, line width=0.6pt, fill=gray!3, inner sep=5.5pt},
  ptitle/.style={anchor=south west, font=\small\bfseries, inner sep=2pt, text=black!75},
  flowlbl/.style={font=\small\itshape, text=black!60, inner sep=1pt}
]
% ---------------- Panel B content: unified graph per city (center) ----------------
% graph motif: Building --ENRICHED_BY--> OsmFeature
\node[gnode, draw=cAuth, fill=green!10] (bldg) at (0,0)
  {\textbf{Building}\\[1.5pt]
   {\scriptsize\ttfamily measured\_height 18.9}\\[1pt]
   {\scriptsize\ttfamily\color{cPred}predicted\_roof\_material}};
% invariants card
\node[card, below=6.5mm of bldg.south west, anchor=north west] (rules)
  {\textbf{Graph invariants}\\[2.5pt] 
  \textbf{\color{black!70}Source-preserving representation}: each external object retains its original identifier and is linked to an authoritative entity rather than merged with it\\[2pt] 
  \textbf{\color{black!70}Non-destructive attribute integration}: derived attributes use source-specific prefixes (\texttt{osm\_*}, \texttt{predicted\_*}) and coexist with authoritative attributes\\[2pt] 
  \textbf{\color{black!70}Confidence and coverage tracking}: predictions and fusion edges carry confidence values in $[0,1]$, and coverage is recorded explicitly, ensuring that missing values are distinguishable from negative evidence\\[2pt]
  \textbf{\color{black!70}Native-language schemas}: attribute names and code lists are preserved in their original natural language; the LLM query layer bridges languages and schema variations};
\node[card, below=5.5mm of rules.south west, anchor=north west] (av)
  {\textbf{Source-layer availability per city}\\[2.5pt]
   {\small\begin{tabular}{@{}l@{\hspace{8pt}}c@{\hspace{8pt}}c@{\hspace{8pt}}c@{\hspace{8pt}}c@{\hspace{8pt}}c@{}}
     & \textbf{HAM} & \textbf{ZRH} & \textbf{HEL} & \textbf{TYO} & \textbf{NYC}\\[2pt]
   \textcolor{cAuth}{Authoritative} & \avF{cAuth} & \avF{cAuth} & \avF{cAuth} & \avF{cAuth} & \avF{cAuth}\\[1.5pt]
   \textcolor{cOsm}{Crowd-sourced} & \avF{cOsm} & \avF{cOsm} & \avF{cOsm} & \avF{cOsm} & \avF{cOsm}\\[1.5pt]
   \textcolor{cPred}{ML-predicted} & \avH{cPred} & \avN & \avN & \avN & \avN\\[1.5pt]
   \textcolor{cRecon}{Reconstructed} & \avS{cRecon} & \avN & \avN & \avN & \avN\\
   \end{tabular}}\\[3pt]
   {\scriptsize \avF{black!55} all buildings \; \avH{black!55} $\sim$50\,\% \; \avS{black!55} selected buildings \; \avN\, not available}\\[1.5pt]
   {\scriptsize\itshape\color{black!60} 180~GiB $\cdot$ 180M nodes $\cdot$ 220M edges $\cdot$ 1.2B properties $\cdot$ 3.6M buildings}};
\node[gnode, draw=cOsm, fill=blue!8, anchor=east] (osmf) at (rules.east |- bldg)
  {\textbf{OsmFeature}\\[1.5pt]
   {\scriptsize\ttfamily osm\_height 20}};
\draw[edgeflow, cOsm] (bldg) -- node[above=2pt, font=\small\ttfamily]{ENRICHED\_BY}
    node[below=2pt, font=\scriptsize\ttfamily]{\{jaccard, is\_primary\}} (osmf);

% ---------------- Panel C content: benchmark (right) ----------------
\node[task, anchor=north west] (famA) at ($(rules.east |- bldg.north)+(12.5mm,0)$)
  {\famb{A}\; \textbf{Text-to-query}\hspace{3pt}{\scriptsize\color{black!55}gold Cypher}\\[2.5pt]
   aggregation $\cdot$ filtering\,/\,top-k $\cdot$ multi-hop traversal\\[1.5pt]
   {\color{cAuth}\textbf{+}} spatial (R-tree predicates)\\[1pt]
   {\color{cAuth}\textbf{+}} cross-source agreement\\[1pt]
   {\color{cAuth}\textbf{+}} provenance-aware queries\\[1pt]
   {\color{cAuth}\textbf{+}} coverage-aware queries\\[1pt]
   {\color{cAuth}\textbf{+}} infeasible queries (must abstain)\\[2.5pt]
   {\color{gray!40}\rule{4.15cm}{0.4pt}}\\[1.5pt]
   {\scriptsize\itshape\color{black!60} execution~accuracy %$\cdot$ structural~match 
   $\cdot$ infeasibility~precision/recall %$\cdot$ latency
   }};
\node[task, below=3mm of famA] (famB)
  {\famb{B}\; \textbf{Representation learning}\\[2.5pt]
   height \& roof-material prediction\\[1pt]
   roof-type \& function classification\\[1pt]
   \texttt{ENRICHED\_BY} link prediction\\[2.5pt]
   {\color{gray!40}\rule{4.15cm}{0.4pt}}\\[1.5pt]
   {\scriptsize\itshape\color{black!60} macro-F1\,/\,$R^2$ on held-out targets $\cdot$ ROC-AUC $\cdot$ provenance-agnostic vs.\ -aware $\Delta$}};
\node[task, below=3mm of famB] (verif) 
  {\textbf{Dataset verification}\\[2.5pt] 
  {\color{cAuth}\textbullet} executable, deterministic gold answers\\[1pt] 
  {\color{cAuth}\textbullet} infeasible questions proven unanswerable\\[1pt] 
  {\color{cAuth}\textbullet} template slots grounded in the released graph\\[1pt] 
  {\color{cAuth}\textbullet} lossless round-trip for elements, attributes, and coordinates};
% \node[task, below=3mm of famB] (verif) 
%   {\textbf{Dataset Verification}\\[2.5pt] % \ding{51} for the following four items
%   {\color{cAuth}\textbullet} executable, deterministic gold answers\\[1pt] %{\color{cAuth}\ding{51}} text\\[1pt] 
%   {\color{cAuth}\textbullet} each infeasible question checked against the graph to confirm no valid answer exists\\[1pt]
%   {\color{cAuth}\textbullet} question templates filled only with values verified to exist in the released graph\\[1pt]
%   {\color{cAuth}\textbullet} source-to-graph-to-export round trip verified lossless at the element, attribute, and coordinate level};

% ---------------- Panel A content: source layers (left) ----------------
% chips vertically centered on the shared panel height
\coordinate (bmid) at ($(bldg.north)!0.5!(verif.south)$);
\node[layer, draw=cOsm, fill=blue!8, anchor=south] (Losm) at ($(bldg.west |- bmid)+(-26mm,4.5mm)$)
  {\textcolor{cOsm}{\textbf{Crowd-sourced}}\\[1pt]
   {\scriptsize\color{black!75}OpenStreetMap: tags, POIs, streets, addresses}};
\node[layer, draw=cAuth, fill=green!10, above=9mm of Losm] (Lauth)
  {\textcolor{cAuth}{\textbf{Authoritative}}\\[1pt]
   {\scriptsize\color{black!75}CityGML 2.0: surveyed geometry, official semantics}};
\node[layer, draw=cPred, fill=orange!13, anchor=north] (Lpred) at ($(bldg.west |- bmid)+(-26mm,-4.5mm)$)
  {\textcolor{cPred}{\textbf{ML-predicted}}\\[1pt]
   {\scriptsize\color{black!75}roof materials, $\sim$50\,\% coverage, confidence per prediction}};
\node[layer, draw=cRecon, fill=violet!8, below=9mm of Lpred] (Lrecon)
  {\textcolor{cRecon}{\textbf{Reconstructed}}\\[1pt]
   {\scriptsize\color{black!75}LoD3 openings from facade images}};

% ---------------- Panels (equal heights via shared y-extremes) ----------------
\begin{scope}[on background layer]
  \node[panel, fit={(Lauth)(Lrecon)($(Lauth.west)+(-6mm,0)$)(Lauth.west |- bldg.north)(Lauth.west |- av.south)(Lauth.west |- verif.south)},
        label={[ptitle]north west:{Source layers}}] (panelA) {};
  \node[panel, fit={(bldg)(osmf)(rules)(av)(bldg.west |- verif.south)},
        label={[ptitle]north west:{One provenance-aware graph per city}}] (panelB) {};
  \node[panel, fit={(famA)(famB)(verif)(famA.west |- av.south)},
        label={[ptitle]north west:{Benchmark (two task families)}}] (panelC) {};
\end{scope}

% trust axis inside panel A
\draw[-{Triangle[width=4.5pt,length=4.5pt]}, gray!65, line width=0.9pt, line cap=round]
  ([xshift=3mm,yshift=3mm]panelA.south west) -- ([xshift=3mm,yshift=-3mm]panelA.north west)
  node[midway, rotate=90, below=0.5pt, font=\small\itshape, text=black!55]{increasing trust};

% flows between panels
\draw[flow] (panelA.east) -- (panelB.west) node[flowlbl, midway, above=1.5pt]{fuse};
\draw[flow] (panelB.east) -- (panelC.west) node[flowlbl, midway, above=1.5pt]{grounds};
\end{tikzpicture}}%
\typeout{OVERVIEWFIG WIDTH: \the\wd\overviewfigbox\space TEXTWIDTH: \the\textwidth}%
\usebox{\overviewfigbox}
\caption{
  Overview of \benchname{}. Four source layers with different trust levels (left) are integrated into a provenance-aware LPG for each city (center). 
  % The graph preserves source identities, authoritative values, confidence scores, and native-language schemas. 
  % The availability matrix summarizes layer coverage across cities. 
  The resulting graphs support two benchmark families with task-specific metrics and dataset-level validation (right). Categories marked (+) are absent from existing text-to-query benchmarks.
}
\Description{Three-panel diagram: four source layers ordered by increasing trust on the left. The fused provenance-aware city graph in the center with a Building to OsmFeature enrichment example, four graph invariants, and a per-city source-layer availability matrix. The two benchmark task families with metrics and dataset verification on the right.}
\label{fig:overview}
\end{figure*}

\begin{table*}[h!]
  \caption{
    Overview of \benchname{}'s five city-scale datasets (full statistics in \Cref{app:dataset_stats}, \Cref{tab:city_sources,tab:city_graph_content}). Hamburg is the Tier~2 deep-fusion instance; the remaining cities form the Tier~1 comparable core. 
    % The five \benchname{} cities at a glance (headline figures; full per-city breakdowns in \Cref{app:dataset_stats}: \Cref{tab:city_sources,tab:city_graph_content,tab:city_enrichment,tab:city_storage}). \emph{Buildings} are lossless graph nodes after identity resolution (\Cref{subsec:fidelity}). \emph{Nodes/bldg.} is the median graph nodes per building, a detail-density measure. \emph{Keys} counts distinct source thematic property keys (the multilingual schema axis). \emph{OSM-enr.} is the share of cadastral buildings with at least one OSM match. \emph{Store}/\emph{Dump} are the full Neo4j store and the released compressed archive. Hamburg is the Tier~2 deep-fusion flagship (all four source layers); the others form the Tier~1 comparable core. $^\dagger$\,reprojected at ingest from a non-metric source CRS.
    }
  \label{tab:city_summary}
  %\footnotesize
  \setlength{\tabcolsep}{5pt}
  \begin{tabular}{l rrrrr r}
    \toprule
    & Hamburg & Helsinki & Zurich & New York & Tokyo & All 5 \\
    & \small DE & \small FI & \small CH & \small US & \small JP & \small corpus \\
    \midrule
    CityGML / LoD        & 2.0 / LoD2 & 2.0 / LoD2 & 2.0 / LoD2.3 & 1.0/2.0 / LoD1--2 & 2.0 / LoD1--3 & -- \\
    Metric CRS           & EPSG:25832 & EPSG:3879 & EPSG:2056 & EPSG:32618$^\dagger$ & EPSG:6677$^\dagger$ & -- \\
    Buildings            & \num{388267} & \num{2980} & \num{102668} & \num{1083437} & \num{2005762} & \num{3583114} \\
    Nodes                & 17.44\,M & 0.41\,M & 41.51\,M & 45.80\,M & 74.41\,M & 179.57\,M \\
    Edges                & 26.30\,M & 0.62\,M & 45.55\,M & 54.21\,M & 91.95\,M & 218.62\,M \\
    Nodes/bldg.\ (med.)  & 30 & 54 & 253 & 31 & 21 & -- \\
    Thematic keys        & 19 & 49 & 22 & 5 & 20 & -- \\
    % OSM-enr.\ (\%)       & 85.3 & 83.1 & 91.5 & 98.8 & 58.9 & 74.8 \\
    OSM-enr.\ (\%)       & 84.9 & 81.1 & 91.2 & 98.7 & 57.8 & 74.1 \\
    Store (GiB)          & 14.8 & 0.3 & 30.2 & 32.4 & 99.2 & 176.9 \\
    Dump (GiB)           & 3.1 & 0.2 & 5.6 & 7.0 & 11.3 & 27.2 \\
    \bottomrule
  \end{tabular}
  
  \smallskip\raggedright\footnotesize
  \emph{Nodes/bldg.} is the median node count per building. \emph{Keys} counts distinct thematic property keys. \emph{OSM-enr.} is the share of buildings matched to at least one OSM feature. \emph{Store}/\emph{Dump} report Neo4j storage size and compressed release size. $^\dagger$ Reprojected at ingest from a non-metric CRS.
\end{table*}

% \begin{figure*}[h!]
%   \centering
%   \begin{subfigure}{0.325\textwidth}
%     \includegraphics[width=\linewidth]{figures/extent_hamburg.jpg}
%     \caption{Hamburg (zoomed-in), 1:84,460}
%   \end{subfigure}\hfill
%   \begin{subfigure}{0.325\textwidth}
%     \includegraphics[width=\linewidth]{figures/extent_tokyo.jpg}
%     \caption{Tokyo (23 wards), 1:46,760}
%   \end{subfigure}\hfill
%   \begin{subfigure}{0.325\textwidth}
%     \includegraphics[width=\linewidth]{figures/extent_nyc.jpg}
%     \caption{New York, 1:91,455}
%   \end{subfigure}\\[2pt]
%   \begin{subfigure}{0.325\textwidth}
%     % previous (city-core capture): \includegraphics{figures/extent_zurich_core.jpg}, Zurich, 1:\num{30120}
%     \includegraphics[width=\linewidth]{figures/extent_zurich_core.jpg} %extent_zurich
%     \caption{Zurich (zoomed-in), 1:120,480}
%   \end{subfigure}\hspace{0.0125\textwidth}
%   \begin{subfigure}{0.325\textwidth}
%     \includegraphics[width=\linewidth]{figures/extent_helsinki.jpg}
%     \caption{Helsinki (Kalasatama Digital Twins), 1:8734}
%   \end{subfigure}
%   \caption{
%     Spatial extent of the five \benchname{} cities. The more than tenfold variation in scale reflects source-specific release boundaries, from Helsinki's inner-city core to Zurich's canton-wide coverage. Basemap: \copyright~OpenStreetMap contributors.
%     % Spatial extent of the five \benchname{} cities: one ground footprint per building, from the source CityGML over the OpenStreetMap basemap. Map scales differ by over an order of magnitude (Helsinki's 1:\num{8734} core versus Zurich's 1:\num{120480} canton-wide extent), reflecting how differently the portals delimit releases: a federal state (Hamburg), a mesh-tiled metropolis (Tokyo), five boroughs (New York), a canton-wide tile scatter (Zurich), and an inner-city core (Helsinki). Basemap: \copyright~OpenStreetMap contributors.
%     }
%   \Description{Five map panels showing building-footprint extents over OpenStreetMap for Hamburg, Tokyo, New York, Zurich, and Helsinki. The maps use different scales, with Helsinki shown at the most zoomed-in scale and Zurich at the broadest extent.}
%   \label{fig:city_extents}
% \end{figure*}

% ===== BACKUP: superseded by the per-building panel of \input{city_graph_content_table.tex}; kept for reference, not compiled =====
\iffalse
\begin{table}
  \caption{Detail density of the authoritative CityGML layer, per fused city: graph nodes and on-disk store attributable to the CityGML-only ingestion stage, normalized per building. The order-of-magnitude spread quantifies the level-of-detail heterogeneity across open national 3D city models: it measures how much surface structure each source models per building, independent of city size. Helsinki includes its 31 bridge features; Tokyo's figures are restricted to the building module (\texttt{bldg} tiles), isolating per-building detail from the co-ingested B-core non-building modules (roads, bridges, city furniture, vegetation, water), which add a further \num{6.6}\,M nodes.}
  \label{tab:detail_density}
  \small
  \begin{tabular}{llrrr}
    \toprule
    City & LoD & Nodes & Nodes/bldg & KiB/bldg \\
    \midrule
    Hamburg & 2 & 15.2\,M & 39 & 34 \\
    Helsinki & 2 & 0.33\,M & 109 & 78 \\
    Zurich & 2.3 & 38.0\,M & 370 & 276 \\
    New York & 1--2 & 40.0\,M & 37 & 27 \\
    Tokyo & 1--2 & 57.9\,M & 29 & 43 \\
    \bottomrule
  \end{tabular}
\end{table}
\fi
% ===== end BACKUP =====

% ===== BACKUP: superseded by \input{city_storage_table.tex}; kept for reference, not compiled =====
\iffalse
\begin{table}
  \caption{On-disk storage per source-layer configuration (Neo4j store in GiB, write-ahead logs excluded; measured additively per ingestion stage). Layer deltas for the ML-predicted and reconstructed layers are estimated from record counts; all other figures are measured. Dashes: layer does not exist for that city (tier asymmetry). Tokyo's CityGML figure covers the full B-core module set (buildings plus roads, bridges, city furniture, vegetation, water), not buildings alone.}
  \label{tab:storage}
  \small
  \begin{tabular}{lrrrrrr}
    \toprule
    City & CityGML & +Roof & +LoD3 & +OSM & Full & Dump \\
    \midrule
    Hamburg & 12.71 & +0.00 & +0.00 & +2.05 & 14.77 & 3.07 \\
    Zurich & 27.06 & -- & -- & +3.15 & 30.21 & 5.57 \\
    Helsinki & 0.22 & -- & -- & +0.06 & 0.28 & 0.15 \\
    New York & 27.51 & -- & -- & +4.84 & 32.35 & 6.95 \\
    Tokyo & 91.9 & -- & -- & +7.3 & 99.2 & 11.3 \\
    \bottomrule
  \end{tabular}
\end{table}
\fi
% ===== end BACKUP =====

\subsection{Source Layers and the Provenance Spectrum}\label{subsec:provenance_spectrum}

\benchname{} is organized around an explicit \emph{provenance spectrum} (\Cref{fig:overview}): every fact belongs to one of four trust classes recoverable directly from the graph structure rather than external documentation. Three rules enforce this. First, a derived value never overwrites an authoritative one; it is stored under a source-specific property (\texttt{osm\_height} alongside the surveyed \texttt{measured\_height}), so both are comparable in one query. Second, each external object is a separate node that retains its source identity, metadata, and geometry, linked to the authoritative anchor by an explicit edge rather than merged into it. Third, every fusion edge stores its match confidence, so consumers make their own trust decisions rather than inheriting a fixed resolution.

% \begin{table}[h!]
%   \caption{The provenance spectrum in \benchname{}. Each trust class is structurally distinguishable in the graph.}
%   \label{tab:provenance_spectrum}
%   \begin{tabularx}{\linewidth}{lXX}
%     \toprule
%     Trust class & Example content & Encoding \\
%     \midrule
%     Authoritative & surveyed geometry, measured height, official function codes & canonical properties on feature nodes \\
%     Crowd-sourced & OSM tags, POIs, addresses, street network & \texttt{OsmFeature} nodes, \texttt{ENRICHED\_BY} edges, \texttt{osm\_*} properties \\
%     ML-predicted & roof materials (5 classes) from aerial imagery~\cite{lukas_roofmaterials} & prediction properties, partial coverage retained \\
%     Reconstructed & LoD3 facades (walls, roofs, windows, doors) from facade imagery & \texttt{Lod3Facade} nodes, \texttt{HAS\_LOD3\_FACADE} edges, \texttt{lod3\_*} properties \\
%     \bottomrule
%   \end{tabularx}
% \end{table}

\subsection{Authoritative Layer and Graph Construction}\label{subsec:construction}

The authoritative layer is derived from open-government CityGML~2.0 LoD2 datasets, using Hamburg's state mapping release~\cite{lgv_hhcitygml_2025} with ALKIS cadastral semantics. \mytoolname{} transforms CityGML into a compact Neo4j LPG~\cite{robinson_graphdb_2015}, where semantically meaningful elements become nodes, syntactic wrappers are merged into edges, and coordinates are preserved verbatim. Each top-level feature is additionally registered in an R-tree spatial index~\cite{neo4j_spatial_2025}, enabling semantic traversals and spatial predicates to be combined in a single query. 
%For Hamburg, the resulting graph comprises \num{388267} buildings, 15.2\,million nodes, and 21.3\,million edges. 
The transformation is idempotent (i.e., re-ingest yields an identical graph) and dataset-independent. Mapping details and losslessness evaluation are reported in the companion system paper~\cite{nguyen_pykci_2026}.
% The authoritative layer is built from open-government CityGML~2.0 LoD2 datasets, for Hamburg from the state mapping agency~\cite{lgv_hhcitygml_2025} with ALKIS cadastral semantics. \mytoolname{} maps CityGML to a compact LPG in Neo4j~\cite{robinson_graphdb_2015}: semantically meaningful elements become nodes centered on top-level features, serialization wrappers collapse into edges, coordinates are stored verbatim, and every top-level feature is registered by its 2D bounding box in an R-tree spatial index~\cite{neo4j_spatial_2025}, so that semantic traversal and spatial predicates (bounding box, intersection, proximity) compose in a single query. For Hamburg, the resulting authoritative layer contains \num{388267} buildings in a graph of about 15.2\,million nodes and 21.3\,million edges. The construction is idempotent and dataset-independent; we refer to the companion system paper for the mapping rules and the losslessness evaluation, which are not contributions of this paper.

% Closed vocabularies, such as ALKIS building-function and roof-type codes~\cite{adv_buildings_2025}, are represented both as scalar properties and as shared hub nodes linked to city features. Scalar values remain the authoritative representation, while vocabulary nodes expose categorical structure for graph queries. Each vocabulary node appears once per city graph and connects all features in the same category, enabling direct retrieval and aggregation through a single traversal. Predicted roof-material labels are currently stored only as source-prefixed scalar properties. 

CityGML datasets use different coordinate reference systems (CRS), including national metric grids (Hamburg, Zurich, and Helsinki), geographic coordinates (Tokyo), and US survey feet (New York). Rather than enforcing a common CRS, each city is maintained in a local metric CRS to preserve accurate distance and area measurements. 
% Non-metric sources are reprojected during ingestion, while original coordinates and CRS identifiers are retained as provenance. 
OSM data are reprojected into the corresponding city-specific CRS before fusion.

% \paragraph{Categorical vocabularies as hub nodes.} Closed vocabularies such as ALKIS building-function and roof-type codes are represented both as scalar properties and as shared vocabulary nodes linked directly to city features. Scalar values remain the authoritative source representation, while vocabulary nodes expose categorical structure to graph queries. Each vocabulary node exists only once in the city graph and connects all features that share the same category, enabling direct retrieval and aggregation through a single graph traversal. Predicted roof-material labels are currently stored only as source-prefixed scalar properties.
%
% \emph{Categorical vocabularies as hub nodes: a deliberate dual representation.} Closed categorical attributes, the ALKIS building-function and roof-form code lists, are represented twice, on purpose. The code itself is stored verbatim as a scalar property on each building (\texttt{function\_code}, \texttt{roof\_type\_code}); this is the authoritative, lossless round-trip representation. In addition, each code that actually occurs in a city materializes once as a \emph{vocabulary hub node} (\texttt{BuildingFunction}, \texttt{RoofType}) carrying the code and its human-readable name, with an explicit edge (\texttt{HAS\_FUNCTION}, \texttt{HAS\_ROOF\_TYPE}) from every building of that class. The hubs pay for themselves in the benchmark setting: the vocabulary becomes visible in the graph schema that a text-to-query model is prompted with (a property value never is), categorical questions compile to one-hop traversals whose per-class counts are node degrees, the hub is the natural home for class-level metadata, and a hub can be explicitly aligned with a neighboring taxonomy, the same mechanism by which the crowd-sourced and ML roof-material vocabularies of \Cref{subsec:conflict} can be related class-to-class rather than string-to-string. The costs are equally real: the class assignment is stored redundantly (scalar and edge must stay consistent), each categorical attribute adds one edge per building to the store, and a handful of high-degree shared nodes is exactly the structure that serializes concurrent writes; every edge creation locks its hub, and two parallel ingest workers merging the same not-yet-existing hub can even duplicate it, since the hub id carries no uniqueness constraint. \mytoolname{} therefore keeps the hubs out of the parallel ingestion hot path entirely: writers persist only the scalar codes, and the hubs with their edges are \emph{derived} from those scalars in a single-threaded pass after parsing, making the derivation race-free by construction and repeatable (the scalar remains the source of truth from which the hub layer can always be rebuilt). The predicted roof-material layer (\Cref{subsec:predicted_layers}) deliberately stays scalar-only with a source-prefixed key in this version; promoting it to a hub with per-edge confidence and model provenance is a natural schema extension, but one that changes graph content and is therefore deferred to a coordinated release rather than applied mid-version.

%\paragraph{City-local metric coordinate systems.} The source datasets use different coordinate reference systems (CRS), including national metric grids (Hamburg, Zurich, and Helsinki), geographic degrees (Tokyo), and US survey feet (New York). Rather than forcing all cities into a shared CRS, each city is maintained in a local metric CRS so that distances and areas remain accurate. Non-metric sources are reprojected during ingestion, while original coordinates and CRS identifiers are retained as provenance. OSM data are reprojected into the same city-specific CRS prior to fusion.
%
% \emph{Coordinate reference systems: one metric CRS per city, not a global one.} The cities come from different national mapping schemes and are published in different CRS and units (\Cref{tab:city_sources}): Hamburg, Zurich, and Helsinki in the metric national grids EPSG:25832, EPSG:2056, and EPSG:3879 respectively, whereas New York uses US~survey feet (EPSG:2263) and Tokyo geographic degrees (EPSG:6697). Crucially, there is \emph{no} single global CRS in which all five cities are simultaneously in true meters: a projected CRS is metric only near its own region, and a global pseudo-Mercator distorts distance and area by up to a factor of two at high latitude, so a shared CRS would corrupt exactly the distances and areas the benchmark measures. We therefore keep \emph{each city in its own local metric CRS}: the three already-metric cities are used as is, and the two non-metric sources are reprojected at ingest to a local metric CRS (Tokyo to EPSG:6677, New York to EPSG:32618), so every distance and area is expressed in true metres. This is sound because absolute positions are never compared across cities; every spatial query is posed within a single city. The OpenStreetMap layer is reprojected into each city's CRS during fusion, so the crowd-sourced and authoritative footprints co-register in the same metric system; footprint matching additionally uses an overlap ratio that is invariant to the CRS. Reprojection is applied only to the geometry and derived bounding boxes; the original coordinates and CRS are retained as provenance, and the already-metric Hamburg dataset is preserved verbatim for lossless round-trip.

\subsection{Identity Resolution and Integrity Check}\label{subsec:fidelity}

% Duplicated or placeholder identifiers occur in multiple input CityGML datasets: Zurich and Helsinki contain non-unique IDs, while Tokyo exhibits duplicate models at overlaps between ward-level CityGML tiles.
Several source datasets violate the assumption that \texttt{gml:id} values are unique, including 100 duplicated identifiers in Zurich and Helsinki and more than \num{975000} (33\,\% of input) ward-boundary identical buildings in Tokyo. 
We resolve these during ingestion without modifying the source files: features sharing an identifier are compared using content and geometry hashes, with identical features merged and distinct features assigned unique identifiers, preventing both erroneous merges and artificial duplication. We further perform pre-ingestion validation to detect identifier, geometry, coordinate, and CRS anomalies, followed by post-ingestion checks of graph integrity, including resolved object counts, identifier uniqueness, provenance completeness, and spatial-index coverage. 
% Applied to every city, these checks, together with \mytoolname{}'s lossless round-trip evaluation, ensure that all source content is either faithfully represented or explicitly accounted for.

% \emph{Identifier resolution without source modification.} Several source datasets violate the assumption that \texttt{gml:id} values are unique. Examples include placeholder or duplicated identifiers in Zurich and Helsinki dataset, as well as ward-boundary duplicates in Tokyo's model. We resolve these during ingestion without modifying the source files. Features sharing an identifier are compared using content and geometry hashes: identical features are merged, whereas distinct features are assigned unique identifiers. Exact geometry matching prevents both erroneous merges and artificial duplication.
% \emph{Identifier resolution without source modification.} Several source datasets violate the assumption that \texttt{gml:id} values are unique. Examples include placeholder and duplicated identifiers in Zurich's swiss\-BUILDINGS\-3D, \num{61} repeated building identifiers in Helsinki, and ward-boundary duplicates in Tokyo's PLATEAU release. During ingestion, \mytoolname{} resolves such collisions directly in the graphs without modifying the source files. Features sharing an identifier are compared using content and geometry hashes: identical features are merged, whereas distinct features are assigned synthesized internal identifiers while retaining the original \texttt{gml:id} as provenance. Exact geometry matching prevents both erroneous merges and artificial duplication, while provenance fields record every synthesized identifier. As a result, the graph preserves all source features, supports lossless round-trip export, and remains faithful to datasets with and without identifier collisions.
% \emph{We never edit the source; we resolve identifier collisions at ingest.} \mytoolname{} builds the graph by \texttt{MERGE}-ing each feature on its \texttt{gml:id}, which presumes that identifier is unique, an assumption several open datasets violate. Zurich's swissBUILDINGS3D stamps the placeholder \texttt{gml:id="ID\_"} on many buildings and even repeats real identifiers within a single tile; Helsinki repeats \num{61} building identifiers; and Tokyo's PLATEAU packages each ward separately, so a mesh tile on a ward boundary is shipped again in every adjacent ward. Merging on a raw duplicate id silently \emph{collapses} distinct features into one node, which is data loss, yet blindly keeping every occurrence would fabricate duplicates of a building that a dataset legitimately ships more than once. \mytoolname{} decides between the two by \emph{content}, without ever editing the id or the file. At ingest each top-level feature (building, building part, city object) is hashed over its full serialized subtree, geometry and attributes together. The first occurrence of a \texttt{gml:id} keeps it verbatim; a later occurrence with an \emph{identical} content hash is a true duplicate and is merged onto the same node; a later occurrence whose content differs is a candidate for a synthesized identifier \texttt{<id>\_\_<source>\_\_<n>}, unless a second, geometry-only signature says otherwise. That second tier exists because a building can be re-shipped with identical geometry but different ancillary attributes, as in Tokyo, where a boundary mesh tile recurs in adjacent ward packages and the copies carry ward-specific PLATEAU \texttt{uro:*} data-quality codes over coordinate-identical geometry; an exact content hash alone would over-split those into phantom buildings. A later occurrence whose full coordinate geometry is identical to an earlier one is therefore the same building and is merged; only differing (or absent) geometry splits. We compare geometry \emph{exactly} rather than by footprint overlap on purpose: two \emph{distinct} buildings that share an id can also share a footprint (Helsinki's duplicates already share bounding boxes), so a fuzzy overlap test would risk merging genuinely different buildings and silently violating the lossless round-trip, whereas coordinate-identical geometry is a safe same-building witness. Two provenance properties, set \emph{only} on a synthesized node, record the outcome (\texttt{source\_gml\_id} holds the original identifier and \texttt{id\_synthesized} flags the node), and on export the loader re-emits \texttt{source\_gml\_id}, so a duplicate-id source round-trips faithfully: two source features sharing an id re-export with that same id while remaining two distinct, independently queryable nodes. A dataset without collisions never dedups or splits and round-trips byte-for-byte.
%
% Structural-property collisions are handled similarly: generic attributes whose names conflict with reserved graph properties are stored under source-prefixed keys while preserving the original names for export.
% \emph{The same content-versus-identity discipline protects attribute keys.} A generic attribute may be named exactly like a structural property: Helsinki, for instance, stores each building's internal cadastral identifier as a \texttt{gen:stringAttribute} named ``ID''. Slugged naively this attribute would overwrite the feature's node identity. Any generic-attribute key that would collide with a reserved structural key is therefore stored under a source-prefixed \texttt{gen\_<key>}, with the original attribute name retained so export re-emits it unchanged; nested \texttt{gen:genericAttributeSet} wrappers, empty attribute values, and repeated attribute names are preserved by the same mechanism. This identity resolution recovers real buildings that a naive \texttt{MERGE} would lose (\num{61} in Helsinki and \num{45} in Zurich, each distinct feature now its own node) and deduplicates redundant copies that a naive keep-all would fabricate: of Tokyo's \num{975545} redundant ward-replicated building occurrences, \num{975077} are byte- or geometry-identical copies collapsed to one node each and the remaining 468 are distinct-content features sharing an id that are split apart, reducing \num{2980839} source elements to \num{2005762} distinct nodes. Faithfulness is verified rather than assumed: the authoritative-layer round-trip is proven lossless not only on synthetic per-module and per-LoD fixtures but on the full \num{2980}-building Helsinki source, whose duplicate identifiers, collision-inducing attribute names, nested attribute sets, multiple addresses per building, and dual LoD1/LoD2 solids exercise every one of these paths; \Cref{app:dupids} records every affected source building with its file, line range, CRS, bounding box, and footprint, so a data provider can locate and, if it chooses, correct the identifiers at the source.

% \emph{Integrity verification.} We employ a pre-ingestion validation to detect identifier, geometry, coordinate, and CRS anomalies in the source data. A post-ingestion check verifies graph integrity, including resolved object counts, identifier uniqueness, provenance completeness, and spatial-index coverage. Both checks are performed for every city. Together with \mytoolname{}'s lossless round-trip evaluation, they ensure that all source content is either faithfully represented or explicitly accounted for.
% \emph{Integrity verification.} We employ a pre-ingestion validation to detect identifier, geometry, coordinate, and CRS anomalies in the source data. A post-ingestion check verifies graph integrity, including resolved object counts, identifier uniqueness, provenance completeness, and spatial-index coverage. Both checks are performed for every city, and produce validation reports alongside the OSM ingestion census (\Cref{subsec:osm_fusion}). Together with the companion system paper's lossless round-trip evaluation, they ensure that all source content is either faithfully represented or explicitly accounted for.
% \emph{Integrity is gated before and after ingestion, not assumed.} Because the source is left unchanged, the graph's correctness must be checked rather than enforced by rewriting. A \emph{precheck} scans each source dataset before ingestion and reports the anomalies that would otherwise corrupt the graph or its indices: duplicate, placeholder, or missing \texttt{gml:id}; missing or degenerate geometry; non-finite or out-of-range coordinates; and a mixed or mis-declared coordinate reference system. A \emph{postcheck} then verifies the constructed graph: that every top-level feature carries a unique node id (no residual collapse), that the building count equals the resolved source count, that provenance (\texttt{source}) is present on every node, and that every spatially indexable feature is registered in the R-tree layer of its class, the conditions under which the schema's semantic traversal and spatial predicates are well defined. The precheck/postcheck pair ships with the release, is re-run per city, and auto-emits each city's identity-resolution summary alongside the OSM ingestion census (\Cref{subsec:osm_fusion}); together with the companion system paper's lossless round-trip of the authoritative layer, they close the ingestion path: what enters the graph is faithful to the source, and what the source contains is either represented or explicitly accounted for.

%\emph{Scope.} \benchname{} verifies faithful representation of source datasets but does not validate their geometric or semantic correctness. Such consistency checks remain the responsibility of data providers and specialized CityGML validation tools.
% \emph{Validating the internal consistency of the source is out of scope.} We check that the source is faithfully represented, not that it is internally consistent, and we draw that line deliberately: internal consistency is a property of the source data, is dataset- and jurisdiction-specific, and enforcing it would blur representing a source with re-authoring it. \benchname{} therefore does \emph{not} perform geometric-consistency checks (for example collinear or non-coplanar surface vertices, self-intersecting or non-watertight solids), semantic-consistency checks (whether coded building functions, usages, or generic attributes conform to a controlled specification such as the AdV product standard for German 3D building models~\cite{adv_buildings_2025}), or mixed geometric--semantic checks (for example a \texttt{measuredHeight} attribute against the geometry's own $z_{\max}-z_{\min}$ span). Such conformance validation belongs to the data providers and to dedicated CityGML checking tools; the traceability record of \Cref{app:dupids} is intended precisely to feed the identifier findings back into that provider-side process.

\subsection{Crowd-Sourced Layer: OSM Fusion}\label{subsec:osm_fusion}

OSM provides an independent description of the city. Fusion focuses on footprint correspondence, where one building may map to multiple polygons in either source (1:1, 1:n, n:1, n:m). Candidate pairs $A, B$ come from R-tree intersections and are retained when  $r=\operatorname{area}(A\cap B)/\min(\operatorname{area}(A),\operatorname{area}(B))\ge\tau$.
%(with \(\tau=0.3\)); 
The overlap coefficient \(r\) serves as the acceptance criterion (\(\tau = 0.3\)). Each accepted edge stores additionally the Jaccard index \(\operatorname{area}(A\cap B)/\operatorname{area}(A\cup B)\). 
%, the intersection area, the correspondence case, and a primary-match flag identifying the dominant OSM contributor by intersection area. 
Connected components of the bipartite overlap graph define correspondence classes, and component-level rules propagate source-prefixed OSM attributes to matched buildings (\Cref{app:enriched_osm}, \Cref{fig:osm_enrich_example}). Geometry and points of interest remain on OSM nodes.
% The overlap coefficient $r$ is the acceptance criterion (\(\tau = 0.3\)) and is stored on the resulting edge together with the Jaccard index \(\operatorname{area}(A\cap B)/\operatorname{area}(A\cup B)\), the intersection area, the correspondence case, and a primary-match flag that marks each building's dominant OSM contributor (by intersection area). Connected components of the bipartite overlap graph classify correspondences, and component-level rules copy source-prefixed OSM attributes to matched buildings (\Cref{app:enriched_osm}, \Cref{fig:osm_enrich_example}). Geometry and points of interest remain on OSM nodes.
% The overlap coefficient $r$ is computed for each candidate pair (\(\tau = 0.3\), selected by threshold sweep). Connected components of the bipartite overlap graph classify correspondences, and component-level rules copy source-prefixed OSM attributes to matched buildings (\Cref{app:enriched_osm}, \Cref{fig:osm_enrich_example}). Geometry and points of interest remain on OSM nodes. 
\Cref{app:enriched_osm} explains how the value of $\tau$ is selected. 
% OSM contributes a second, independently produced description of the same city. Fusion is dominated by footprint correspondence: the same building may be one OSM polygon and several cadastral ones, or vice versa (1:1, 1:n, n:1, n:m). We resolve all four with one mechanism. Candidate pairs are generated by R-tree intersection; a pair is kept if the overlap ratio $r = \operatorname{area}(A \cap B) / \min(\operatorname{area}(A), \operatorname{area}(B)) \geq \tau$ (default $\tau = 0.3$, from a threshold sweep), with the Jaccard index on every surviving edge. Connected components of the bipartite overlap graph classify the case, and component-level rules decide which OSM values are copied (source-prefixed) onto which buildings (\Cref{fig:osm_enrich_example}); geometry and points of interest stay on the OSM node and are reached by traversal.
%
Unlike traditional conflation pipelines, we retain the full weighted overlap graph as \texttt{ENRICHED\_BY} edges rather than enforcing a final matching, which is an NP-hard problem~\cite{naumann_nnmatching_2024}. 
%In Hamburg, fusion links \num{354162} OSM buildings to \num{388267} CityGML buildings: \num{331318} (85.3\,\%) receive at least one correspondence, while 94.4\,\% of OSM buildings find a counterpart, and \num{22219} points of interest connect via \texttt{HAS\_POI}. The remaining \num{831432} unmatched OSM features are preserved as spatially indexed standalone nodes. \Cref{subsec:conflict} analyzes the resulting correspondence structure and attribute redundancy.
% Unlike one-shot conflation pipelines, we do not commit to a hard partition, which would require an NP-complete optimization~\cite{naumann_nnmatching_2024}; the full weighted overlap graph is retained as \texttt{ENRICHED\_BY} edges carrying overlap ratio, Jaccard, intersection area, case label, and a primary-match flag. On Hamburg the fusion confronts \num{354162} OSM buildings (\num{1188139} OSM features) with \num{388267} cadastral buildings: \num{331318} (85.3\,\%) receive at least one edge, 94.4\,\% of OSM buildings find a counterpart, \num{22219} points of interest attach by \texttt{HAS\_POI} to their containing building, and \num{831432} OSM features with no cadastral counterpart in the buildings-only LoD2 source (streets, furniture, land use, vegetation, water) are kept as spatially indexed standalone nodes rather than discarded. \Cref{subsec:conflict} analyses the correspondence structure and attribute-level redundancy this creates.
All OSM features are additionally registered in a second, dedicated R-tree layer, so crowd-sourced and authoritative geometry remain independently and jointly queryable.

% Candidate footprints are scoped by the same R-tree index in which every feature is registered at ingest, so the fusion runs whole-city rather than tile-by-tile, which makes per-city re-runs and threshold sweeps practical. All OSM features are additionally registered in a second, dedicated R-tree layer, so crowd-sourced and authoritative geometry remain independently and jointly queryable.

%\emph{Ingestion completeness verification.} An ingestion census re-derives the expected graph from the source extract and validates five properties: feature existence, tag losslessness, geometry fidelity, placement (either anchored to a CityGML feature or flagged unmatched), and spatial indexing with provenance. All fused instances pass with zero violations, from Helsinki's \num{55930} admitted features to Tokyo's \num{3830408} (\Cref{tab:osm_census}). The census also makes admission boundaries explicit: points and polygons are ingested without thematic filtering, lines only for transport and water networks, while remaining features are non-network linear objects or relation views over already-ingested members.
%\emph{Ingestion completeness is gated, not assumed.} An ingestion census verifies the crowd-sourced layer with the rigor the companion system paper applies to the authoritative round-trip. It re-derives the expected graph content from the source extract through the identical code path and checks the stored graph at five levels: feature existence (each admitted feature present exactly once under its stable OSM id), tag losslessness, geometry fidelity (coordinates verbatim in the city's metric CRS, with exact bounding boxes), placement (every feature either anchored to a CityGML feature or flagged unmatched with a machine-readable reason), and spatial indexing with full provenance. Every fused instance passes with zero violations on all five checks, from Helsinki's \num{55930} admitted features to Tokyo's \num{3830408} (\Cref{tab:osm_census}). The census also makes the admission boundary explicit, accounting for every feature the source exposes in each city's bounding box: points and polygons enter without a thematic filter, lines only for the transport and water network, and the remainder is non-network linear features and relation views over already-ingested members.

\subsection{ML-Predicted and Reconstructed Layer}\label{subsec:predicted_layers}

The ML layer adds roof-material classes (concrete, metal, glass, roof tiles, tar paper) predicted from aerial orthophotos~\cite{lukas_roofmaterials}. 
% Rather than retaining only the majority label as in prior work~\cite{kanna2026semantic}, \benchname{} preserves each building's full material distribution, storing every predicted class with its pixel coverage. 
\benchname{} preserves each building's full material distribution, storing every predicted class with its pixel coverage. 
%This exposes both the dominant material and minority classes (\num{570} of \num{194799} predicted Hamburg buildings contain multiple materials). 
Fusion is a direct ID join, with all values stored under source-prefixed keys. Predictions cover \num{194799} of \num{388267} buildings (50.2\,\%), limited by orthophoto availability. 
% This information allows coverage-aware queries (\Cref{subsec:query_tasks}) and imputation task (\Cref{subsec:rl_tasks}).
% The ML-predicted layer attaches roof-material classes (concrete, metal, glass, roof tiles, tar paper) predicted from aerial orthophotos~\cite{lukas_roofmaterials}. The classifier is per-pixel, so a building's native output is a \emph{distribution} over classes. Where the source semantic-mapping work~\cite{kanna2026semantic} keeps only the majority label, \benchname{} retains the full per-building composition: each predicted material with its pixel-coverage fraction, which doubles as a per-material confidence and orders the classes, so the dominant class stays directly available while minority materials remain queryable (\num{570} of \num{194799} predicted Hamburg buildings carry more than one). Fusion is by \texttt{gml:id} with no spatial matching, and every value lives under a source-prefixed key, so the layer is additive. Predictions exist for \num{194799} of \num{388267} buildings (50.2\,\%), bounded by orthophoto availability; we keep this partial coverage deliberately, since the covered and uncovered subsets are cleanly separable and drive the coverage-aware query category (\Cref{subsec:query_tasks}) and the imputation task (\Cref{subsec:rl_tasks}).

The reconstructed layer adds LoD3 building models derived from facade imagery for 17 buildings in Hamburg's HafenCity district. The models contribute 416 windows and 83 doors as explicit \texttt{Opening} nodes, represented by 499 interior rings in wall geometries. %The source CityGML was not schema-valid, so a reproducible, non-destructive transformation reclassified levels of detail, canonicalized element order, and recomputed envelopes, reducing 17 validation errors to zero while preserving all coordinates and the original files. 
Corrected facades are linked to their LoD2 buildings via \texttt{HAS\_LOD3\_FACADE} edges as separate nodes in a dedicated spatial layer, never replacing the measured geometry. 
% (\Cref{fig:lod3_facades})
% Owing to their limited number, these models support showcase queries only and are excluded from statistical analyses.
% The reconstructed layer integrates LoD3 building models derived from facade imagery \verify{method and citation TBD} for 17 buildings in Hamburg's HafenCity district, whose walls carry 416 windows and 83 doors as explicit \texttt{Opening} nodes, with 499 interior rings cutting those openings as holes in the wall geometry. As delivered, the source CityGML was not schema-valid (building solid emitted after its boundary surfaces, LoD3 detail tagged as LoD2, an envelope not enclosing the geometry); a non-destructive, reproducible transform reclassifies the level of detail from opening and hole content, canonicalizes element order against the CityGML schema, and recomputes the envelope, turning 17 validation errors into none with every coordinate preserved and the original file retained. The corrected facades attach to their LoD2 anchors by explicit \texttt{HAS\_LOD3\_FACADE} edges as separate nodes in a dedicated spatial layer, never overwriting the measured geometry (\Cref{fig:lod3_facades}). Given their small count they power a showcase question category only and are excluded from statistical claims.


\subsection{Cross-Source Agreement and Conflict}\label{subsec:conflict}

\benchname{}'s datasets are highly diverse 
%despite sharing a common CityGML and OSM schema 
(see~\Cref{app:dataset_stats},~\Cref{fig:city_stats}). Fusing these urban data sources creates both coverage and redundancy. We analyze geometric correspondences and dual-sourced attributes in Hamburg, Helsinki, Zurich, New York, and Tokyo. 
% All statistics are computed directly from the fused graph via Cypher queries released with the benchmark.
% Fusing independent descriptions of the same city does not only fill gaps: wherever two sources describe the same object, the graph keeps both opinions side by side. We quantify, for Hamburg in detail and with Helsinki, Zurich, and New York as contrasting points, how the sources partition the built fabric into building objects (geometric level) and where an authoritative or ML value and a crowd-sourced value describe the same property (attribute level). All statistics are computed on the fused graph with Cypher over the \texttt{ENRICHED\_BY} edges and source-prefixed property pairs, and the queries ship with the release, so every number below is reproducible.

% \begingroup
% \definecolor{cham}{HTML}{2A78D6}\definecolor{chel}{HTML}{008300}%
% \definecolor{czur}{HTML}{E87BA4}\definecolor{cnyc}{HTML}{EDA100}\definecolor{ctok}{HTML}{1BAF7A}%
% \colorlet{cgray}{black!7}%
% \tikzset{cbar/.style={rounded corners=0.5pt},
%   rl/.style={left,font=\normalsize,inner sep=2.5pt},
%   vl/.style={right,font=\small,text=black!55,inner sep=2.5pt}}%
% \begin{figure*}[ht!]
%   \centering
%   % Each of the four quadrant charts is forced to an identical bounding box
%   % (-2.1..6.5 wide) inside a 0.5\textwidth box, so all four have the SAME width,
%   % exactly half of the full-width composition chart below.
%   % ---- row 1: detail density | schema richness ----
%   \begin{minipage}[t]{0.5\textwidth}\centering
%     {\normalsize\bfseries Detail density: nodes per building}\par\vspace{2pt}
%     \begin{tikzpicture}\def\MX{369.8}\def\AW{5.0}\def\bh{0.26}
%       \useasboundingbox(-2.1,0.15)rectangle(6.5,-3.35);
%       \foreach \lbl/\md/\mn/\c [count=\i] in {Hamburg/30/41/cham,Helsinki/54/107.4/chel,Zurich/253/369.8/czur,New York/31/36.9/cnyc,Tokyo/21/28.9/ctok}{
%         \pgfmathsetmacro\yb{-\i*0.62}\pgfmathsetmacro\wm{\md/\MX*\AW}\pgfmathsetmacro\wa{\mn/\MX*\AW}
%         \fill[cgray,cbar](0,\yb+0.32)rectangle(\AW,\yb+0.32+\bh);\fill[\c,cbar](0,\yb+0.32)rectangle(\wm,\yb+0.32+\bh);\node[vl]at(\wm,\yb+0.32+\bh/2){\md};
%         \fill[cgray,cbar](0,\yb)rectangle(\AW,\yb+\bh);\fill[\c!42!white,cbar](0,\yb)rectangle(\wa,\yb+\bh);\node[vl]at(\wa,\yb+\bh/2){\mn};
%         \node[rl]at(0,\yb+0.29){\lbl};}
%     \end{tikzpicture}\par\vspace{-2pt}{\small\color{black!55}Darker = median \qquad lighter = mean}
%   \end{minipage}%
%   \begin{minipage}[t]{0.5\textwidth}\centering
%     {\normalsize\bfseries Storage by source layer}\par\vspace{2pt}
%     \begin{tikzpicture}\def\MX{99.18}\def\AW{5.0}\def\bh{0.38}
%       \useasboundingbox(-2.1,0.15)rectangle(6.5,-3.35);
%       \foreach \lbl/\cg/\tot/\tt/\c [count=\i] in {Hamburg/12.71/14.77/15/cham,Helsinki/0.22/0.28/0.28/chel,Zurich/27.06/30.21/30/czur,New York/27.51/32.35/32/cnyc,Tokyo/91.92/99.18/99/ctok}{
%         \pgfmathsetmacro\yy{-\i*0.62}\pgfmathsetmacro\wcg{\cg/\MX*\AW}\pgfmathsetmacro\wt{\tot/\MX*\AW}
%         \fill[cgray,cbar](0,\yy)rectangle(\AW,\yy+\bh);\fill[\c,cbar](0,\yy)rectangle(\wcg,\yy+\bh);\fill[\c!38!white,cbar](\wcg,\yy)rectangle(\wt,\yy+\bh);
%         \node[rl]at(0,\yy+\bh/2){\lbl};\node[vl]at(\wt,\yy+\bh/2){\tt\,GiB};}
%     \end{tikzpicture}\par\vspace{-2pt}{\small\color{black!55}Darker = CityGML \qquad lighter = $+$ OSM}
%   \end{minipage}
%   \par\vspace{4pt}
%   % ---- row 2: storage | enrichment ----
%   \begin{minipage}[t]{0.5\textwidth}\centering
%     {\normalsize\bfseries Source schema richness: thematic keys}\par\vspace{2pt}
%     \begin{tikzpicture}\def\MX{49}\def\AW{5.0}\def\bh{0.38}
%       \useasboundingbox(-2.1,0.15)rectangle(6.5,-3.35);
%       \foreach \lbl/\v/\c [count=\i] in {Hamburg/19/cham,Helsinki/49/chel,Zurich/22/czur,New York/5/cnyc,Tokyo/20/ctok}{
%         \pgfmathsetmacro\yy{-\i*0.62}\pgfmathsetmacro\bw{\v/\MX*\AW}
%         \fill[cgray,cbar](0,\yy)rectangle(\AW,\yy+\bh);\fill[\c,cbar](0,\yy)rectangle(\bw,\yy+\bh);
%         \node[rl]at(0,\yy+\bh/2){\lbl};\node[vl]at(\bw,\yy+\bh/2){\v};}
%     \end{tikzpicture}%\par\vspace{0pt}{\small\strut}
%   \end{minipage}%
%   \begin{minipage}[t]{0.5\textwidth}\centering
%     {\normalsize\bfseries OSM enrichment rate}\par\vspace{2pt}
%     \begin{tikzpicture}\def\MX{100}\def\AW{5.0}\def\bh{0.38}
%       \useasboundingbox(-2.1,0.15)rectangle(6.5,-3.35);
%       \draw[dashed,black!30](\AW,0.02)--(\AW,-3.15);\node[font=\small,text=black!45,below]at(\AW,-3.15){100\%};
%       \foreach \lbl/\v/\c [count=\i] in {Hamburg/85.3/cham,Helsinki/83.1/chel,Zurich/91.5/czur,New York/98.8/cnyc,Tokyo/58.9/ctok}{
%         \pgfmathsetmacro\yy{-\i*0.62}\pgfmathsetmacro\bw{\v/\MX*\AW}
%         \fill[cgray,cbar](0,\yy)rectangle(\AW,\yy+\bh);\fill[\c,cbar](0,\yy)rectangle(\bw,\yy+\bh);
%         \node[rl]at(0,\yy+\bh/2){\lbl};\node[vl]at(\bw,\yy+\bh/2){\v\%};}
%     \end{tikzpicture}%\par\vspace{1pt}{\small\strut}
%   \end{minipage}
%   % \par\vspace{6pt}
%   % ---- full-width: node composition (monochrome per city); width = 2x a quadrant ----
%   % {\normalsize\bfseries Node composition: every graph is geometry-dominated}\par\vspace{2pt}
%   % \begin{tikzpicture}\def\AW{15.0}\def\bh{0.40}
%   %   \useasboundingbox(-2.1,0.78)rectangle(15.1,-3.32);
%   %   \foreach \nm/\mix/\lx in {Geometry/100/0,Semantic/64/2.9,{Crowd-sourced OSM}/42/5.9,{Index / containers}/22/10.6}{
%   %     \fill[cham!\mix!white,cbar](\lx,0.30)rectangle(\lx+0.32,0.62);\node[right,font=\small,text=black!55]at(\lx+0.38,0.46){\nm};}
%   %   \foreach \lbl/\sem/\geo/\osm/\idx/\gt/\c [count=\i] in {Hamburg/2.2/90.8/6.8/0.1/91/cham,Helsinki/0.7/85.6/13.5/0.2/86/chel,Zurich/0.3/95.2/4.5/0.1/95/czur,New York/2.4/91.9/5.6/0.1/92/cnyc,Tokyo/3.6/91.1/5.1/0.1/91/ctok}{
%   %     \pgfmathsetmacro\yy{-\i*0.60}\pgfmathsetmacro\tot{\sem+\geo+\osm+\idx}
%   %     \pgfmathsetmacro\wg{\geo/\tot*\AW}\pgfmathsetmacro\ws{\sem/\tot*\AW}\pgfmathsetmacro\wo{\osm/\tot*\AW}
%   %     \pgfmathsetmacro\xb{\wg+\ws}\pgfmathsetmacro\xc{\wg+\ws+\wo}
%   %     \fill[\c,cbar](0,\yy)rectangle(\wg,\yy+\bh);\fill[\c!64!white,cbar](\wg,\yy)rectangle(\xb,\yy+\bh);\fill[\c!42!white,cbar](\xb,\yy)rectangle(\xc,\yy+\bh);\fill[\c!22!white,cbar](\xc,\yy)rectangle(\AW,\yy+\bh);
%   %     \node[left,font=\normalsize,inner sep=3pt]at(0,\yy+\bh/2){\lbl};\node[font=\small,text=white]at(\wg/2,\yy+\bh/2){\gt\%};}
%   % \end{tikzpicture}
%   \caption{
%     Comparison of the five \benchname{} graphs. Zurich is the densest, Tokyo the largest, Helsinki the most attribute-rich, and New York the most OSM-enriched. Geometry accounts for 85--95\,\% of graph content in every city. Full statistics in \Cref{app:dataset_stats}.
%     % Per-city comparison of the five \benchname{} graphs (fixed color per city). \emph{Detail density}: median (solid) and mean (faded) nodes per building; level of detail, not city size, drives graph scale (a Zurich building costs $\sim$12$\times$ a Tokyo one). \emph{Schema richness}: distinct source thematic keys (Helsinki~49 vs New York~5), the multilingual axis. \emph{Storage}: the authoritative CityGML layer (darker) dominates, OSM fusion (lighter) adds a modest slice. \emph{OSM enrichment rate}: \num{58.9}--\num{98.8}\,\%. \emph{Node composition}: every graph is 85--95\,\% geometry. Full per-city figures in \Cref{app:dataset_stats}.
%     }
%   \Description{Five side-by-side bar-chart panels comparing Hamburg, Helsinki, Zurich, New York, and Tokyo across graph characteristics: detail density (median and mean nodes per building), schema richness (distinct thematic keys), storage split between CityGML and OSM, OSM enrichment rate percentages, and node composition dominated by geometry.}
%   \label{fig:city_stats}
% \end{figure*}
% \endgroup

% \begingroup
% \definecolor{cham}{HTML}{2A78D6}\definecolor{chel}{HTML}{008300}%
% \definecolor{czur}{HTML}{E87BA4}\definecolor{cnyc}{HTML}{EDA100}\definecolor{ctok}{HTML}{1BAF7A}%
% \colorlet{cgray}{black!7}%
% \tikzset{cbar/.style={rounded corners=0.5pt},
%   rl/.style={left,font=\normalsize,inner sep=2.5pt},
%   vl/.style={right,font=\small,text=black!55,inner sep=2.5pt}}%
% \begin{figure*}[ht!]
%   \centering
%   % Each of the four quadrant charts is forced to an identical bounding box
%   % (-2.1..6.5 wide) inside a 0.5\textwidth box, so all four have the SAME width,
%   % exactly half of the full-width composition chart below.
%   % ---- row 1: detail density | schema richness ----
%   \begin{minipage}[t]{0.25\textwidth}\centering
%     {\footnotesize\bfseries Nodes / building}\par\vspace{1pt}\resizebox{\linewidth}{!}{%
%     \begin{tikzpicture}\def\MX{369.8}\def\AW{5.0}\def\bh{0.26}
%       \useasboundingbox(-2.1,0.15)rectangle(6.5,-3.35);
%       \foreach \lbl/\md/\mn/\c [count=\i] in {Hamburg/30/41/cham,Helsinki/54/107.4/chel,Zurich/253/369.8/czur,New York/31/36.9/cnyc,Tokyo/21/28.9/ctok}{
%         \pgfmathsetmacro\yb{-\i*0.62}\pgfmathsetmacro\wm{\md/\MX*\AW}\pgfmathsetmacro\wa{\mn/\MX*\AW}
%         \fill[cgray,cbar](0,\yb+0.32)rectangle(\AW,\yb+0.32+\bh);\fill[\c,cbar](0,\yb+0.32)rectangle(\wm,\yb+0.32+\bh);\node[vl]at(\wm,\yb+0.32+\bh/2){\md};
%         \fill[cgray,cbar](0,\yb)rectangle(\AW,\yb+\bh);\fill[\c!42!white,cbar](0,\yb)rectangle(\wa,\yb+\bh);\node[vl]at(\wa,\yb+\bh/2){\mn};
%         \node[rl]at(0,\yb+0.29){\lbl};}
%     \end{tikzpicture}}
%   \end{minipage}\hfill
%   \begin{minipage}[t]{0.25\textwidth}\centering
%     {\footnotesize\bfseries Storage (GiB)}\par\vspace{1pt}\resizebox{\linewidth}{!}{%
%     \begin{tikzpicture}\def\MX{99.18}\def\AW{5.0}\def\bh{0.38}
%       \useasboundingbox(-2.1,0.15)rectangle(6.5,-3.35);
%       \foreach \lbl/\cg/\tot/\tt/\c [count=\i] in {Hamburg/12.71/14.77/15/cham,Helsinki/0.22/0.28/0.28/chel,Zurich/27.06/30.21/30/czur,New York/27.51/32.35/32/cnyc,Tokyo/91.92/99.18/99/ctok}{
%         \pgfmathsetmacro\yy{-\i*0.62}\pgfmathsetmacro\wcg{\cg/\MX*\AW}\pgfmathsetmacro\wt{\tot/\MX*\AW}
%         \fill[cgray,cbar](0,\yy)rectangle(\AW,\yy+\bh);\fill[\c,cbar](0,\yy)rectangle(\wcg,\yy+\bh);\fill[\c!38!white,cbar](\wcg,\yy)rectangle(\wt,\yy+\bh);
%         \node[rl]at(0,\yy+\bh/2){\lbl};\node[vl]at(\wt,\yy+\bh/2){\tt\,GiB};}
%     \end{tikzpicture}}
%   \end{minipage}\hfill
%   \begin{minipage}[t]{0.25\textwidth}\centering
%     {\footnotesize\bfseries Thematic keys}\par\vspace{1pt}\resizebox{\linewidth}{!}{%
%     \begin{tikzpicture}\def\MX{49}\def\AW{5.0}\def\bh{0.38}
%       \useasboundingbox(-2.1,0.15)rectangle(6.5,-3.35);
%       \foreach \lbl/\v/\c [count=\i] in {Hamburg/19/cham,Helsinki/49/chel,Zurich/22/czur,New York/5/cnyc,Tokyo/20/ctok}{
%         \pgfmathsetmacro\yy{-\i*0.62}\pgfmathsetmacro\bw{\v/\MX*\AW}
%         \fill[cgray,cbar](0,\yy)rectangle(\AW,\yy+\bh);\fill[\c,cbar](0,\yy)rectangle(\bw,\yy+\bh);
%         \node[rl]at(0,\yy+\bh/2){\lbl};\node[vl]at(\bw,\yy+\bh/2){\v};}
%     \end{tikzpicture}}
%   \end{minipage}\hfill
%   \begin{minipage}[t]{0.25\textwidth}\centering
%     {\footnotesize\bfseries OSM enrichment}\par\vspace{1pt}\resizebox{\linewidth}{!}{%
%     \begin{tikzpicture}\def\MX{100}\def\AW{5.0}\def\bh{0.38}
%       \useasboundingbox(-2.1,0.15)rectangle(6.5,-3.35);
%       \draw[dashed,black!30](\AW,0.02)--(\AW,-3.15);\node[font=\small,text=black!45,below]at(\AW,-3.15){100\%};
%       \foreach \lbl/\v/\c [count=\i] in {Hamburg/85.3/cham,Helsinki/83.1/chel,Zurich/91.5/czur,New York/98.8/cnyc,Tokyo/58.9/ctok}{
%         \pgfmathsetmacro\yy{-\i*0.62}\pgfmathsetmacro\bw{\v/\MX*\AW}
%         \fill[cgray,cbar](0,\yy)rectangle(\AW,\yy+\bh);\fill[\c,cbar](0,\yy)rectangle(\bw,\yy+\bh);
%         \node[rl]at(0,\yy+\bh/2){\lbl};\node[vl]at(\bw,\yy+\bh/2){\v\%};}
%     \end{tikzpicture}}
%   \end{minipage}
%   % \par\vspace{6pt}
%   % ---- full-width: node composition (monochrome per city); width = 2x a quadrant ----
%   % {\normalsize\bfseries Node composition: every graph is geometry-dominated}\par\vspace{2pt}
%   % \begin{tikzpicture}\def\AW{15.0}\def\bh{0.40}
%   %   \useasboundingbox(-2.1,0.78)rectangle(15.1,-3.32);
%   %   \foreach \nm/\mix/\lx in {Geometry/100/0,Semantic/64/2.9,{Crowd-sourced OSM}/42/5.9,{Index / containers}/22/10.6}{
%   %     \fill[cham!\mix!white,cbar](\lx,0.30)rectangle(\lx+0.32,0.62);\node[right,font=\small,text=black!55]at(\lx+0.38,0.46){\nm};}
%   %   \foreach \lbl/\sem/\geo/\osm/\idx/\gt/\c [count=\i] in {Hamburg/2.2/90.8/6.8/0.1/91/cham,Helsinki/0.7/85.6/13.5/0.2/86/chel,Zurich/0.3/95.2/4.5/0.1/95/czur,New York/2.4/91.9/5.6/0.1/92/cnyc,Tokyo/3.6/91.1/5.1/0.1/91/ctok}{
%   %     \pgfmathsetmacro\yy{-\i*0.60}\pgfmathsetmacro\tot{\sem+\geo+\osm+\idx}
%   %     \pgfmathsetmacro\wg{\geo/\tot*\AW}\pgfmathsetmacro\ws{\sem/\tot*\AW}\pgfmathsetmacro\wo{\osm/\tot*\AW}
%   %     \pgfmathsetmacro\xb{\wg+\ws}\pgfmathsetmacro\xc{\wg+\ws+\wo}
%   %     \fill[\c,cbar](0,\yy)rectangle(\wg,\yy+\bh);\fill[\c!64!white,cbar](\wg,\yy)rectangle(\xb,\yy+\bh);\fill[\c!42!white,cbar](\xb,\yy)rectangle(\xc,\yy+\bh);\fill[\c!22!white,cbar](\xc,\yy)rectangle(\AW,\yy+\bh);
%   %     \node[left,font=\normalsize,inner sep=3pt]at(0,\yy+\bh/2){\lbl};\node[font=\small,text=white]at(\wg/2,\yy+\bh/2){\gt\%};}
%   % \end{tikzpicture}
%   \caption{
%     Comparison of the five \benchname{} graphs. Zurich is the densest, Tokyo the largest, Helsinki the most attribute-rich, and New York the most OSM-enriched. Geometry accounts for 85--95\,\% of graph content in every city. Full statistics in \Cref{app:dataset_stats}.
%     % Comparison of the five \benchname{} graphs. Zurich is the densest, Tokyo the largest, Helsinki the most attribute-rich, and New York the most OSM-enriched. Geometry accounts for 85--95\,\% of graph content in every city. In the first two panels, darker bars denote the median and the CityGML layer and lighter bars the mean and the added OSM layer. Full statistics in \Cref{app:dataset_stats}.
%     % Per-city comparison of the five \benchname{} graphs (fixed color per city). \emph{Detail density}: median (solid) and mean (faded) nodes per building; level of detail, not city size, drives graph scale (a Zurich building costs $\sim$12$\times$ a Tokyo one). \emph{Schema richness}: distinct source thematic keys (Helsinki~49 vs New York~5), the multilingual axis. \emph{Storage}: the authoritative CityGML layer (darker) dominates, OSM fusion (lighter) adds a modest slice. \emph{OSM enrichment rate}: \num{58.9}--\num{98.8}\,\%. \emph{Node composition}: every graph is 85--95\,\% geometry. Full per-city figures in \Cref{app:dataset_stats}.
%     }
%   \Description{Five side-by-side bar-chart panels comparing Hamburg, Helsinki, Zurich, New York, and Tokyo across graph characteristics: detail density (median and mean nodes per building), schema richness (distinct thematic keys), storage split between CityGML and OSM, OSM enrichment rate percentages, and node composition dominated by geometry.}
%   \label{fig:city_stats}
% \end{figure*}
% \endgroup

\paragraph{Geometric disagreement.} Most matched buildings align cleanly: 87.2\,\% of Hamburg's \num{293895} correspondences are 1:1 matches (\Cref{tab:footprint_cases}). 
%The remaining 12.8\,\% involve \num{73615} cadastral buildings and reflect systematic differences in how sources partition buildings. 
In 1:n components, one OSM building covers 2.5 CityGML buildings on average; in n:1 components, one CityGML building corresponds to 3.0 OSM buildings. The \num{2363} n:m components motivate retaining the full overlap graph rather than forcing a single matching. Fragmentation also varies by city, from 0.4\,\% in New York to 
%9.5\,\% in Helsinki.
9.8\,\% in Tokyo.
% \emph{Geometric and topological disagreement.} The two sources agree on a building's delineation far more often than not, but the disagreement is substantial in absolute terms (\Cref{tab:footprint_cases}): 87.2\,\% of Hamburg's \num{293895} correspondence components are clean 1:1 matches, yet the remaining 12.8\,\% involve \num{73615} cadastral buildings, more than a fifth of all matched ones. The direction is interpretable: in 1:n components one OSM polygon spans on average 2.5 cadastral buildings (a mapper digitized a perimeter block or terraced row the cadastre splits into parcels), while in n:1 components one cadastral building is covered by 3.0 OSM polygons on average (wings, annexes, and \texttt{building:part} detail the cadastre merges). The \num{2363} n:m components resist any clean assignment, which is exactly why the graph keeps the full weighted overlap subgraph rather than forcing a partition. The case mix is itself a per-city property spanning an order of magnitude, from New York's 0.4\,\% fragmented (its crowd mapping tracks the cadastre almost one-to-one) through Zurich's 3.3\,\% and Hamburg's 6.3\,\% to Helsinki's 9.5\,\%, where Finnish mapping splits more buildings into tagged parts; the same mechanism absorbs all four.

\begin{table}[h!]
  \caption{Footprint correspondence between CityGML and OSM (share of
  components per city). Full counts in Table~\ref{tab:footprint_cases}.}
  \label{tab:footprint_cases_lite}
  \begin{tabular}{lrrrrr}
    \toprule
    City & 1:1 & 1:$n$ & $n$:1 & $n$:$m$ & Fragmented \\
    \midrule
    Hamburg  & 87.2 & 6.5 & 5.5 & 0.8 & 6.3 \\
    Helsinki & 84.2 & 6.3 & 6.4 & 3.1 & 9.5 \\
    Zurich   & 90.8 & 5.8 & 2.4 & 0.9 & 3.3 \\
    New York & 99.3 & 0.3 & 0.3 & 0.1 & 0.4 \\
    Tokyo    & 85.9 & 4.3 & 2.3 & 7.5 & 9.8 \\
    \bottomrule
  \end{tabular}
  
  \smallskip\raggedright\footnotesize
  Fragmented comprises the \(n{:}1\) and \(n{:}m\) cases, that is, CityGML buildings split across multiple OSM building footprints. 
  % The case mix varies by more than an order of magnitude across cities, motivating retention of the full weighted overlap graph instead of a single matching.
\end{table}

\begin{table}[h!]
  \caption{OSM ingestion census.}
  \label{tab:osm_census}
  %\small
  \begin{tabular}{lrrrrrr}
    \toprule
    City & Ingested & Coverage & Anchored & Standalone \\
    \midrule
    Hamburg  & \num{1188139} & 95.6\,\% & \num{356707} & \num{831432} \\
    Helsinki & \num{55930}   & 89.7\,\% & \num{4281}   & \num{51649} \\
    Zurich   & \num{1860401} & 95.8\,\% & \num{102794} & \num{1757607} \\
    New York & \num{2555468} & 97.6\,\% & \num{1121901} & \num{1433567} \\
    Tokyo    & \num{3830408} & 97.9\,\% & \num{1235353} & \num{2595055} \\
    \bottomrule
  \end{tabular}

  \smallskip\raggedright\footnotesize
  % All ingested OSM features are \emph{verified}, meaning they pass all five completeness gates (stored exactly once under a stable OSM ID; every tag reconstructed losslessly; geometry stored verbatim in the city's metric CRS with exact bounding boxes; placed as either an anchored feature or a reasoned standalone node; registered in the R-tree with full provenance and dataset-membership edge). 
  %\emph{In bbox} counts every source feature GDAL exposes inside the city bounding box (dataset extent plus fusion buffer) across all five OSM layers. 
  \emph{Ingested} is how many OSM features enter the graph, not including non-network lines like barriers, building-annotation open ways, and man-made linear features.
  %Not ingested, by design: non-network lines (Hamburg \num{44341}: \num{16949} barriers, \num{13454} building-annotation open ways, \num{4151} man-made linear features, the rest power/aeroway/administrative; Helsinki \num{5781}: \num{2999} barriers, \num{1353} man-made, \num{555} building-annotation, the rest natural/administrative; Zurich \num{68518}: \num{39430} barriers, \num{4994} building-annotation, \num{3870} man-made, the rest power/cable/golf/natural/administrative; New York \num{45680}: \num{1814} power, the rest barrier/man-made/coastline/building-annotation service ways) and the two relation layers (Hamburg \num{10038}, Helsinki \num{616}, Zurich \num{12107}, New York \num{15894}), which are views over member geometry already ingested through the base layers. 
  \emph{Coverage}: ingested features as a share of all source OSM features in the city's bounding box. 
  % The remainder is non-network lines and relation layers excluded by design. 
  \emph{Anchored} vs.\ \emph{Standalone} split the ingested features by whether they attach to a CityGML feature or are kept as net-new nodes.
\end{table}

\paragraph{Attribute agreement and conflict.} Height and storey counts largely agree when they are present (\Cref{tab:dual_attrs}): storeys match exactly for 87.2\,\% of \num{144364} buildings (98.6\,\% within $\pm1$), and the median height difference over \num{3404} buildings is 0.27\,m. Disagreements remain informative, including 219 buildings differing by more than 5\,m. Cross-city behavior differs substantially: Zurich shows much lower height agreement (median $|\Delta h|=2.3$\,m), likely reflecting differing measurement conventions. 
% rather than survey error.
% \emph{Attribute redundancy and conflict.} Three building properties are dual-sourced at scale in Hamburg (\Cref{tab:dual_attrs}). For storeys and height the crowd-sourced value largely confirms the authoritative one: storey counts agree exactly on 87.2\,\% of \num{144364} buildings (98.6\,\% within $\pm$1), and the median absolute height difference over \num{3404} buildings is 0.27\,m. The disagreeing tail is the interesting part: 219 buildings differ in height by more than 5\,m, each a genuine, queryable open question (outdated survey, wrong tag, or a building changed between acquisitions) rather than a data-cleaning casualty. The profile differs sharply by city. Helsinki carries a surveyed height on nearly every building but no roof material, so its dual coverage is height (\num{76} buildings, median $|\Delta h|$ 0.57\,m) and storeys (\num{48}) only. Zurich narrows to height alone (\num{770} buildings), and its agreement is far looser (median $|\Delta h|$ 2.3\,m, 22.6\,\% differing by more than 5\,m), less survey error than definition: the Swiss height and the OSM \texttt{height} tag need not share a roof datum. That the same comparison is tight in Hamburg and loose in Zurich is precisely the discrepancy the dataset exposes rather than silently reconciles.

Roof material provides the richest comparison.
%because the second source is an ML prediction. 
Among \num{3641} buildings with both ML and OSM labels, agreement reaches 75.2\,\%.
%after mapping OSM tags to the five-class taxonomy. 
Agreement is highest for roof tiles (92.6\,\%) and lower for flat-roof materials.
% , where some disagreements reflect different semantic perspectives rather than errors. 
OSM also contributes \num{5059} roof-material values without predictions and 134 values outside the prediction taxonomy.
% Roof material is the richest case, because here the second opinion is an ML \emph{prediction} rather than a survey. For \num{3641} buildings both an ML class and a matched OSM \texttt{roof:material} tag exist. After folding the free-text OSM vocabulary onto the five prediction classes, the sources agree on 75.2\,\% of mappable pairs, with a strongly class-dependent structure: agreement is high for the visually distinctive pitched-roof class (92.6\,\% of model-predicted roof tiles are tagged tiles in OSM) and lower for flat-roof classes, where the model's \emph{concrete} often meets OSM's \emph{tar paper} or \emph{gravel}. This is not pure error: a flat concrete roof with a bituminous covering is correctly described by both terms, the model reporting the structure from the orthophoto and the mapper the covering from the street, so part of the disagreement is a semantic layering difference the graph keeps inspectable. OSM also extends the layer: \num{5059} buildings carry an OSM roof material where no prediction exists, and 134 carry values outside the five-class taxonomy (slate, thatch, green roofs) that the model cannot express.
%
%As a result, dual-sourced subsets provide natural evaluation data for imputation and provenance-aware queries (\Cref{subsec:query_tasks}), while coverage itself varies by city, from Helsinki's heights and storeys to Zurich's height-only overlap and New York's purely geometric redundancy.
%
% Two consequences follow. First, redundancy is a feature: the dual-sourced subsets are freely obtained evaluation sets (the agreement rates above are what an imputation model must beat), and the cross-source and provenance question categories (\Cref{subsec:query_tasks}) instantiate directly over them. Second, dual coverage is a per-city variable, not a constant of the construction, depending on local tagging culture and which layers exist, from Helsinki's height and storeys through Zurich's height only to New York's none at all (no surveyed attribute and no ML layer, so its redundancy is purely geometric and locational). The release reports the per-city equivalents of \Cref{tab:footprint_cases,tab:dual_attrs}.

% =====================================================================
% Flagship HafenCity teaser variants (rendering of the same scene).
% These are provided as four SEPARATE figures so any of them can be
% commented out independently.
% =====================================================================
% \begin{figure*}[t]
%   \centering
%   \includegraphics[width=\textwidth]{figures/Teaser_AuthentiCity_opaque.png}
%   \caption{Flagship HafenCity (Versmannstra\ss e) instance of \benchname{}, rendered as OGC~3D~Tiles over an OpenStreetMap basemap with fully opaque buildings colored by their \emph{machine-learned} roof material (blue metal, green glass, graphite tar paper, sand concrete, red tiles) and the \emph{reconstructed} LoD3 building in gold, with the spatially aligned knowledge graph (nodes and edges) overlaid.}
%   \label{fig:teaser_opaque_kg}
% \end{figure*}
%
% \begin{figure*}[t]
%   \centering
%   \includegraphics[width=\textwidth]{figures/Teaser_AuthentiCity_opaque_noKG.png}
%   \caption{The same opaque multi-source 3D city model as \Cref{fig:teaser_opaque_kg} \emph{without} the knowledge-graph overlay, showing the machine-learned roof-material coloring on its own (blue metal, green glass, graphite tar paper, sand concrete, red tiles; reconstructed LoD3 building in gold).}
%   \label{fig:teaser_opaque_nokg}
% \end{figure*}
%
% \begin{figure*}[t]
%   \centering
%   \includegraphics[width=\textwidth]{figures/Teaser_AuthentiCity.png}
%   \caption{Translucent variant of the flagship scene: the same buildings rendered semi-transparent so the aligned knowledge graph reads \emph{through} them, over the OpenStreetMap basemap. Buildings are colored by machine-learned roof material (muted palette) with the reconstructed LoD3 building in gold.}
%   \label{fig:teaser_soft_kg}
% \end{figure*}
%
% \begin{figure*}[t]
%   \centering
%   \includegraphics[width=\textwidth]{figures/Teaser_AuthentiCity_closeup.png}
%   \caption{Close-up on the \emph{photogrammetrically reconstructed} LoD3 building (gold) and its dense per-opening subgraph, with the surrounding HafenCity streets legible; translucent buildings expose the aligned knowledge graph.}
%   \label{fig:teaser_closeup}
% \end{figure*}

% \Cref{fig:city_stats} gives a comparison of all five city graphs after enrichment of OSM data, predicted roof materials and reconstructed LoD3 facade geometries.

\subsection{Multilingual Urban Knowledge Graphs}\label{subsec:multilingual}

Helsinki, Hamburg, and especially Tokyo illustrate why natural-language interfaces matter for urban KGs. Tokyo's PLATEAU data uses Japanese attribute names and values, for example \jp{地区計画} (district plan) and \jp{市谷柳町地区}, which \mytoolname{} preserves verbatim as Unicode property keys. Traditionally, querying such data required familiarity with both the local schema and natural language. With LLM-grounded querying, users can ask questions in their own language, 
%while the model maps them to the underlying Japanese attributes, 
as shown in this work, making Tokyo's 100\,GiB graph accessible to both local residents and international analysts.

% \subsection{Two-Tier Composition and Statistics}\label{subsec:tiers}

% \benchname{} is released in two tiers. \textbf{Tier~1 (comparable core)} applies the identical CityGML-plus-OSM construction to every city, so core questions instantiate per city and generalization becomes measurable; \textbf{Tier~2 (deep fusion)} is the Hamburg instance carrying all four source layers. Asymmetric data availability is thereby a documented property of the release rather than a hidden artifact.

% \emph{Why these five cities.} Each contributes a distinct axis of real-world heterogeneity, and each yields an enriched \mytoolname{} graph interesting in its own right. Hamburg supplies the ALKIS cadastre of an entire German city-state and is the only instance carrying all four provenance layers, making it the deep-fusion flagship. Zurich contributes the highest geometric fidelity, swisstopo's swissBUILDINGS3D~3.0 at LoD2.3. Helsinki is a compact, fully inspectable CC~BY core of \num{2980} buildings that doubles as a reproducibility reference. Tokyo is the opposite extreme: \num{2980839} buildings across the 23 special wards, a Japanese attribute schema, and a multi-module PLATEAU core rather than buildings alone. New York adds a fifth national scheme over five boroughs, ingested from a US-survey-feet source. Together they span five national CityGML schemes, five metric coordinate reference systems, three continents, and both Latin and Japanese scripts, so \benchname{} measures cross-schema and cross-language generalization rather than single-city tuning.

% \emph{Why CityGML and OSM as the base pairing.} CityGML is the ISO/OGC standard for semantic 3D city models~\cite{groeger_citygml2_2012}, and national CityGML products are already the authoritative distillation of the underlying government registers: German LoD2 takes footprints and function codes from the ALKIS cadastre and adds laser-scanning heights~\cite{adv_buildings_2025}, and PLATEAU attaches use, structure, and construction-year attributes from the statutory urban-planning survey~\cite{plateau_3dcitymodel_2021}. Ingesting the cadastre separately would largely duplicate this layer, so we leave a parcel layer to future work rather than break the Tier~1 contract. CityJSON~\addedcite{ledoux_cityjson_2019} is a lossless alternative encoding, so CityJSON-native products such as the Dutch 3D~BAG~\addedcite{peters_3dbag_2022} remain reachable by conversion, and global footprint corpora~\addedcite{microsoft_footprints_2022,sirko_openbuildings_2021,overture_maps_2024} offer coverage but neither roof geometry nor thematic semantics with fine provenance. OSM is the complement at the opposite end of the trust spectrum: the largest crowd-sourced corpus, globally uniform, rich in the use-level detail cadastral products omit. The pairing is the one configuration both reproducible in every city and genuinely multi-source.

% National mapping schemes make the cities non-identical along three axes (\Cref{tab:city_sources}): \emph{different code lists} (ALKIS building-function and roof-type codes have no exact PLATEAU or US counterpart); varying \emph{levels of detail}; and, most visibly, \emph{attribute names and values written in the local language}. In the Tokyo PLATEAU data the generic-attribute names and many values are Japanese, for instance \jp{地区計画} (``district plan'') with a value such as \jp{市谷柳町地区} (Ichigaya-Yanagich\=o); \mytoolname\ preserves them verbatim as Unicode property keys, so they stay queryable in their original script. For Tokyo we take the PLATEAU 2025 (v5) release over the 23 special wards and a multi-module core (buildings, roads, bridges, city furniture, vegetation, water), excluding the area-tessellating and ADE extension modules; even so it is the largest instance, with \num{2005762} building nodes and \num{659006} non-building city objects in one R-tree layer. Cross-city questions are defined over a \emph{common attribute subset} present in every city (geometry, height, a coarse function class, the OSM layer), with city-specific attributes marked as such, so the ``comparable core'' is an explicit contract. Together the five cities contribute \num{4558196} building features, from \num{2980} in Helsinki to \num{2980839} across Tokyo's wards (\Cref{fig:city_extents}).

% This heterogeneity is where a natural-language interface earns its place: rather than learning each city's code list and attribute language, a user queries in their own language and an LLM maps the question onto the local schema, so a German- or English-speaking analyst and a Japanese resident can both query the Japanese-annotated Tokyo graph. Multilingual, code-list-agnostic querying is thus an evaluation dimension \benchname{} uniquely enables among 3D-city benchmarks, and a direct argument for an LLM-grounded query interface rather than fixed templates.

% \emph{Storage cost per source layer.} Assembled layer by layer, the graph's storage is attributable per layer via the \texttt{source} property on every node (\Cref{tab:city_storage}). For Hamburg the authoritative CityGML layer dominates at 12.71\,GiB; OSM fusion adds 2.05\,GiB (2.26\,M nodes, 5.05\,M edges, all tags, geometry, and its own R-tree layer), while the ML and LoD3 layers are essentially free ($<$\,0.01\,GiB each), quantifying how cheaply a derived opinion layer rides on an existing anchor graph. The fully fused instance compresses to a 3.07\,GiB dump.

% \emph{Detail density: geometric detail, not building count, drives graph size.} Normalizing per building makes the level-of-detail heterogeneity of the national sources measurable (\Cref{tab:city_graph_content}): a median building costs 30 nodes in Hamburg, 31 in New York, and 21 in Tokyo, but 54 in Helsinki and 253 in Zurich (LoD2.3), an order-of-magnitude spread that reflects how much roof and facade structure each agency models rather than city size. This is why the \num{102673}-building Zurich instance (30.2\,GiB) outweighs the \num{388267}-building Hamburg one, and it is benchmark-relevant: 3D-structure questions traverse proportionally more geometry per answer in a high-detail city (\Cref{fig:city_stats}).

% \subsection{Release, Documentation, and Licensing}\label{subsec:release}

% The release comprises (1) the graph per city as a Neo4j dump plus a backend-neutral export (node and edge tables with a documented schema) so that consumers are not bound to one database; the dumps are compact (the fully fused Hamburg instance, 17.4\,M nodes with all four source layers and its spatial indexes, compresses to 3.07\,GiB, \Cref{tab:city_storage}), (2) loaders and task splits, (3) the question suite with gold queries and materialized gold answers, (4) the evaluation harness, including the verification gates that every instance must pass before release (the OSM ingestion completeness census of \Cref{subsec:osm_fusion} and the cross-source statistics queries of \Cref{subsec:conflict}), and (5) a datasheet following~\addedcite{gebru_datasheets_2021} documenting collection, preprocessing, recommended uses, and known biases. The source CityGML data for Hamburg is published under the dl-de/by-2-0 open data license by the state agency~\cite{lgv_hhcitygml_2025,dlde_by20_2026}; OSM data is under ODbL; our annotations, question suite, and documentation are released under CC~BY~4.0. The artifact is deposited on Zenodo with a DOI \verify{[DOI to be minted at deposit time]}.

\section{Benchmark A: Natural-Language-to-Query}\label{subsec:query_tasks}

To show that the released graphs can be queried in natural language, and that doing so demands reasoning beyond what conventional text-to-SQL and text-to-Cypher benchmarks test, we define a natural-language-to-query benchmark. Given a natural-language question and the graph schema, a model must either produce an executable Cypher query whose result matches the gold answer or explicitly declare the question infeasible. We evaluate on LPGs by design. While a relational CityGML store can support aggregate, filter, and spatial queries, the benchmark's cross-source agreement, provenance- and confidence-filtered retrieval, and coverage-aware aggregation rely on LPG-native fusion structures absent from relational schemas.
%such as \texttt{ENRICHED\_BY} edges, source-prefixed \texttt{osm\_}/\texttt{predicted\_} properties, and per-fact source and coverage metadata that are absent from the relational schema.
% Given a natural-language question and the graph schema as context, a model must either produce an executable Cypher query whose result matches the gold answer or explicitly declare the question infeasible. We evaluate against the LPG by design: a relational CityGML store expresses the aggregate, filter, and spatial categories, but not this benchmark's contribution, cross-source agreement, provenance- and confidence-filtered retrieval, and coverage-aware aggregation, which read fusion structures (the \texttt{ENRICHED\_BY} edges, the source-prefixed \texttt{osm\_*}/\texttt{predicted\_*} scalars, and per-fact source and coverage) that a relational schema has no column for.

\subsection{Categories}

We design \num{1394} query questions spanning 84 templates and nine categories (\Cref{tab:categories}; full inventory in \Cref{app:questions}), five of which are our contributions. \emph{Spatial} questions, the largest category, cover window, proximity, nearest-neighbor, geometric multi-hop, 3D-structure, and density queries. \emph{Cross-source} questions compare authoritative and crowd-sourced values and ask what the crowd-sourced layer reports where authoritative data is absent, a common case in New York, where \num{97.5}\,\% of buildings have an OSM height but none an authoritative one. \emph{Provenance-filtered} questions constrain answers by source or confidence, including traps (``using only authoritative data\ldots'') where reading a crowd-sourced value yields a plausible but incorrect answer. \emph{Coverage-aware} questions require recognizing that aggregates over partially covered attributes are meaningful only with respect to the covered subset; the flagship trap asks for a city-wide roof-material count, where the gold answer couples count and coverage. \emph{Infeasible} questions have no valid answer and test whether a model declines rather than fabricates a query, each paired with a guard query proving the information absent. 

Two design choices increase difficulty. First, as feasibility depends on city-specific coverage, the same template can be coverage-aware in one city and infeasible in another (e.g., roof material is a coverage question in Hamburg but infeasible elsewhere). Second, ten templates are additionally posed in German and Japanese, including three Tokyo templates whose gold queries filter on Japanese property keys. A thematic subset targets city-specific schema content, including Helsinki's floor area, volume, and construction year; Zurich's data-vintage fields; and Tokyo's Japanese-keyed attributes 
%, identified through an audit of each city's property inventory 
(\Cref{app:thematic}).

% Questions fall into categories (\Cref{tab:categories}; full inventory in \Cref{app:questions}), of which five are the distinctive contribution. \emph{Spatial} questions, the largest category, span window, proximity, and nearest-neighbor queries, geometric multi-hop questions that compose a spatial predicate with semantic traversal (the tallest building within 200\,m of a caf\'e), 3D-structure questions over boundary surfaces, and density. \emph{Cross-source} questions compare authoritative against crowd-sourced values in both directions and ask what the crowd-sourced layer reports where the authoritative source is silent, which matters concretely on New York, where \num{97.5}\,\% of buildings carry an OSM height but none an authoritative one. \emph{Provenance-filtered} questions constrain by source or match confidence, including traps (``using only authoritative data\ldots'') where reading the crowd-sourced property yields a plausibly close but wrong answer. \emph{Coverage-aware} questions require recognizing that an aggregate over a partially covered attribute is only meaningful relative to the covered subset; the flagship trap asks a city-wide count over the roof-material attribute, and the gold answer couples count with coverage, so a bare count scores wrong. \emph{Infeasible} questions have no correct answer and score whether the model says so rather than fabricating a query, each with a guard query proving the referenced information absent. Two devices multiply what the categories measure: feasibility is a property of the data, so the same question flips category across cities (roof material is a coverage question on Hamburg, infeasible elsewhere), and a multilingual axis poses ten templates additionally in German and Japanese, with three Tokyo templates whose gold query filters on Japanese property keys. A further subset targets thematic content native to one city's schema (Helsinki's floor area, volume, and construction year; Zurich's data-vintage fields; Japanese-keyed Tokyo attributes), added after auditing the per-city property inventory (\Cref{app:thematic}).

\subsection{Construction, Verification, and Metrics}

The benchmark is generated from 84 schema-grounded templates, each paired with a hand-authored gold Cypher query whose slot values are read directly from the released graph. Feasible templates declare the graph tokens they require and instantiate only where those tokens exist; missing attributes are exercised through the infeasible subfamily. Every question passes a five-stage verification gate: gold queries execute, are deterministic, and are non-empty unless permitted; infeasible questions include a guard query; and each rebuild re-materializes all gold answers. This extends the methodology of~\cite{chauhan_mindthequery_2025}. Additionally, \num{282} questions (20.2\,\%) were manually reviewed, revealing a proximity-predicate defect that escaped automated checks and was fixed before release.

The benchmark is split by template into a public development set and a held-out test set with unpublished gold queries. The test set contains unseen templates and holds out 370 of \num{1394} questions (26.5\,\%). We report execution accuracy (result-set equivalence under canonicalization), decomposed into executed-and-correct (EX), executed-but-wrong, errored, and over-refusal. This distinction separates Cypher-generation failures from semantic failures. We also report infeasibility precision, recall, and F1, 
%together with query latency, 
following prior NL-to-Cypher benchmarks~\cite{feng_cypherbench_2025, li_bird_2023}.

% Questions are generated from 84 schema-grounded templates, each with a hand-authored gold Cypher query whose every slot value is read from the released graph, so no question references absent data and the suite regenerates against new versions. Each feasible template declares the graph tokens its gold query reads and instantiates on a city only when the graph provides them; the same missing attribute is otherwise exercised by the absent-layer infeasible subfamily. Every question passes a five-stage gate (the gold query executes, is deterministic, and is non-empty unless permitted; infeasible questions carry a guard query; each rebuild re-materializes all gold answers), extending the discipline of~\addedcite{chauhan_mindthequery_2025}; \num{282} questions (20.2\,\%) were additionally human-reviewed, which caught a proximity-predicate defect the executable gates had passed (five radius templates mis-scaled against the metric CRS), fixed before release. The suite is split by template into a public development set and a held-out test set whose gold queries are unpublished, so the test set holds unseen templates; it holds out 370 of \num{1394} questions (26.5\,\%). We report execution accuracy (result-set equivalence under canonicalization), decomposed on the feasible questions into executed-and-correct (EX), executed-but-wrong, errored, and over-refusal, since these separate a model that cannot emit valid Cypher from one that emits valid but semantically wrong Cypher, plus infeasibility precision/recall/F1 and query latency, keeping the metric set compatible with~\addedcite{feng_cypherbench_2025, li_bird_2023}.

\subsection{Setup and Results}\label{subsec:experiments_setup}

We evaluate a local open-weight model (qwen2.5-coder:7b, Ollama \texttt{Q4\_K\_M}) and a commercial frontier model (Claude Sonnet~5, run as a closed-book Claude Code subagent). Both models receive byte-identical, schema-only context with no database or tool access. Outputs are cached and scored offline against the gold queries (examples in~\Cref{app:questions_examples}).
%(full protocol in \Cref{app:eval_protocol}). 
Even the frontier model leaves substantial headroom. Claude achieves 54--69\,\% EX across the five cities (\Cref{tab:results_query}). Nearly all remaining cases are \emph{executed-but-wrong} rather than errored, indicating semantic rather than syntax failures. qwen reaches only 6--19\,\% EX and is dominated by \emph{errored} queries (40--54\,\%), largely Cypher syntax errors in spatial tasks where it hallucinates PostGIS \texttt{ST\_*} functions. This decomposition reveals qualitatively different failure modes that aggregate EX obscures.
\begin{table*}[ht!]
  \caption{
  Task family A results on all five city dev and held-out test splits (\%). The four feasible outcomes sum to 100: EX (correct), Ex-wr (executed but incorrect), Err (execution failure), and O-ref (incorrect refusal). Inf-F1 measures infeasibility detection. 
  %Full results are provided in \Cref{app:full_results_a}.
  % Task family A results on all five Tier-1 city dev and held-out test splits (percentages). The four \emph{feasible} columns are mutually exclusive and sum to 100: EX (executed and result-correct), Ex-wr (executed but wrong), Err (query did not execute), O-ref (wrongly declined a feasible question). Inf-F1 is infeasibility-detection F1 (over-refusals are the false positives). qwen2.5-coder never refuses, so its Inf-F1 is undefined (n/a) and its recall is~0; its Err mass is overwhelmingly Cypher syntax errors concentrated in the spatial category (hallucinated PostGIS \texttt{ST\_*} functions). NYC's and Tokyo's test-split Claude Err are spatial query timeouts on the largest graphs (\Cref{subsec:experiments_setup}). Full per-category and per-language results are in \Cref{app:full_results}.
  }
  \label{tab:results_query}
  % \footnotesize
  \setlength{\tabcolsep}{4pt}
  \begin{tabular}{ll ccccc c ccccc}
    \toprule
    & & \multicolumn{5}{c}{Development set} & & \multicolumn{5}{c}{Held-out test set} \\
    \cmidrule(lr){3-7}\cmidrule(lr){9-13}
    Model & City & EX\,$\uparrow$ & Ex-wr & Err & O-ref\,$\downarrow$ & Inf-F1 & & EX\,$\uparrow$ & Ex-wr & Err & O-ref\,$\downarrow$ & Inf-F1 \\
    \midrule
    Claude   & Hamburg  & 60.7 & 32.6 &  0.0 & 6.7 & 72.0 & & 58.6 & 41.4 &  0.0 & 0.0 & 100.0 \\
    Sonnet 5 & Helsinki & 56.1 & 39.0 &  0.0 & 4.9 & 78.3 & & 42.9 & 57.1 &  0.0 & 0.0 &  66.7 \\
             & Zurich   & 61.5 & 38.5 &  0.0 & 0.0 & 94.7 & & 60.0 & 40.0 &  0.0 & 0.0 & 100.0 \\
             & New York & 69.2 & 28.2 &  0.0 & 2.6 & 90.0 & & 50.0 & 35.7 & 14.3 & 0.0 & 100.0 \\
             & Tokyo    & 53.7 & 37.3 &  7.5 & 1.5 & 90.0 & & 41.7 & 37.5 & 20.8 & 0.0 & 100.0 \\
    \addlinespace
    qwen2.5- & Hamburg  & 19.1 & 36.0 & 44.9 & 0.0 & n/a & & 20.7 & 48.3 & 31.0 & 0.0 & n/a \\
    coder:7b & Helsinki & 14.6 & 45.1 & 40.2 & 0.0 & n/a & & 10.7 & 53.6 & 35.7 & 0.0 & n/a \\
             & Zurich   &  9.2 & 41.5 & 49.2 & 0.0 & n/a & & 15.0 & 40.0 & 45.0 & 0.0 & n/a \\
             & New York & 10.3 & 35.9 & 53.8 & 0.0 & n/a & & 14.3 & 28.6 & 57.1 & 0.0 & n/a \\
             & Tokyo    &  6.0 & 43.3 & 50.7 & 0.0 & n/a & & 12.5 & 33.3 & 54.2 & 0.0 & n/a \\
    \bottomrule
  \end{tabular}

  \smallskip\raggedright\footnotesize
  Inf-F1 is the F1 score for infeasibility detection, where over-refusals count as false positives. qwen2.5-coder never refuses, so Inf-F1 is undefined (n/a) and recall is 0. For qwen2.5-coder, most Err cases are Cypher syntax errors in spatial queries caused by hallucinated PostGIS \texttt{ST\_*} functions. For Claude, most Err cases on the NYC and Tokyo test splits are spatial-query timeouts on the largest graphs (\Cref{subsec:experiments_setup}). Baselines are evaluated on a stratified subsample of the released suite; per-split counts are in the repository. 
\end{table*}
On infeasibility, qwen never refuses (recall~0), generating a query for every unanswerable question, whereas Claude declines appropriately (Inf-F1 72--95 on dev). 
% Errors cluster in the novel provenance, coverage, and infeasibility categories rather than standard aggregation or filtering. 
% Held-out test results track development performance on smaller graphs. 
Declines on New York and Tokyo stem primarily from spatial-query timeouts on the largest graphs (1--2M buildings) rather than semantic drift.
% (\Cref{app:full_results_a}).

% We evaluate a local open-weight model (qwen2.5-coder:7B, Ollama \texttt{Q4\_K\_M}) and a commercial frontier model (Claude Sonnet~5), with pinned versions, cached outputs, and a fixed per-city schema prompt. Even the frontier model leaves substantial headroom: Claude reaches 54--69\,\% EX across the five cities (\Cref{tab:results_query}), its non-EX mass almost entirely \emph{executed-but-wrong} rather than errored, so it emits valid Cypher and fails semantically. qwen reaches only 6--19\,\% EX, dominated by \emph{errored} queries (40--54\,\%), overwhelmingly Cypher syntax errors in the spatial category where it hallucinates PostGIS \texttt{ST\_*} functions; the decomposition is the point, as aggregate EX alone hides that these are two different failures, one semantic and one an inability to emit valid Cypher at all. On infeasibility qwen never refuses (recall~0), fabricating a query for every unanswerable question, exactly the failure the metric exposes, whereas Claude refuses appropriately (Inf-F1 72--95 on dev). Failures concentrate in the novel provenance-, coverage-, and infeasibility categories rather than in aggregation or filtering. Held-out test tracks dev on the smaller graphs; the New York and Tokyo drops are spatial query timeouts on the largest graphs (1--2\,M buildings) rather than semantic drift (\Cref{app:full_results}).

% %%%%%%%%%%%%%%%%%%%%%%%%%%LUKAS STUFF%%%%%%%%%%%%%%%%%%%%%%%%%%%%%%%
% %%%%%%%%%%%%%%%%%%%%%%%%%%%%%%%%%%%%%%%%%%%%%%%%%%%%%%%%%%%%%%%%%%%%
% %%%%%%%%%%%%%%%%%%%%%%%%%%%%%%%%%%%%%%%%%%%%%%%%%%%%%%%%%%%%%%%%%%%%

% \section{Benchmark B: Representation Learning}\label{subsec:rl_tasks}

% To highlight that the graph supports applications beyond text-to-query tasks and is applicable to a variety of downstream analytical tasks, including energy efficiency, environmental sustainability, and urban planning~\citep{LIU2022102936}, we provide an embedding-based evaluation following~\citet{dwivedi_benchmarkgnn_2023} for heterogeneous urban knowledge graphs. The evaluation comprises three tasks, each conducted under fixed, reproducible splits spanning random, spatial-block, and cross-city protocols. The first task is \emph{attribute imputation}, which predicts the roof-material class for buildings that lack this attribute. Roof-material information is available for only approximately \num{50}\% of buildings and is provided by an ML model trained on crowd-sourced OSM labels~\citep{lukas_roofmaterials}. The second task is \emph{node classification}, predicting the ALKIS building-function class and in addition, the roof-type class for held-out buildings using administrative ground-truth labels. By design, the two roof-information targets span the trust spectrum, with authoritative roof-type labels at one end and ML-predicted, OSM-derived roof-material labels at the other. Consequently, only the latter requires a circularity check (Section~\ref{sec:repL:results}). The third task is \emph{link prediction}, predicting held-out \texttt{ENRICHED\_BY} edges and thereby casting footprint correspondence as a learning task~\cite{naumann_nnmatching_2025}. Each task is evaluated under both a \emph{provenance-agnostic} protocol, in which source distinctions are hidden as in existing urban knowledge graphs, and a \emph{provenance-weighted} (aware) protocol, in which trust classes and fusion-edge confidence scores are exposed, for example by weighting message passing using the Jaccard index. The key result is the $\Delta$ between these protocols, indicating whether provenance information improves performance or is ignored by standard embedding methods. To the best of our knowledge, this is the first such measurement on a real city-scale multi-source graph. We position provenance-aware representation learning as an open problem for which our \benchname{} provides both the benchmark substrate and an initial empirical data point (\ref{tab:results_rl}), rather than a complete solution, with baseline results released.

% % \begin{table}[t]
% %   \caption{Task family B (representation learning), evaluated under a provenance-agnostic protocol (source distinctions hidden) and a provenance-weighted protocol (trust classes and fusion-edge confidences exposed); the $\Delta$ measures whether provenance information helps. Metrics: macro-F1 for roof-material imputation and building-function classification, MRR for \texttt{ENRICHED\_BY} link prediction. Baseline results are forthcoming.}
% %   \label{tab:results_rl}
% %   \setlength{\tabcolsep}{6pt}
% %   \begin{tabular}{lccc}
% %     \toprule
% %     & Imputation & Node class. & Link pred. \\
% %     Protocol & (macro-F1) & (macro-F1) & (MRR) \\
% %     \midrule
% %     Provenance-agnostic & -- & -- & -- \\
% %     Provenance-weighted & -- & -- & -- \\
% %     \addlinespace
% %     $\Delta$ (weighted $-$ agnostic) & -- & -- & -- \\
% %     \bottomrule
% %   \end{tabular}
% % \end{table}

% \subsection{Representation Learning Benchmark}\label{sec:repL}

% % The provenance model of \benchname{}---canonical urban entities linked to source-specific evidence (authoritative CityGML, crowd-sourced OSM, and, for Hamburg, machine-learned roof materials and reconstructed LoD3 geometry) by confidence-weighted correspondence edges---makes \emph{source}, \emph{confidence}, and \emph{coverage} first-class signals rather than metadata to be discarded at load time. 
% % This section asks whether graph representation learning can exploit those signals: 
% We define three learning tasks over the graph, compare \emph{provenance-agnostic} and \emph{provenance-aware} embeddings, and report baseline results together with the downstream applications enabled by the learned representations. Building-height prediction (\textbf{T1}) is evaluated across all five cities. Roof-material imputation represents a cross-source completion variant of \textbf{T1} using Hamburg's machine-learned roof-material layer and is therefore evaluated only in Hamburg. Roof-type prediction (\textbf{T2}) is evaluated only in the two cities that provide this label, namely Hamburg and Helsinki. building-function classification (\textbf{T2}) has sufficient label coverage only in Hamburg. Cross-source matching (\textbf{T3}) is evaluated independently in each city. Section~\ref{sec:repL:setup} specifies exactly which cities contribute to each reported result, as this variation in label coverage is itself a characteristic of open, multi-source urban knowledge graphs.

% \subsection{Tasks (T1 -T3)}
% \label{sec:repL:tasks}
% We instantiate the representation-learning family as three tasks that probe
% complementary abilities.

% \paragraph{\textbf{T1 --- Multi-source attribute prediction.}} Predict attributes of a canonical entity from its multi-source evidence and graph context, including building height and storey count (regression, evaluated using MAE, RMSE, and $R^2$) as well as \emph{cross-source completion}, which predicts an attribute generated by one source or process for entities where it is unavailable. Evaluation is restricted to entities for which the authoritative value is absent and never replaces existing authoritative information. Our instance of cross-source completion is roof-material imputation in Hamburg. An external ML model~\citep{lukas_roofmaterials} predicts roof material from aerial imagery for approximately 50\% of buildings, and the task is to infer the roof-material class for the remaining buildings using only the graph structure. Performance is reported as macro-F1 (Section \ref{sec:repL:results}).

% \paragraph{\textbf{T2 --- Node classification.}} Assign semantic categories to canonical entities. We evaluate two classification tasks: canonical roof-type prediction, which is available only in the cities where the authoritative source provides this information (Hamburg and Helsinki), and building-function/function classification, for which sufficient label coverage exists only in Hamburg (see Section~\ref{sec:repL:setup}). As both tasks exhibit class imbalance, we report macro-F1.


% \paragraph{\textbf{T3 --- Cross-source matching.}} Predict whether candidate correspondences between source-specific evidence nodes, such as a CityGML building and an OSM feature, refer to the same canonical entity. We formulate this task as link prediction, or equivalently entity resolution, over confidence-weighted correspondence edges~\cite{wang2018gcnalign}. Negative examples include \emph{hard} non-matches consisting of spatially adjacent but distinct entities. We report ROC-AUC, average precision (AP), and Hits@$k$/MRR using per-source-node ranking.

% % \subsection{Provenance-agnostic vs.\ provenance-aware representations}
% \subsection{Provenance-agnostic vs.\ provenance-aware}
% \label{sec:repL:prov}
% We compare two protocols that differ only in how much of the provenance structure the encoder may use, so that any performance gap is attributable to provenance awareness rather than to model capacity.

% \paragraph{\textbf{Provenance-agnostic.}} The multi-source graph is flattened into canonical entities by merging source-specific evidence, removing edge-confidence values, and discarding source labels. This corresponds to the conventional setting encountered by generic graph models and serves as the basis for our homogeneous and heterogeneous GNN baselines.

% \paragraph{\textbf{Provenance-aware.}} The encoder retains \emph{(i) source-typed nodes and relations}, including CityGML, OSM, and ML-derived sources, \emph{(ii) confidence scores as edge weights during message passing}, and \emph{(iii) coverage and agreement features at the node level}, encoding which sources are present, the number of contributing sources, and pairwise attribute agreement. Specifically, a canonical entity $c$ aggregates evidence from its neighboring evidence nodes $e \in \mathcal{N}(c)$ as
% \begin{equation}
% \mathbf h_c' \;=\; \sigma\Big(\mathbf W_{\text{self}}\,\mathbf h_c
% \;+\!\!\sum_{e\in\mathcal N(c)}\!\frac{w_{ce}}{\textstyle\sum_{e'} w_{ce'}}\;
% \mathbf W_{s(e)}\,\mathbf h_e\Big),
% \label{eq:confagg}
% \end{equation}
% where $w_{ce}\in[0,1]$ denotes the correspondence confidence and $s(e)$ denotes the source of $e$, with $\mathbf W_{s(e)}$ representing the corresponding source-specific weight matrix. The agnostic variant sets $w_{ce}=1$ and $\mathbf W_{s(e)}=\mathbf W$, thereby recovering mean-pooled GraphSAGE~\citep{hamilton2017inductive}. Our implementation follows \cref{eq:confagg} directly. The value of $w_{ce}$ is the geometric overlap (Jaccard index) stored on the CityGML--OSM correspondence edges and is normalized per target node, resulting in a confidence-weighted mean aggregation. The source-specific projections $\mathbf W_{s(e)}$ are instantiated as relation-specific weight matrices, with one matrix per typed relation using \texttt{to\_hetero}. Coverage and agreement features, including source availability, source count, and best-match confidence, are appended to the canonical-entity feature vectors. Non-correspondence relations, such as building--POI proximity edges, are assigned unit weight because their edge attributes, such as distance, do not represent match confidence.


% \subsection{Models}
% \label{sec:repL:models}
% We evaluate models spanning two orthogonal dimensions of the encoder--decoder design, namely \textbf{shallow} versus \textbf{message-passing} methods and, independently, \textbf{provenance-agnostic} vs.\ \textbf{provenance-aware} representations.

% \begin{itemize}
%   \item \textbf{Attribute-only MLP}: node attributes only, without graph structure, providing a lower bound on the contribution of graph-based information.
%   \item \textbf{Shallow embeddings}: node2vec~\cite{grover2016node2vec} in the provenance-agnostic setting, and metapath2vec~\cite{dong2017metapath2vec} with provenance-typed metapaths, for example \textsc{Canonical}--\textsc{CityGML}--\textsc{Canonical} and \textsc{Canonical}--\textsc{OSM}--\textsc{Canonical}, evaluated using a frozen embedding probe~\citep{chen2020simple}.
%   \item \textbf{GNNs (provenance-agnostic)}: GraphSAGE~\cite{hamilton2017graphsage} and GAT~\cite{velickovic2018gat} on the flattened graph, providing strong, well-tuned provenance-agnostic baselines following the recommendations of~\citet{lv2021hgb}.
%   \item \textbf{GNNs (provenance-aware)}: Source-typed R-GCN~\cite{schlichtkrull2018rgcn}, HAN~\cite{wang2019han}, HGT~\cite{hu2020hgt}, and Simple-HGN~\cite{lv2021hgb}, together with a confidence-weighted variant implementing the aggregation scheme defined in \cref{eq:confagg}.
%   \item \textbf{Matching decoders}: a shared GNN encoder with a bilinear or dot-product link decoder. We compare provenance-agnostic and confidence-weighted encoders against non-learned baselines based on spatial-distance and attribute-similarity thresholds.
%   \item \textbf{Self-supervised}: DGI~\cite{velickovic2019dgi} pretrains an inductive GNN encoder using contrastive learning and is evaluated with frozen embeddings and a downstream probe, analogous to the shallow embedding methods but using a message-passing encoder instead of a per-node embedding table. GraphMAE~\cite{hou2022graphmae} is included for completeness but is not evaluated. Together, these methods provide a natural bridge toward few-shot cross-city transfer.
% \end{itemize}

% This list defines the complete benchmark design. The experiments reported below evaluate the attribute-only MLP, node2vec and metapath2vec (Section \ref{sec:repL:setup}), DGI, GraphSAGE, GAT, HGT, HAN, R-GCN, the confidence-weighted provenance-aware encoder defined in \cref{eq:confagg}, and both non-learned baselines for T3 across tasks T1--T3. Simple-HGN and GraphMAE are included in the benchmark specification but are not evaluated in this round.


% \subsection{Experimental setup}
% \label{sec:repL:setup}
% \paragraph{\textbf{Splits.}} We evaluate all models under three data-splitting strategies. The \emph{random} split is stratified and reported only as a reference because it is susceptible to leakage arising from spatial autocorrelation. The \emph{spatial} split mitigates this effect by partitioning each city into spatial tiles and holding out entire tiles, preventing nearby buildings with similar attributes from appearing in both training and test sets. Because a single spatial hold-out exhibits high variance at the scale of an individual city, for example when a tile contains few or no instances of a minority class, we employ \emph{spatial $K$-fold cross-validation} with $K=5$. Out-of-fold predictions are pooled before evaluation so that every entity is tested exactly once. Finally, the \emph{cross-city} split follows a leave-one-city-out protocol in which models are trained on all but one city and evaluated on the held-out city, measuring generalization across different urban morphologies and geographic regions. Building-height prediction (T1) uses all five cities, with each city serving as the test set in turn. Roof-type prediction (T2) supports only Hamburg--Helsinki transfer because these are the only two cities that provide roof-type labels (Section~\ref{sec:repL:results}). We therefore report this as a two-city experiment rather than forcing it into a five-city evaluation. For T3, correspondence edges are partitioned such that held-out positive edges are removed from the message-passing graph, preventing information leakage through the links being predicted. Negative examples include \emph{hard} non-matches consisting of spatially adjacent but distinct entities. T3 is evaluated independently for each city rather than in a cross-city setting (Section~\ref{sec:repL:discussion}). All split definitions are fixed and shared across models to ensure a controlled comparison between provenance-agnostic and provenance-aware methods.

% \paragraph{\textbf{Protocol.}} Shallow embedding methods are trained in an unsupervised manner, frozen, and evaluated using a downstream probe, whereas GNNs are trained end-to-end. Unless stated otherwise, we report the mean and standard deviation over \num{3} random seeds. The random-walk training used by node2vec and metapath2vec does not scale to million-node graphs with our implementation and batch size. Consequently, these methods are evaluated only at subset scale on Hamburg ($n=81$), while all GNN and attribute-only baselines are evaluated on the full city-scale graphs. Hyperparameters are selected using the validation set only. Provenance-agnostic and provenance-aware variants share the same backbone architecture, network depth, and training budget, including an embedding and hidden dimension of \num{64}, \num{2} message-passing layers, \num{4} attention heads for GAT and HGT, dropout of \num{0.3}, and Adam optimization with a learning rate of \num{0.01}, weight decay of $5$$\times$$10^{-4}$, and \num{150} training epochs. For node2vec and metapath2vec, we use a walk length of \num{20}, a context window of \num{7}, \num{10} walks per node, \num{50} training epochs, and a frozen-embedding logistic regression probe. Continuous targets are standardized during training and transformed back to their original scale for evaluation. To prevent target leakage, the predicted attribute is always removed from the input features. By default, its correlated counterpart, for example storey count when predicting building height, is retained because it represents a legitimate multi-source attribute rather than a derivation of the prediction target. A dedicated ablation study (Section~\ref{sec:repL:results}) additionally removes this feature to isolate the contribution of the graph structure. Regression tasks use an MSE loss, whereas the matching task uses a dot-product decoder optimized with binary cross-entropy. Because New York City and Zurich do not provide authoritative building-height attributes, the target is derived instead from the vertical extent of the LoD2 building solid. The same quantity is never used simultaneously as both an input feature and a prediction target, preserving the leakage guard. 
% \\\\
% All models are implemented in PyTorch Geometric\footnote{\url{https://pytorch-geometric.readthedocs.io/en/latest/}} and trained on an Nvidia RTX PRO~6000 workstation with 96GB of VRAM. Data are streamed directly from Neo4j, and the data loader automatically handles missing node types, features, and labels, allowing the same implementation to run unchanged across all cities. 
% % \textcolor{red}{We release both the implementation and the data splits.}

% \subsection{Experimental results}
% \label{sec:repL:results}
% We report results at two evaluation scales. The subset study, conducted on a Hamburg subset containing $n=81$ buildings (\cref{tab:repL-subset}), was used to validate the experimental pipeline before scaling to the complete benchmark. The full-scale results are presented in the remaining tables. \Cref{tab:repL-main} reports Hamburg results for T1 and T2, \cref{tab:repL-crosscity} reports five-city cross-city transfer for T1 together with Hamburg--Helsinki transfer for the T2 roof-type task, \cref{tab:repL-match} reports T3 cross-source matching for each city, and \cref{tab:repL-roofmat} reports the T1 cross-source completion task, namely roof-material imputation, which is evaluated only in Hamburg. Unless otherwise stated, full-scale results are averaged over three random seeds. The T1 cross-city experiments in \cref{tab:repL-crosscity} are averaged over ten seeds because they constitute a central empirical result of this work (Section~\ref{sec:repL:results}. Boldface indicates the best result within each column, except in \cref{tab:repL-crosscity}, where each row represents an independent held-out city and the best result is highlighted within that row.

% \paragraph{\textbf{Preliminary subset (Hamburg, $n{=}81$ buildings).}}
% \Cref{tab:repL-subset} summarizes a single-city subset study that validated the complete experimental pipeline, including spatial $K$-fold cross-validation and leakage-free evaluation, using the only label space available in the pre-full-dataset extract, namely the canonical CityGML \emph{roof type}. The original five ALKIS roof-type classes were collapsed into the categories {Flat, Gable-family, Other} because the minority classes were too rare to support stratified spatial splits. Three observations emerged from the subset study. Two were confirmed by the full-scale evaluation, whereas one was overturned. \emph{(i) Topology provides information beyond node attributes.} For building-height regression, removing the highly correlated storey-count feature caused the attribute-only model to perform worse than predicting the mean ($R^2=\num{-0.86}$). In contrast, GraphSAGE recovered substantial predictive performance from graph structure alone ($R^2=\num{0.52}$), demonstrating that neighborhood information contributes meaningful signal. \emph{(ii) Provenance awareness provides no measurable benefit at subset scale.} The confidence-weighted, source-typed encoder achieved the highest macro-F1 score for roof-type classification (\num{0.36} compared with \num{0.34} for provenance-agnostic GraphSAGE), but the difference remained within the variance across random seeds. More importantly, the provenance-aware model showed no improvement on the seven buildings for which OSM and CityGML provided conflicting roof evidence, directly testing the central hypothesis. We attribute this result to the limited scale of the subset, which contains too few conflicting cases and only \num{92} correspondence pairs for T3 to produce a reliable estimate. \emph{(iii) Cross-source matching is the most challenging task.} Message passing reduces performance because hard negative examples share much of the same spatial context as the corresponding positive pairs. Consequently, an attribute-and-position baseline outperforms the learned graph models in this setting. Overall, the subset study motivated the subsequent multi-city evaluation (Section~\ref{sec:repL:setup}).


% \begin{table}[t]
%   \centering
%   \caption{Preliminary results on the Hamburg subset ($n=81$; spatial $K$-fold cross-validation with $K=5$; mean over \num{3} random seeds). Roof type is evaluated using macro-F1 (majority-class baseline \num{0.29}), height using $R^2$ with the collinear storey feature withheld, and matching using ROC-AUC with hard negatives. "—" indicates that the experiment was not run.}
%   \label{tab:repL-subset}
%   \small\setlength{\tabcolsep}{5pt}
%   \begin{tabular}{@{}lccc@{}}
%     \toprule
%     & T2 Roof type & T1 Height & T3 Matching\\
%     Model & macro-F1 $\uparrow$ & $R^2$ $\uparrow$ & ROC-AUC $\uparrow$\\
%     \midrule
%     MLP (attr-only)             & 0.30 & $-0.86$ & \textbf{0.59}\\
%     node2vec (probe)            & 0.26 & --      & --\\
%     metapath2vec (prov.\ paths) & 0.30 & --      & --\\
%     \midrule
%     GraphSAGE (agnostic)        & 0.34 & \textbf{0.52} & 0.39\\
%     GAT (agnostic)              & 0.31 & unstable\textsuperscript{$\dagger$} & --\\
%     HGT (agnostic)              & 0.27 & 0.20 & --\\
%     Conf.-weighted GNN (aware)  & \textbf{0.36} & 0.45 & 0.46\\
%     \bottomrule
%   \end{tabular}\\[2pt]
%   {\footnotesize\raggedright\textsuperscript{$\dagger$}GAT is high-variance on
%   the small-$n$ regression ($R^2{=}-2.35\pm3.09$); GraphSAGE is the stable GNN
%   at this scale.\par}
% \end{table}

% \begin{table}[t]
%   \centering
%   \caption{Full-scale (single-city) Hamburg results (388k buildings; spatial cross-validation; mean over \num{3} seeds). Height prediction retains the collinear storey count. Building function uses the eight most frequent classes (Section~\ref{sec:repL:setup}). HGT, HAN, and R-GCN are source-typed but not confidence-weighted.}
%   \label{tab:repL-main}
%   \small\setlength{\tabcolsep}{4pt}
%   \begin{tabular}{@{}lccc@{}}
%     \toprule
%     & Roof type & Building use & Height\\
%     Model & macro-F1 $\uparrow$ & macro-F1 $\uparrow$ & $R^2$ $\uparrow$\\
%     \midrule
%     MLP (attr-only)             & 0.504 & 0.329 & 0.666\\
%     DGI (self-supervised)       & 0.516 & 0.409 & --\\
%     GraphSAGE (agnostic)        & 0.556 & 0.462 & 0.730\\
%     GAT (agnostic)              & 0.531 & --    & 0.676\\
%     \multicolumn{4}{@{}l}{\emph{source-typed, not confidence-weighted:}}\\
%     HGT                         & 0.541 & --    & 0.714\\
%     HAN                         & 0.460 & 0.208 & 0.416\\
%     R-GCN                       & 0.558 & 0.470 & 0.713\\
%     Conf.-weighted GNN (aware)  & \textbf{0.559} & \textbf{0.476} & \textbf{0.733}\\
%     \bottomrule
%   \end{tabular}
% \end{table}

% \begin{table}[t]
%   \centering
% \caption{T3 cross-source matching per city (higher is better; mean over \num{3} seeds; std $\leq$ \num{0.03} unless shown). \emph{spatial} and \emph{attrsim} are non-learned baselines, \emph{attr} is a no-graph encoder, and \emph{agnostic}/\emph{aware} are GraphSAGE encoders. $n$ denotes correspondence pairs, with 30\% held out for testing.}
%   \label{tab:repL-match}
%   \small\setlength{\tabcolsep}{3pt}
%   \begin{tabular}{@{}lrccccc@{}}
%     \toprule
%     City & $n$ & spatial & attrsim & attr & agnostic & aware\\
%     \midrule
%     Hamburg  & 364{,}684     & \textbf{0.967} & 0.539 & 0.553 & 0.212 & 0.269\\
%     Helsinki & 2{,}712       & \textbf{0.952} & 0.563 & 0.574 & 0.307 & 0.388\\
%     NYC      & 1{,}073{,}423 & \textbf{0.996} & 0.500 & 0.501 & 0.175 & 0.168\\
%     Tokyo    & 1{,}349{,}003 & \textbf{0.952} & 0.501 & 0.505 & 0.266 & 0.266\\
%     Zurich   & 97{,}818      & \textbf{0.974} & 0.494 & 0.535 & 0.267 & 0.300\\
%     \bottomrule
%   \end{tabular}
% \end{table}

% \begin{table}[t]
%   \centering
% \caption{T1 cross-city height transfer with leave-one-city-out evaluation ($R^2$; mean$\pm$std over \textbf{10} seeds). Bold indicates the best cross-city result per row. \emph{probe} is a no-graph pooled-data baseline, and \emph{own-city} is the single-city reference (Table~\ref{tab:repL-main}; Section~\ref{sec:repL:results}). Bottom: T2 roof type transfer on the only valid two-city pair (merged 3-class macro-F1; mean over \num{3} seeds).}
%   \label{tab:repL-crosscity}
%   \small\setlength{\tabcolsep}{4pt}
%   \begin{tabular}{@{}lrcccc@{}}
%     \toprule
%     Held-out city & $n$ & probe & agnostic & aware & own-city\\
%     \midrule
%     Hamburg  & 388{,}267   & $-0.145$ & $0.270\pm.170$ & $\mathbf{0.548\pm.071}$ & 0.733\\
%     Helsinki & 2{,}980     & $-0.000$ & $\mathbf{0.516\pm.016}$ & $0.500\pm.029$ & 0.461\\
%     NYC      & 1{,}083{,}437 & $\mathbf{0.258}$ & $0.167\pm.034$ & $0.188\pm.024$ & 0.407\\
%     Tokyo    & 2{,}005{,}762 & $-0.484$ & $0.011\pm.130$ & $\mathbf{0.082\pm.063}$ & 0.399\\
%     Zurich   & 102{,}668   & $-0.011$ & $0.581\pm.035$ & $\mathbf{0.601\pm.037}$ & 0.269\\
%     \bottomrule
%   \end{tabular}\\[4pt]
%   \begin{tabular}{@{}lcc@{}}
%     \toprule
%     Held-out (roof type) & agnostic GNN & aware GNN\\
%     \midrule
%     Hamburg  & 0.203 & \textbf{0.208}\\
%     Helsinki & \textbf{0.336} & \textbf{0.337}\\
%     \bottomrule
%   \end{tabular}
% \end{table}

% \paragraph{\textbf{Full run (Hamburg, $n{=}388k$ buildings).}}
% The full evaluation both confirms and revises the subset findings. Within a single city, the subset's observation that provenance awareness provides limited benefit remains valid at scale. The agnostic-to-aware performance difference remains small for Hamburg T1 height prediction ($\num{0.730}\rightarrow\num{0.733}$) and T2 roof-type classification ($\num{0.556}\rightarrow\num{0.559}$). The exception is T2 building-function classification, which has the richest label structure with 131 raw classes, where the provenance-aware model achieves the largest single-city improvement observed ($\num{0.462}\rightarrow\num{0.476}$). A closer analysis in \cref{tab:repL-main} explains this limited difference. R-GCN, which incorporates source-specific relations but does not use confidence weighting, performs within \num{0.02} of the full provenance-aware encoder across all three Hamburg tasks, which is roof type ($\num{0.558}$ vs.\ $\num{0.559}$), building function ($\num{0.470}$ vs.\ $\num{0.476}$), and height regression ($\num{0.713}$ vs.\ $\num{0.733}$). This indicates that source typing captures most of the single-city benefit, while confidence weighting provides only a marginal additional contribution in this setting. HAN is the exception as it underperforms even the attribute-only baseline for building-function classification ($\num{0.208}$ vs.\ $\num{0.329}$). This may result from a mismatch between HAN's intended semantic-path assumptions and our graph schema since PyG's HANConv constructs metapaths from direct one-hop relations, whereas the original HAN formulation relies on longer, semantically meaningful paths. Our schema may therefore not expose sufficient metapath structure for HAN to exploit. 
% % We report this result as observed rather than applying task-specific tuning that would obscure the comparison.

% On the T1 building-height task under the leave-one-city-out protocol, the provenance-aware encoder now matches or outperforms the provenance-agnostic baseline for every held-out city (\cref{fig:crosscity}). Performance improves from $\num{0.270}$ to $\num{0.548}$ for Hamburg, with the standard deviation decreasing from \num{0.170} to \num{0.071}, indicating not only higher accuracy but also greater stability. Further improvements are observed for Zurich ($\num{0.581}\rightarrow\num{0.601}$), New York City ($\num{0.167}\rightarrow\num{0.188}$), and, in contrast to the earlier three-seed experiment, Tokyo ($\num{0.011}\rightarrow\num{0.082}$). Helsinki shows no meaningful difference within the experimental variance ($\num{0.516}$ vs.\ $\num{0.500}$). These results suggest that a single-city evaluation provides too little cross-source disagreement and no distribution shift for source-typed, confidence-weighted aggregation to demonstrate a clear advantage, consistent with the null result on the seven-building disagreement subset in our Hamburg subset ($n=81$; Section~\ref{sec:repL:results}). In contrast, holding out an entire city introduces both cross-source disagreement and distribution shift, which is precisely the setting for which the aggregation scheme in \cref{eq:confagg} was designed.

% Interestingly, the two smallest cities by building count, Helsinki ($n=2{,}980$) and Zurich ($n=102{,}668$), benefit from training on the other four cities rather than on their own data alone. Performance improves from \num{0.461} to \num{0.500} for Helsinki and from \num{0.269} to \num{0.601} for Zurich, the largest improvement observed in the table. In contrast, the three largest cities, Hamburg ($388{,}267$ buildings), New York City ($1{,}083{,}437$), and Tokyo ($2{,}005{,}762$), all perform worse under cross-city transfer than when trained on their own data, for example Tokyo drops from \num{0.399} to \num{0.082}. This pattern suggests that data-scarce cities gain more from the additional training examples available in other cities than they lose from differences in urban morphology. Conversely, data-rich cities already contain sufficient local examples, making the distribution shift introduced by other cities a source of noise rather than additional signal. This observation supports the downstream application of cross-city transfer for low-resource urban knowledge graphs, where models pretrained on data-rich cities can be used to bootstrap prediction performance in cities with limited labeled data. Notably, this conclusion emerges directly from the primary experiments rather than from a dedicated transfer-learning study.

% For T3 cross-source matching, the provenance-aware encoder matches or outperforms provenance-agnostic GraphSAGE in four of the five cities, with the largest ROC-AUC improvements observed in Helsinki (\num{+0.081}) and Zurich (\num{+0.033}). However, this comparison proved not to be the most informative result for T3. A non-learned \emph{spatial-distance} baseline, which predicts correspondences solely from footprint proximity without any training, achieves ROC-AUC scores between \num{0.95} and \num{0.996} across all cities (\cref{tab:repL-match}). This substantially exceeds the performance of all learned models, including the attribute-only baseline, which remains below \num{0.57} in every case. In hindsight, this result is unsurprising because the centroids of matching building footprints are almost coincident, whereas the hard negatives generated by the k-d tree are spatially close but still distinct. Consequently, Euclidean distance alone nearly solves the task by construction. Confidence weighting improves the GNN encoder relative to the provenance-agnostic variant, but it does not alter the fundamental observation that none of the learned models approach the predictive power of a simple coordinate-based lookup (Section~\ref{sec:repL:discussion}).
% % (\cref{tab:repL-match}, \cref{fig:match-pr})

% \paragraph{Ablations.} The collinear-partner ablation, in which storey count is removed when predicting building height, replicates cleanly at full Hamburg scale and now includes the complete range of evaluated methods (\cref{tab:repL-ablation}), rather than only the attribute probe, GraphSAGE, and the provenance-aware encoder. Once the shortcut feature is removed, every graph-based model degrades less than the attribute-only probe, although the magnitude of the effect differs across architectures. The attribute-only probe exhibits the largest absolute decline ($\Delta=\num{-0.352}$) because it has no alternative source of predictive information. In contrast, GraphSAGE, R-GCN, and the provenance-aware encoder experience only modest reductions ($\Delta\approx\num{-0.13}$ to $\num{-0.15}$), indicating that they recover building height from neighborhood context rather than from the withheld attribute. GAT represents an intermediate case. At full scale, its performance remains stable ($\Delta=\num{-0.295}$), in sharp contrast to the highly variable subset result ($R^2=\num{-2.35\pm3.09}$ at $n=\num{81}$), but it remains more dependent on the shortcut feature than the other GNNs. This suggests that its attention mechanism relies more heavily on the readily available correlated attribute than mean-aggregation methods do. HAN, by contrast, is almost unaffected by the ablation ($\num{0.416}\rightarrow\num{0.411}$, $\Delta=\num{-0.005}$), consistent with the results in \cref{tab:repL-main}, which indicate that it does not effectively exploit the graph schema. A decomposition of the provenance-aware encoder into its three constituent components, namely confidence weighting, coverage features, and source typing, is currently available only at subset scale. There, coverage features improved the featureless probe but not the GNN, while the full confidence-weighted aggregation produced the only clear single-city improvement under the leakage-controlled (spatial cross-validation) evaluation. Repeating this three-way ablation on the full multi-city benchmark remains future work.
% % \textcolor{red}{As a result, we conclude that a model that does not make effective use of the shortcut has little performance to lose when it is removed.}

% \begin{table}[t]
%   \centering
%   \caption{Collinear-partner ablation on Hamburg height (spatial cross-validation; $R^2$; mean over \num{3} seeds). \emph{With} retains the storey count, \emph{without} removes it. $\Delta$ = without $-$ with.}
%   \label{tab:repL-ablation}
%   \small\setlength{\tabcolsep}{4pt}
%   \begin{tabular}{@{}lccc@{}}
%     \toprule
%     Model & with & without & $\Delta$\\
%     \midrule
%     Mean baseline                & $-0.017$ & $-0.017$ & $0.000$\\
%     MLP (attr-only)              & 0.666 & 0.314 & $-0.352$\\
%     GraphSAGE (agnostic)         & 0.730 & 0.582 & $-0.148$\\
%     GAT (agnostic)               & 0.676 & 0.381 & $-0.295$\\
%     HGT                          & 0.714 & 0.549 & $-0.165$\\
%     HAN                          & 0.416 & 0.411 & $\mathbf{-0.005}$\\
%     R-GCN                        & 0.713 & 0.580 & $-0.133$\\
%     Conf.-weighted GNN (aware)   & 0.732 & \textbf{0.604} & $-0.128$\\
%     \bottomrule
%   \end{tabular}
% \end{table}


% \paragraph{Task: roof-material imputation (Hamburg only).} This experiment provides the concrete instance of the \emph{cross-source completion} task introduced in Section~\ref{sec:repL:tasks}. An external ML model, which is not part of the representation-learning benchmark, predicts roof-material classes ({roof\_tiles, tar\_paper, concrete, metal, glass}) for \num{194799} of Hamburg's \num{388267} buildings, corresponding to \num{50.2}\% coverage. This constitutes a genuine transductive setting because both covered and uncovered buildings belong to the same graph, and coverage is determined by the upstream process rather than by random sampling. We treat the ML-generated roof-material labels as the prediction target and evaluate whether graph-based representations can recover these labels for the buildings on which the external model was not applied.


% \paragraph{Circularity check.} Because the roof-material prediction model and the OSM fusion layer both encode information related to roof appearance, a graph encoder could, in principle, exploit duplicated information rather than learning meaningful structural representations. We perform two checks to assess this possibility. First, the feature sets for \texttt{Building} and \texttt{OsmBuilding} do not include the OSM roof-\emph{material} tag, only the distinct roof-\emph{shape} tag. Moreover, only \num{6873} of the \num{194799} labeled buildings have a matched OSM roof-material tag, corresponding to just \num{3.5}\% of the dataset. Consequently, a direct duplicate-label pathway is available for only a very small fraction of the graph. Second, if the graph encoder were primarily exploiting duplicated labels, graph-based methods would be expected to achieve near-perfect performance. Instead, the attribute-only baseline, which does not use graph structure or OSM information, already achieves a macro-F1 score of \num{0.283}, while the best graph-based model reaches only \num{0.316}, an absolute improvement of \num{0.033} (\cref{tab:repL-roofmat}). Although we cannot determine from the graph alone which data were used to train the external roof-material model, we can conclude that our encoder does not have direct access to duplicate roof-material labels, and its empirical behavior is consistent with learning from graph structure rather than exploiting label leakage. Per-class analysis reveals the source of this modest improvement. The two majority classes, roof\_tiles (\num{55.8}\%) and tar\_paper (\num{26.3}\%), are predicted reasonably well by all methods, including the attribute-only probe, with F1 scores ranging from \num{0.57} to \num{0.82}. The largest contribution from graph structure appears for the concrete class (\num{11.8}\%), where the provenance-aware encoder improves F1 from \num{0.01} for the probe to \num{0.10}. This indicates a genuine structural benefit for a class that cannot be effectively separated using node attributes alone. The minority classes metal (\num{4.4}\%) and glass (\num{1.7}\%) remain effectively unpredictable across all methods, with F1 scores close to zero, likely due to their limited representation at the available label density. HAN again underperforms the attribute-only probe (\num{0.241} vs.\ \num{0.283}), marking the third evaluated task in which the model appears poorly matched to this graph schema (\cref{tab:repL-main}).


% \begin{table}[t]
%   \centering
% \caption{Roof-material imputation on Hamburg (spatial cross-validation; macro-F1; mean over \num{3} seeds). MLP baselines exclude OSM features, while graph models use matched OSM features via message passing.}
%   \label{tab:repL-roofmat}
%   \small\setlength{\tabcolsep}{4pt}
%   \begin{tabular}{@{}lc@{}}
%     \toprule
%     Model & macro-F1 $\uparrow$\\
%     \midrule
%     MLP (attr-only, no OSM)      & 0.283\\
%     MLP + coverage stats         & 0.286\\
%     GraphSAGE (agnostic)         & 0.314\\
%     GAT (agnostic)                & 0.286\\
%     HGT                          & 0.296\\
%     HAN                          & 0.241\\
%     R-GCN                       & 0.314\\
%     DGI (self-supervised)        & 0.304\\
%     Conf.-weighted GNN (aware)   & \textbf{0.316}\\
%     \bottomrule
%   \end{tabular}
% \end{table}

% \paragraph{Task: storey count prediction.} The \texttt{storeys\_above\_ground} attribute is available only in Hamburg and Tokyo (Section~\ref{sec:repL:setup}). We therefore evaluate it as a separate T1 target rather than including it in \cref{tab:repL-main}, since the results are not directly comparable across the two cities. Hamburg represents the easier case, with the attribute-only probe achieving an $R^2$ of \num{0.624}, GraphSAGE reaching \num{0.832}, and the provenance-aware encoder achieving \num{0.838}. The strong performance of the probe indicates that storey count is already highly correlated with the other available Hamburg attributes, while graph structure provides a further consistent improvement. Tokyo presents a substantially more challenging prediction problem. The attribute-only probe achieves an $R^2$ of only \num{0.110}, whereas GraphSAGE and the provenance-aware encoder improve performance to \num{0.435} and \num{0.420}, respectively. Although graph structure approximately quadruples the predictive performance relative to the probe, the overall performance ceiling remains considerably lower than in Hamburg. A plausible explanation is that Tokyo's building stock is more uniformly low-rise (Section~\ref{sec:repL:setup}), leaving less target variance for any model to explain. This auxiliary task reinforces the conclusion drawn from building-height prediction: graph structure consistently improves performance, but the amount of recoverable signal depends not only on the model but also on the characteristics of the city.

% % \begin{figure}[t]
% %   \centering
% %   \includegraphics[width=0.95\linewidth]{figures/fig_prov_results.pdf}
% % \caption{Performance difference between the provenance-aware and agnostic encoders across tasks and evaluation settings (positive = aware performs better). Cross-city height shows the largest and most consistent improvement.}
% %   \label{fig:prov-results}
% % \end{figure}

% % \begin{figure*}[t]
% %   \centering
% %   \includegraphics[width=0.95\linewidth]{figures/fig_matching_pr.pdf}
% % \caption{T3 precision--recall curves with hard negatives, one panel per city. The spatial-distance baseline consistently outperforms the learned models, while the attribute-only, agnostic, and aware encoders remain close to the base rate (cf.~\cref{tab:repL-match}).}
% %   \label{fig:match-pr}
% % \end{figure*}

% \begin{figure}[t]
%   \centering
%   \includegraphics[width=0.95\linewidth]{figures/fig_crosscity.pdf}
% \caption{T1 cross-city height transfer under leave-one-city-out evaluation ($R^2$; 10 seeds). The aware encoder matches or outperforms the agnostic encoder in every city, with the largest improvement on Hamburg.}
%   \label{fig:crosscity}
% \end{figure}

% % \begin{figure}[t]
% %   \centering
% %   \includegraphics[width=\linewidth]{figures/fig_prov_tsne.pdf}
% % \caption{2D t-SNE visualization of Hamburg building embeddings learned by the T3 encoders. Columns show roof type, OSM coverage, and roof-material coverage (grey = unavailable). Compared with the agnostic encoder, the aware encoder more clearly separates OSM-covered buildings and yields a more structured embedding space.}
% %   \label{fig:prov-tsne}
% % \end{figure}

% \subsection{Downstream tasks and applications}
% \label{sec:repL:downstream}
% Beyond the three benchmark tasks (T1--T3), provenance-aware urban representations support a range of downstream applications. Cross-city transfer is one example directly supported by experimental evidence in this study, while the remaining applications represent motivated directions for future work.
% %machen wir die Tasks alle fett im text?
% %kuerzen wir Figure, Table und Sections in Klammern ab?
% % where cross-city transfer (\cref{fig:prov-results})

% \begin{itemize}
% \item \textbf{Data fusion and map conflation.} T3 directly models the automated matching of authoritative and crowd-sourced city models, replacing brittle rule-based conflation with confidence-aware entity resolution.

% \item \textbf{Quality assurance and error detection.} Low-confidence predictions and disagreements between sources can identify likely data quality issues, such as errors in OSM, enabling provenance-aware anomaly detection.

% \item \textbf{LoD enrichment and geometry completion.} Predicting semantic attributes, such as storey count or façade characteristics, from LoD2 models augmented with OSM information enables data-driven enrichment of sparse city models.

% \item \textbf{Provenance-respecting attribute completion.} Missing authoritative attributes can be inferred from crowd-sourced and ML-derived evidence without overwriting existing authoritative values, thereby improving coverage while preserving provenance. The T1 roof-material prediction task is a direct instance of this application.

% % \item \textbf{Coverage-guided data acquisition.} Predictive uncertainty and sparse source coverage can identify areas where additional surveys or street-level data collection would most improve the completeness of the digital twin.

% \item \textbf{Cross-city transfer for data-scarce urban models.} Models pretrained on data-rich cities can be transferred to cities with limited authoritative data, providing an effective strategy for bootstrapping high-detail semantic 3D city models and digital twins in data-scarce environments.

% \item \textbf{Embedding-based urban analytics.} Learned representations support downstream analyses including energy-efficiency estimation, urban heat-island assessment, solar-suitability analysis, functional zoning through embedding clustering, and nearest-neighbor retrieval of similar buildings across cities.
% \end{itemize}

% \subsection{Discussion and limitations}
% \label{sec:repL:discussion}
% % This study posed three questions but could only answer the first. The full five-city evaluation resolves the remaining two, albeit in different directions. 
% Topology provides useful information beyond node attributes when the available attributes are not enough. At full Hamburg scale, removing the correlated storey-count feature reduces the attribute-only probe to \num{0.314}, while GraphSAGE still reaches \num{0.582}, showing that the graph structure provides additional signal. Provenance-aware representations help when source differences matter. Within a single city, however, source disagreements are too rare to create clear improvements. The small gains observed there mainly come from source typing rather than confidence weighting, as R-GCN performs within \num{0.02} of the full provenance-aware encoder across all Hamburg tasks (\cref{tab:repL-main}). Confidence weighting becomes more important in cross-city transfer, where models face different data distributions. Across ten seeds, the confidence-weighted, source-typed encoder matches or outperforms the provenance-agnostic baseline in every held-out city (\cref{fig:crosscity}), improving Hamburg ($\num{0.270}\rightarrow\num{0.548}$) while also reducing variance by 2.4-fold.

% % (\cref{tab:repL-match}, \cref{fig:match-pr})
% Cross-source matching is mainly limited by the task setup rather than the model. A simple spatial-distance baseline already achieves ROC-AUC scores between \num{0.95} and \num{0.996} across all cities (\cref{tab:repL-match}), leaving little room for learned methods to improve. In this setting, message-passing models generally perform worse. This effect becomes stronger at larger scale because more data also means more spatially similar but incorrect candidate matches. Confidence weighting provides a small but consistent improvement over provenance-agnostic message passing, but it does not approach the performance of the spatial baseline. These results suggest that future T3 models should focus on pairwise features that capture cross-source agreement, or use a hybrid approach where spatial proximity handles obvious matches and learned models focus on ambiguous cases.

% % Qualitatively, the provenance-aware encoder separates OSM-matched buildings into a distinct embedding region, whereas the provenance-agnostic encoder does not. The visualization also reveals that ML and OSM coverage gaps overlap, indicating that the two missingness patterns are correlated rather than independent. This structural relationship is not captured by the classification metrics alone.

% However, several limitations remain. Derived data products, including ML-predicted roof material and reconstructed LoD3 geometry, are currently available only for Hamburg. Two semantic targets that were expected to generalize across cities do not provide sufficient coverage. Roof type is available in only two of the five cities, Hamburg and Helsinki, while building function/use has usable coverage only in Hamburg. For example, Helsinki contains labels for only 94 of its \num{2980} buildings, which is too sparse for reliable evaluation. These results show that open multi-source city models, even across CityGML datasets from the same continent, are not uniformly annotated. A benchmark based on such data should make these differences explicit rather than silently focusing on the best-covered city.

% The shallow embedding methods node2vec and metapath2vec are transductive, preventing them from embedding previously unseen entities. DGI does not share either limitation. It is an inductive encoder, was evaluated successfully at full city scale (\cref{tab:repL-main}), and consistently outperformed the attribute-only baseline on every single-city task without using labels during pretraining, for example on building-function classification (\num{0.409} versus \num{0.329}). However, it remains below the fully supervised GraphSAGE baseline (\num{0.462}), and we did not evaluate it under the cross-city protocol in this study. Extending DGI, together with GraphMAE, which is included in the benchmark specification but not yet evaluated, to the cross-city setting is a natural next step toward few-shot transfer for data-scarce cities.

% For T3, negative correspondences are generated heuristically as spatial hard negatives. More fundamentally, embedding-based matching is intrinsically constrained because the geometric overlap that defines a correspondence cannot also be provided as a decoder feature without introducing circularity. Finally, one practical lesson emerged during data integration. Tokyo encodes missing height and storey values using numeric sentinel values ($\pm9999$) rather than nulls. When combined with the other cities, these placeholders silently corrupted every regression task until detected (Section~\ref{sec:repL:setup}). Missing attributes represented as null values are straightforward to identify, whereas implausible numeric placeholders are much more difficult to detect. Any pipeline that ingests third-party open city data should therefore include explicit checks for sentinel values as well as missing attributes.

% %%%%%%%%%%%%%%%%%%%%%%%%%%%%%%%%%%%%%%%%%%%%%%%%%%%%%%%%%%%%%%%%%%%%
% %%%%%%%%%%%%%%%%%%%%%%%%%%%%%%%%%%%%%%%%%%%%%%%%%%%%%%%%%%%%%%%%%%%%
% %%%%%%%%%%%%%%%%%%%%%%%%%%LUKAS STUFF%%%%%%%%%%%%%%%%%%%%%%%%%%%%%%%

% %%%%%%%%%%%%%%%%%%%%%%%%%%LUKAS STUFF (KURZ)%%%%%%%%%%%%%%%%%%%%%%%%
% %%%%%%%%%%%%%%%%%%%%%%%%%%%%%%%%%%%%%%%%%%%%%%%%%%%%%%%%%%%%%%%%%%%%
% %%%%%%%%%%%%%%%%%%%%%%%%%%%%%%%%%%%%%%%%%%%%%%%%%%%%%%%%%%%%%%%%%%%%

\section{Benchmark B: Representation Learning}\label{subsec:rl_tasks}

To show that the graph supports applications beyond text-to-query, from urban analytics such as energy efficiency~\citep{Yap2025} and urban planning~\citep{LIU2022102936} to provenance-aware uses such as data fusion, quality assurance, and cross-city transfer, we provide an embedding-based evaluation following the protocol conventions~\citep{dwivedi_benchmarkgnn_2023}. It comprises three tasks under spatial-block and cross-city splits, namely \emph{attribute imputation} of ML-predicted roof material (available for only ${\sim}\num{50}\%$ of buildings~\citep{lukas_roofmaterials}) and authoritative building height (T1), \emph{node classification} of administrative building-function and roof-type classes (T2), and \emph{link prediction} of held-out \texttt{ENRICHED\_BY} correspondence edges (T3). Each task is run under a \emph{provenance-agnostic} protocol that hides source distinctions and a \emph{provenance-aware} one that exposes source types and fusion-edge confidences, for example through Jaccard-weighted message passing, and the difference between them measures whether provenance helps. To our knowledge this is the first such evaluation on a real city-scale multi-source graph, although label coverage varies by task, with height spanning all five cities, matching per city, roof type in Hamburg and Helsinki, and building function only Hamburg (Section~\ref{sec:repL:setup}).

% To highlight that the graph supports applications beyond text-to-query tasks, from analytical uses such as energy efficiency, environmental sustainability, and urban planning~\citep{LIU2022102936} to provenance-aware uses such as data fusion and confidence-aware conflation, quality assurance, level-of-detail enrichment, provenance-respecting attribute completion, and cross-city transfer for data-scarce cities, we provide an embedding-based evaluation following~\citet{dwivedi_benchmarkgnn_2023} for heterogeneous urban KGs. The evaluation comprises three tasks, each conducted under fixed, reproducible random, spatial-block, and cross-city splits. The first is \emph{attribute imputation}, predicting roof-material labels for buildings where they are unavailable. Roof material is available for only ${\sim}\num{50}\%$ of buildings and is provided by an ML model trained on crowd-sourced OSM labels~\citep{lukas_roofmaterials}. The second is \emph{node classification}, predicting administrative building-function and roof-type classes. These two roof targets span the trust spectrum, from authoritative roof-type labels to ML-predicted, OSM-derived roof material, so only the latter requires a circularity check (Section~\ref{sec:repL:results}). The third is \emph{link prediction}, predicting held-out \texttt{ENRICHED\_BY} edges to learn footprint correspondence~\cite{naumann_nnmatching_2025}. Each task is evaluated under both a \emph{provenance-agnostic} protocol, which hides source distinctions, and a \emph{provenance-aware} protocol, which exposes source types and fusion-edge confidence scores, for example by weighting message passing with the Jaccard index. The key result is the performance difference between these protocols, measuring whether provenance information improves representation learning. To the best of our knowledge, this is the first such evaluation on a real city-scale multi-source graph. Label coverage itself varies by task, with building height \textbf{(T1)} spanning all five cities, cross-source matching \textbf{(T3)} independently per city, roof type \textbf{(T2)} available in Hamburg and Helsinki, and building function \textbf{(T2)} only in Hamburg (Section~\ref{sec:repL:setup}).

% We position provenance-aware representation learning as an open problem for which our \benchname{} provides both the benchmark substrate and an initial empirical data point (\cref{tab:repL-main},~\cref{tab:repL-crosscity}), rather than a complete solution.
% This variation in label coverage is itself a characteristic of open, multi-source urban knowledge graphs, and Section~\ref{sec:repL:setup} specifies which cities contribute to each result.

% Building height \textbf{(T1)} is evaluated across all five cities. Roof-material imputation, a cross-source-completion variant of \textbf{T1} using Hamburg's machine-learned roof-material layer, and building-function classification \textbf{(T2)} are evaluated only in Hamburg. Roof type \textbf{(T2)} is evaluated in Hamburg and Helsinki, and cross-source matching \textbf{(T3)} independently per city. This variation in label coverage is itself a characteristic of open, multi-source urban knowledge graphs, and Section~\ref{sec:repL:setup} specifies which cities contribute to each result.

% \textcolor{blue}{Den Einführungstext kompakter machen, auf 50 Prozent Länge oder mehr. Verwende KGs, LPGs, anstatt Knowledge Graphs, usw.}

\subsection{Provenance-Agnostic vs.\ Provenance-Aware}
\label{sec:repL:prov}
The two protocols differ only in how much of the provenance structure the encoder may use, so that any performance gap is attributable to provenance awareness rather than to model capacity. The \emph{provenance-agnostic} setting flattens the multi-source graph into canonical entities by merging source-specific evidence, removing edge-confidence values, and discarding source labels. The \emph{provenance-aware} encoder instead retains (i) source-typed nodes and relations (CityGML, OSM, ML-derived), (ii) confidence scores as edge weights, and (iii) coverage and agreement node features (source presence, count, and best-match confidence). Specifically, a canonical entity $c$ aggregates evidence from its neighbors $e\in\mathcal N(c)$ as
\begin{equation}
\mathbf h_c' \;=\; \sigma\Big(\mathbf W_{\text{self}}\,\mathbf h_c
\;+\!\!\sum_{e\in\mathcal N(c)}\!\frac{w_{ce}}{\textstyle\sum_{e'} w_{ce'}}\;
\mathbf W_{s(e)}\,\mathbf h_e\Big),
\label{eq:confagg}
\end{equation}
where $w_{ce}\in[0,1]$ is the correspondence confidence (the Jaccard overlap stored on the CityGML--OSM edges, normalized per target node) and $\mathbf W_{s(e)}$ is a source-specific projection, instantiated as one relation-specific weight matrix per typed relation. The agnostic variant sets $w_{ce}=1$ and $\mathbf W_{s(e)}=\mathbf W$, recovering mean-pooled GraphSAGE~\citep{hamilton2017graphsage}.

\subsection{Experimental Setup}
\label{sec:repL:setup}
\paragraph{Models.} At full scale we evaluate an attribute-only MLP, provenance-agnostic GraphSAGE~\cite{hamilton2017graphsage} and GAT~\cite{velickovic2018gat}, source-typed R-GCN~\cite{schlichtkrull2018rgcn}, HAN~\cite{wang2019han}, and HGT~\cite{hu2020hgt}, the confidence-weighted encoder of \cref{eq:confagg}, self-supervised DGI~\cite{velickovic2019dgi}, and non-learned spatial-distance and attribute-similarity baselines for T3. The shallow node2vec~\cite{grover2016node2vec} and metapath2vec~\cite{dong2017metapath2vec} baselines are transductive and cannot embed unseen entities, so we exclude them from the full-scale comparison.

\paragraph{Splits and protocol.} The \emph{spatial} split holds out entire spatial tiles, and because a single hold-out is high-variance at city scale we use \emph{spatial $K$-fold cross-validation} ($K{=}5$) with out-of-fold predictions pooled so every entity is tested once. The \emph{cross-city} split is leave-one-city-out. Shallow methods are trained unsupervised, frozen, and probed, whereas GNNs are trained end-to-end. We report mean$\pm$std over \num{3} seeds for the survey (\Cref{tab:repL-main}) and over \num{10} seeds for the cross-city result. For the focused provenance-agnostic versus provenance-aware paired comparisons and the confidence-weighting ablation, we repeat over \num{10} seeds and pair runs by seed, reporting a paired $t$-test~\citep{gosset1908}, the Wilcoxon signed-rank test~\citep{wilcoxon1945}, Cohen's $d_z$~\citep{lakens2013}, and a 95\% confidence interval on the per-seed difference, and we treat a comparison as a null when the interval contains zero or the effect size is negligible. Provenance-agnostic and aware variants share backbone, depth, and budget (hidden dimension \num{64}, \num{2} message-passing layers, dropout \num{0.3}, Adam at learning rate \num{0.01}, \num{150} epochs), and are trained in PyTorch Geometric on an Nvidia RTX PRO~6000 (96\,GB). To prevent leakage, each predicted attribute is removed from the inputs, and a dedicated ablation additionally removes its correlated counterpart, such as storey count when predicting height. Data are streamed from Neo4j, with the loader handling missing node types, features, and labels so the same code runs unchanged across cities.

\subsection{Experimental Results}
\label{sec:repL:results}
A single-city subset (Hamburg, $n=81$ buildings) first validated the pipeline and previewed both headline findings, that topology adds signal beyond attributes (with the collinear storey feature withheld, the attribute probe falls to $R^2{=}\num{-0.86}$ while GraphSAGE recovers $\num{0.52}$) and that provenance-awareness gives no measurable single-city benefit, including no gain on the seven buildings with conflicting OSM/CityGML roof evidence.

\paragraph{Full run (Hamburg, $n{=}388$k buildings).}
Within a single city, the aware encoder's advantage over the agnostic baseline is statistically consistent but negligible in magnitude. Over \num{10} seeds it improves T1 height ($\num{0.729}\rightarrow\num{0.733}$), T2 roof type ($\num{0.556}\rightarrow\num{0.558}$), and T2 building function ($\num{0.464}\rightarrow\num{0.474}$), each significant under a paired $t$-test ($p<\num{0.02}$) yet at most \num{0.010} in absolute terms (cf.~\Cref{tab:repL-main}, \num{3}-seed run). R-GCN, source-typed but not confidence-weighted, stays within about \num{0.02} of the full encoder on all three tasks, which points to source typing rather than confidence weighting as the origin of even this small effect. We therefore state a falsifiable hypothesis: source typing and confidence weighting are separable, and confidence weighting is redundant with source typing in-distribution but not under distribution shift. We test it with an ablation that removes only the confidence weights, leaving source typing, coverage features, and architecture unchanged. In-distribution, removing confidence weighting changes every task by at most \num{0.001}, with no significant effect on height ($p=\num{0.21}$) or building function ($p=\num{0.59}$) and only a negligible \num{0.0006} on roof type ($p=\num{0.02}$), confirming that confidence weighting is redundant with source typing in-distribution. Whether it also contributes under distribution shift, where the aware encoder's cross-city gains appear (\Cref{tab:repL-crosscity}), requires the same ablation under leave-one-city-out and remains future work. In the T3 matching embedding space (\cref{fig:prov-tsne}), the aware encoder separates OSM-matched buildings into a distinct region where the agnostic encoder does not, and the ML- and OSM-coverage gaps co-locate, indicating that the two missingness patterns are correlated rather than independent, a restructuring the supervised metrics alone do not reveal. Under leave-one-city-out on T1 height the aware encoder improves on four of the five held-out cities (\cref{tab:repL-crosscity}), with Hamburg rising from \num{0.270} to \num{0.548} and its standard deviation dropping from \num{0.170} to \num{0.071}, plus gains on Zurich, New York City, and Tokyo, and Helsinki tied within noise. Roof-type transfer, possible only on the Hamburg--Helsinki pair, shows no meaningful gap between the two encoders (\Cref{tab:repL-crosscity-roof}). For T3 matching, every learned encoder falls far short of non-learned baselines, for a structural reason. \texttt{ENRICHED\_BY} correspondences are defined by footprint overlap (median Jaccard \num{0.842}), so on the default nearest-neighbor negatives a spatial-distance rule reaches ROC-AUC \num{0.95}--\num{0.996} while the encoders reach at most \num{0.57}. On the ambiguous n:m subset, where a building overlaps several OSM candidates and distance is uninformative, the distance rule falls to \num{0.75} AUC, a trivial attribute rule (building height versus OSM levels) still reaches \num{0.94}, and every learned encoder collapses to chance (\num{0.47}--\num{0.49} AUC on \num{18234} Hamburg cases). T3 is therefore not a coordinate lookup, but it exposes a concrete gap, namely that current encoders exploit geometry and ignore the cross-source attribute signal that actually disambiguates correspondences. We provide the distance and attribute heuristics as reference baselines and pose attribute-aware n:m matching as an open challenge.

\begin{table}[t]
  \centering
  \caption{Full-scale (single-city) Hamburg results (388k buildings, spatial cross-validation, mean$\pm$std over \num{3} seeds). Height retains the collinear storey count.}
  \label{tab:repL-main}
  \small
  \setlength{\tabcolsep}{4pt}
  \begin{tabular}{@{}lccc@{}}
    \toprule
    & Roof type & Bldg function & Height\\
    Model & macro-F1 $\uparrow$ & macro-F1 $\uparrow$ & $R^2$ $\uparrow$\\
    \midrule
    MLP (attr-only)             & $0.504\pm.000$ & $0.329\pm.000$ & $0.666\pm.000$\\
    DGI (self-supervised)       & $0.516\pm.007$ & $0.409\pm.002$ & $0.683\pm.002$\\
    GraphSAGE (agnostic)        & $0.556\pm.000$ & $0.462\pm.000$ & $0.730\pm.002$\\
    GAT (agnostic)              & $0.531\pm.000$ & $0.365\pm.001$ & $0.676\pm.001$\\
    \multicolumn{4}{@{}l}{\emph{source-typed, not confidence-weighted}}\\
    HGT                         & $0.541\pm.002$ & $0.393\pm.007$ & $0.714\pm.005$\\
    HAN                         & $0.460\pm.003$ & $0.208\pm.004$ & $0.416\pm.026$\\
    R-GCN                       & $0.558\pm.000$ & $0.470\pm.003$ & $0.713\pm.004$\\
    Conf.-weighted GNN (aware)  & $\mathbf{0.559\pm.001}$ & $\mathbf{0.476\pm.001}$ & $\mathbf{0.732\pm.000}$\\
    \bottomrule
  \end{tabular}
\end{table}

\begin{table}[t]
  \centering
  \caption{T1 cross-city height transfer, leave-one-city-out ($R^2$, mean$\pm$std over \textbf{10} seeds, best cross-city result per row in bold), where \emph{own-city} is the single-city reference. The \emph{probe} is a no-graph baseline pooled over the training cities and is not comparable to the single-city attribute probe of \Cref{tab:repL-main}.}
  \label{tab:repL-crosscity}
  \small
  \setlength{\tabcolsep}{4pt}
  \begin{tabular}{@{}lrcccc@{}}
    \toprule
    Held-out & $n$ & probe & agnostic & aware & own-city\\
    \midrule
    Hamburg  & 388{,}267     & $-0.145$ & $0.270\pm.170$ & $\mathbf{0.548\pm.071}$ & 0.733\\
    Helsinki & 2{,}980       & $-0.000$ & $\mathbf{0.516\pm.016}$ & $0.500\pm.029$ & 0.461\\
    NYC      & 1{,}083{,}437 & $\mathbf{0.258}$ & $0.167\pm.034$ & $0.188\pm.024$ & 0.407\\
    Tokyo    & 2{,}005{,}762 & $-0.484$ & $0.011\pm.130$ & $\mathbf{0.082\pm.063}$ & 0.399\\
    Zurich   & 102{,}668     & $-0.011$ & $0.581\pm.035$ & $\mathbf{0.601\pm.037}$ & 0.269\\
    \bottomrule
  \end{tabular}
\end{table}

\paragraph{Cross-source completion and auxiliary targets.} Unlike the three survey tasks, roof-material imputation, which predicts the ${\sim}\num{50}\%$ of Hamburg buildings the external ML model did not label, gives the aware encoder its clearest single-city edge (macro-F1 $\num{0.283}\rightarrow\num{0.316}$, concentrated in the \texttt{concrete} class). A circularity check confirms this is structural rather than leakage, as only \num{3.5}\% of labeled buildings carry a matched OSM roof-material tag.

% \subsection{Limitations \textcolor{blue}{Schau dir Appendix B an und ggfs. diesen Text reduzieren und dahin schieben / löschen}}
% \label{sec:repL:discussion}
% ML-derived layers (roof material, LoD3 geometry) exist only for Hamburg, and even across same-continent CityGML data the coverage is uneven, with roof type spanning two of five cities and building function only one, so a multi-source benchmark should make this explicit rather than report only the best-covered city. The shallow baselines are also transductive and cannot embed unseen entities.

% \subsection{Discussion and limitations}
% \label{sec:repL:discussion}
% Beyond the benchmark tasks, provenance-aware representations support data fusion and confidence-aware conflation (T3), quality assurance, LoD enrichment, provenance-respecting attribute completion (the roof-material task is a direct instance), and cross-city transfer for data-scarce cities. 
% However, several limitations remain. Derived layers (ML roof material, LoD3 geometry) exist only for Hamburg, and even across same-continent CityGML data the coverage is uneven, with roof type spanning two of five cities and building function only one, so a multi-source benchmark should make this explicit rather than report only the best-covered city. Additionally, the shallow baselines are transductive, whereas DGI is inductive and beats the attribute-only baseline on every single-city task without labels ($\num{0.409}$ vs.\ $\num{0.329}$ on building function) but was not run cross-city, an obvious next step toward few-shot transfer. 
% Finally, embedding-based matching is intrinsically bounded because the overlap that \emph{defines} a correspondence cannot be a decoder feature, and Tokyo's numeric sentinel heights ($\pm9999$) silently corrupt regression until detected, so ingestion pipelines must check for implausible placeholders, not only nulls.

% %%%%%%%%%%%%%%%%%%%%%%%%%%%%%%%%%%%%%%%%%%%%%%%%%%%%%%%%%%%%%%%%%%%%
% %%%%%%%%%%%%%%%%%%%%%%%%%%%%%%%%%%%%%%%%%%%%%%%%%%%%%%%%%%%%%%%%%%%%
% %%%%%%%%%%%%%%%%%%%%%%%%%%LUKAS STUFF (KURZ)%%%%%%%%%%%%%%%%%%%%%%%%

\section{Conclusion}\label{sec:conclusion}

We presented \benchname{}, a multi-source, provenance-aware 3D city knowledge graph and a two-family benchmark on top of it. The dataset makes origin, confidence, and coverage of urban facts explicit and queryable across authoritative, crowd-sourced, ML-predicted, and reconstructed layers. The benchmark evaluates capabilities that are not captured by existing text-to-query or urban reasoning benchmarks, including spatial reasoning over indexed 3D geometry, cross-source comparison, coverage-aware aggregation, infeasibility detection, and provenance-aware representation learning. We release the complete artifact under open licenses, including the enriched LPGs, question suite, gold queries, data splits, loaders, evaluation harness, and datasheet, together with an archival DOI. We hope \benchname{} will serve both as a rigorous testbed for evaluating query-generation models on realistic, heterogeneous urban data and as a foundation for developing provenance-aware embedding methods, an important gap highlighted by our baseline results. 
The project is available at \url{https://github.com/hcu-cml/pykci/authenticity}.
% \lukascomment{Koennten wir noch explizit highlighten, dass sich die codes fuers embedding (splits, encoder, etc. dort auch befinden)?}

% \begin{acks}
% This research was supported by the project ``Next Generation City Networking'' (Grant No. 19DZ24004) at the Hanseatic Wireless Innovation Competence Center (HAWICC). The project is funded by the Federal Ministry of Transport via the German Center for Future Mobility (DZM).
% \end{acks}

%%
%% The next two lines define the bibliography style to be used, and
%% the bibliography file.
\bibliographystyle{ACM-Reference-Format}
\bibliography{sample-base}

\appendix

\begin{figure}[h!]
  \centering
  \begin{subfigure}{0.49\linewidth}
    \includegraphics[width=\linewidth]{figures/extent_hamburg.jpg}
    \caption{Hamburg, 1:84,460}
  \end{subfigure}\hfill
  \begin{subfigure}{0.49\linewidth}
    \includegraphics[width=\linewidth]{figures/extent_tokyo.jpg}
    \caption{Tokyo (23 wards), 1:46,760}
  \end{subfigure}\\[2pt]
  \begin{subfigure}{0.49\linewidth}
    \includegraphics[width=\linewidth]{figures/extent_nyc.jpg}
    \caption{New York, 1:91,455}
  \end{subfigure}\hfill
  \begin{subfigure}{0.49\linewidth}
    % previous (city-core capture): \includegraphics{figures/extent_zurich_core.jpg}, Zurich, 1:\num{30120}
    \includegraphics[width=\linewidth]{figures/extent_zurich_core.jpg} %extent_zurich
    \caption{Zurich, 1:120,480}
  \end{subfigure}\\[2pt]
  \begin{subfigure}{0.49\linewidth}
    \includegraphics[width=\linewidth]{figures/extent_helsinki.jpg}
    \caption{Helsinki (Kalasatama Digital Twins), 1:8734}
  \end{subfigure}
  \caption{
    Spatial extent of the five \benchname{} cities. 
    %The more than tenfold variation in scale reflects source-specific release boundaries, from Helsinki's inner-city core to Zurich's canton-wide coverage. 
    Basemap: \copyright~OpenStreetMap contributors.
    % Spatial extent of the five \benchname{} cities: one ground footprint per building, from the source CityGML over the OpenStreetMap basemap. Map scales differ by over an order of magnitude (Helsinki's 1:\num{8734} core versus Zurich's 1:\num{120480} canton-wide extent), reflecting how differently the portals delimit releases: a federal state (Hamburg), a mesh-tiled metropolis (Tokyo), five boroughs (New York), a canton-wide tile scatter (Zurich), and an inner-city core (Helsinki). Basemap: \copyright~OpenStreetMap contributors.
    }
  \Description{Five map panels showing building-footprint extents over OpenStreetMap for Hamburg, Tokyo, New York, Zurich, and Helsinki. The maps use different scales, with Helsinki shown at the most zoomed-in scale and Zurich at the broadest extent.}
  \label{fig:city_extents}
\end{figure}

\section{Datasheet}\label{app:datasheet}

We document \benchname{} following the datasheets framework of~\citet{gebru_datasheets_2021}. Per-city figures referenced below are given in \Cref{tab:city_sources,tab:city_graph_content}. The cities' spatial extents are shown in~\Cref{fig:city_extents}. 

\subsection{Motivation} 

\dsq{For what purpose was the dataset created?} To provide a multi-source, provenance-aware 3D city knowledge graph and a benchmark that evaluates reasoning over source origin, confidence, coverage, and cross-source agreement, which existing text-to-query and urban benchmarks do not exercise (\Cref{sec:related_work}). 

\dsq{Who created it?} \benchname{} was created by the authors of this paper as part of an academic research effort on multi-source urban data integration. 
% \dsq{Who created and funded it?} \benchname{} was created by the authors of this paper as part of an academic research effort on multi-source urban data integration; author names and affiliations are withheld for anonymous review and disclosed at camera-ready. Funding is detailed in the Acknowledgements.

\subsection{Composition}

\dsq{What do the instances represent?} Nodes are buildings and other city objects (building parts, boundary surfaces, geometry primitives, and, per city, bridges, roads, vegetation, water bodies, and city furniture), together with fused OpenStreetMap features and, for Hamburg, ML-predicted roof materials and reconstructed LoD3 facades.

\dsq{How many instances are there?} Five city graphs totaling roughly \num{180}~GiB, \num{180}M nodes, \num{220}M edges, \num{1.2}B properties, and \num{3.6}M buildings, ranging from \num{2980} buildings (Helsinki) to \num{2005762} (Tokyo).

\dsq{Are relationships between instances made explicit?} Yes, relationships are the dataset. Typed, directed edges connect each building to its parts, boundary surfaces, and geometry primitives, and to its district and city, while cross-source fusion is expressed as explicit \texttt{ENRICHED\_BY}/\texttt{HAS\_POI} edges carrying overlap and confidence attributes rather than as silently merged scalars (\Cref{sec:dataset}). Every edge records its source, so provenance is queryable at the relationship level, not just the node level.

\dsq{Is any information missing, and does the dataset contain all instances or a sample?} It contains all buildings of each released source extent; coverage of derived layers is deliberately partial and flagged (ML roof materials cover \num{50.2}\,\% of Hamburg buildings; LoD3 covers 17 Hamburg buildings), so absence is never a negative label.

\dsq{Does the dataset contain confidential or personal data?} No. Instances are buildings, not people. All sources are already public. Address-level strings are retained only where OpenStreetMap itself publishes them. The remaining datasheet questions concerning human data subjects (consent, ethical review, data retention, offensive content) are therefore not applicable.

\dsq{Are there errors, noise, or redundancies?} Yes, and they are documented and preserved rather than silently altered. Some source datasets ship non-unique \texttt{gml:id}s, resolved at ingest by content and geometry hashing. 
%(\Cref{app:dupids}). 
Tokyo's PLATEAU uses a $\pm$\texttt{9999} height sentinel, kept verbatim but excluded from all statistics. The Helsinki source export carries upstream-corrupted Finnish diacritics (\Cref{sec:limitations}).

\subsection{Collection Process}

\dsq{How was the data acquired?} Authoritative CityGML was downloaded from the open-government portals in \Cref{tab:city_sources}. OpenStreetMap was obtained as regional Geofabrik extracts. The roof-material predictions are produced by our own imagery-based model~\cite{lukas_roofmaterials}, and the LoD3 facades are reconstructed from our own facade imagery.

\dsq{Over what timeframe?} The source vintages are listed per city in \Cref{tab:city_sources}. The graph was constructed and fused in 2026.

\subsection{Preprocessing, Cleaning, and Labeling}

\dsq{Was any preprocessing done?} CityGML is mapped to a labeled property graph with coordinates stored verbatim. Non-metric sources (New York, Tokyo) are reprojected to a per-city metric CRS at ingest (\Cref{tab:city_sources}). OpenStreetMap footprints are matched to authoritative footprints by a confidence-weighted overlap graph and attached as \texttt{ENRICHED\_BY} edges without overwriting authoritative values. Derived scalars are prefixed (\texttt{osm\_*}, \texttt{predicted\_*}).

\dsq{Is the raw source available?} Yes, through the sources given in \Cref{tab:city_sources}. The original coordinates, CRS, and \texttt{gml:id}s are retained as provenance so the mapping is auditable.

\dsq{Was the data validated?} Every released instance passes automated integrity gates: a lossless round-trip census on the authoritative layer and an OpenStreetMap ingestion-completeness census on the fused layer (\Cref{subsec:osm_fusion}).

\subsection{Uses}

\dsq{Has the dataset been used for any tasks already?} Yes, the two benchmark families reported in this paper, natural-language-to-query translation (Task~A) and graph representation learning (Task~B), are evaluated on it, and their baseline results are released alongside the dataset (\Cref{subsec:query_tasks,subsec:rl_tasks}). We are aware of no third-party uses at time of release.

\dsq{What tasks is the dataset intended for?} The two benchmark families (\Cref{subsec:query_tasks,subsec:rl_tasks}): natural-language-to-query translation and graph representation learning, both emphasizing provenance-, coverage-, and cross-source-aware reasoning. Additionally, the dataset can serve as a foundation for future graph-based urban analyses and data fusion.

\dsq{What uses should be avoided?} The dataset describes the built environment, not individuals, and should not be repurposed to infer information about residents. Derived layers are marked as predicted or reconstructed precisely so consumers can exclude them where authoritative-only evidence is required.

\subsection{Distribution}

\dsq{How is the dataset distributed and under what license?} As a Neo4j dump (27.2 GiB) plus a backend-neutral node/edge-table export, with loaders, task splits, gold queries and materialized answers, the evaluation harness, and this datasheet (\Cref{sec:dataset}). Each source retains its own license (Hamburg dl-de/by-2-0~\cite{dlde_by20_2026}, Zurich swisstopo~\cite{swisstopo_ogdterms_2022}, Helsinki CC~BY~4.0~\cite{creativecommons_ccby40_2026}, New York NYC OpenData~\cite{nyc_opendata_2026}, Tokyo PLATEAU Public Data License 1.0~\cite{tokyo_license_2024}). OpenStreetMap is licensed under the Open Data Commons Open Database License (ODbL)~\cite{opendatacommons_odbl_2026} by the OpenStreetMap Foundation (OSMF). Our annotations, question suite, and documentation are released under CC~BY~4.0.

\dsq{Is there a DOI?} Yes, the artifact is archived on Zenodo under DOI \href{https://doi.org/10.5281/zenodo.21547211}{\texttt{10.5281/zenodo.21547211}}.

\subsection{Maintenance}

\dsq{Who maintains it and how is it versioned?} The authors host the leaderboard and version the suite. Any change to a gold answer on a dataset rebuild bumps the suite version under a changelog (gate G5, \Cref{app:questions_gate}).

\dsq{Will it be extended?} Yes: further Tier~1 cities and a sensed layer are planned as inference data becomes available (\Cref{sec:limitations}).

\dsq{Contact.} The corresponding author of this paper, Huynh Duc An Son Nguyen (\href{mailto:son.nguyen@hcu-hamburg.de}{son.nguyen@hcu-hamburg.de}).

\section{Limitations}\label{sec:limitations}

\dsq{Geographic scope.} Tier~2 exists for one city and Tier~1 spans five across three continents. Findings may not transfer to regions with different cadastral traditions or OSM community density. 
%We could not include an Australian city because no open CityGML dataset is published there. 

\dsq{Cross-city heterogeneity.} National schemes differ in detail and vocabulary, so cross-city questions use a common attribute subset (\Cref{sec:dataset}). Per-city attributes remain queryable but outside the comparable core. The binding constraint for the ML-predicted layer is inference-imagery availability. 

\dsq{Upstream source defects.} Source values are preserved verbatim rather than silently repaired: the Helsinki export carries upstream-corrupted Finnish diacritics (\num{1490} Unicode replacement characters across \num{283} values), which we retain as-is while the integrity gates confirm faithful preservation. 

\dsq{Prediction bias.} The roof-material predictions inherit their training-imagery biases~\cite{lukas_roofmaterials} and are marked as predicted. 

\section{CityGML-OSM Knowledge Graph}\label{app:enriched_osm}

\Cref{fig:osm_enrich_example} illustrates the structure and content of \benchname{}'s LPGs, which extend authoritative CityGML graphs with complementary OSM information. The enrichment process is strictly additive: authoritative CityGML entities and attributes are preserved, while OSM-derived information is attached either as source-prefixed properties on CityGML feature nodes or as separate OSM feature nodes linked through explicit relationships. Building-level correspondences between the two sources are represented through enrichment edges that record spatial matching statistics, including overlap ratio and Jaccard similarity, and identify the primary match based on the largest overlap. To avoid introducing ambiguity, only attributes from the primary OSM match are propagated to the corresponding CityGML building, while information from secondary matches remains associated with its original OSM feature. 
%Additional contextual information, such as points of interest and building footprints, is connected through dedicated relationships that encode relevant spatial measures. 
Throughout the graph, provenance is retained at the source level, enabling users to distinguish authoritative CityGML information from OSM-derived enrichments and trace the origin of all integrated data. 

We choose \(\tau\) (\Cref{subsec:osm_fusion}) through a sensitivity analysis of the correspondence structure. Since no city-scale ground-truth CityGML--OSM correspondences exist, we sweep \(\tau \in [0.1, 0.7]\) and evaluate the mean Jaccard of primary matches (precision proxy) and the number of enriched buildings (recall proxy). The mean primary-match Jaccard remains constant (0.775), indicating that \(\tau\) affects only marginal edges. 
% Higher thresholds remove valid partial correspondences, whereas lower thresholds admit low-Jaccard slivers that spuriously connect components and increase \(n{:}m\) correspondences. 
We therefore select \(\tau = 0.3\), which provides near-maximal enrichment while preserving correspondence quality. Because each retained edge stores its Jaccard score, downstream methods can reweight or rethreshold matches without repeating the alignment. 

\begin{figure*}[h!]
  \centering
  % Neo4j Browser palette
  \definecolor{osmFill}{HTML}{57C7E3}\definecolor{osmEdge}{HTML}{23B3D7}%
  \definecolor{cgFill}{HTML}{ECB5C9}\definecolor{cgEdge}{HTML}{DA7298}%
  \definecolor{geoFill}{HTML}{A5ABB6}\definecolor{geoEdge}{HTML}{9AA1AC}%
  \definecolor{poiFill}{HTML}{8DCC93}\definecolor{poiEdge}{HTML}{5DB665}%
  \begin{tikzpicture}[
    gn/.style={circle, draw=geoEdge, fill=geoFill, line width=1pt, minimum size=1.3cm, inner sep=1pt, font=\scriptsize, text=white},
    on/.style={gn, draw=osmEdge, fill=osmFill, text=black!80},
    cn/.style={gn, draw=cgEdge, fill=cgFill, text=black!80},
    pn/.style={gn, draw=poiEdge, fill=poiFill, text=black!80},
    rel/.style={-{Latex[length=5pt,width=4.5pt]}, draw=black!35, line width=0.8pt},
    enr/.style={-{Latex[length=5.5pt,width=5pt]}, draw=black!70, line width=1pt},
    rlbl/.style={font=\scriptsize\ttfamily, text=black!75, inner sep=1.5pt},
    plbl/.style={font=\scriptsize\itshape, text=black!55, inner sep=1.5pt},
    albl/.style={font=\scriptsize\ttfamily, inner sep=1pt}
  ]
    % ---- left: OSM geometry chain (exterior ring <- polygon <- footprint) ----
    \node[gn] (ext1)  at (0,3.8)   {exterior};
    \node[gn] (ext2)  at (0,0)     {exterior};
    \node[gn] (poly1) at (3.2,3.8) {};
    \node[gn] (poly2) at (3.2,0)   {};
    % ---- OSM building nodes (blue) ----
    \node[on] (osm1) at (6.4,3.8) {commercial};
    \node[on] (osm2) at (6.4,0)   {residential};
    \node[albl, text=osmEdge!80!black, above=1.5mm of osm1] {building = commercial};
    \node[albl, text=osmEdge!80!black, below=1.5mm of osm2] {building = residential};
    % ---- CityGML building nodes (pink) ----
    \node[cn] (cg1) at (11.4,3.8) {commercial};
    \node[cn] (cg2) at (11.4,0)   {residential};
    \node[albl, text=cgEdge!75!black, above=1.5mm of cg1] {osm\_building = commercial};
    \node[albl, text=cgEdge!75!black, below=1.5mm of cg2] {osm\_building = residential};
    % ---- left R-tree node serving both OSM features ----
    \node[gn] (rtL) at (3.6,1.9) {};
    % ---- right: CityGML geometry / index fan ----
    \node[gn] (rt1)    at (15.0,5.0)  {};
    \node[gn] (surf1)  at (16.1,3.8)  {};
    \node[gn] (solid1) at (15.0,2.6)  {};
    \node[gn] (rt2)    at (15.0,-1.2) {};
    \node[gn] (surf2)  at (16.1,0)    {};
    \node[gn] (solid2) at (15.0,1.2)  {};
    % ---- structural edges ----
    \draw[rel] (poly1) -- node[rlbl, above] {HAS\_EXTERIOR\_RING} (ext1);
    \draw[rel] (poly2) -- node[rlbl, above] {HAS\_EXTERIOR\_RING} (ext2);
    \draw[rel] (osm1)  -- node[rlbl, above] {HAS\_FOOTPRINT} (poly1);
    \draw[rel] (osm2)  -- node[rlbl, above] {HAS\_FOOTPRINT} (poly2);
    \draw[rel] (rtL)   -- node[rlbl, above, sloped] {RTREE\_REFERENCE} (osm1);
    \draw[rel] (rtL)   -- node[rlbl, below, sloped] {RTREE\_REFERENCE} (osm2);
    \draw[rel] (rt1)   -- node[rlbl, above, sloped] {RTREE\_REFERENCE} (cg1);
    \draw[rel] (cg1)   -- node[rlbl, above] {HAS\_LOD\_SURFACE} (surf1);
    \draw[rel] (cg1)   -- node[rlbl, below, sloped] {HAS\_LOD\_SOLID} (solid1);
    \draw[rel] (rt2)   -- node[rlbl, below, sloped] {RTREE\_REFERENCE} (cg2);
    \draw[rel] (cg2)   -- node[rlbl, below] {HAS\_LOD\_SURFACE} (surf2);
    \draw[rel] (cg2)   -- node[rlbl, above, sloped] {HAS\_LOD\_SOLID} (solid2);
    % ---- enrichment edges (primary solid, non-primary dashed) ----
    \draw[enr] (cg1) -- node[rlbl, above] {ENRICHED\_BY}
                        node[plbl, below] {is\_primary = true} (osm1);
    \draw[enr] (cg2) -- node[rlbl, above] {ENRICHED\_BY}
                        node[plbl, below] {is\_primary = true} (osm2);
    \draw[enr] (cg2) -- node[rlbl, above, sloped] {ENRICHED\_BY}
                                node[plbl, below, sloped] {is\_primary = false} (osm1);
    % ---- OSM point-of-interest attached to its containing building ----
    \node[pn] (poi1) at (11.7,1.8) {restaurant};
    \node[albl, text=poiEdge!60!black, below=1.5mm of poi1] {amenity = restaurant};
    \draw[enr] (cg1) -- node[rlbl, right=1pt, pos=0.35] {HAS\_POI}
                        node[plbl, right=1pt, pos=0.72] {distance\_m = 4.2} (poi1);
  \end{tikzpicture}
  \caption{
    OSM enrichment: CityGML buildings (pink) are linked to OSM buildings (blue) by \texttt{ENRICHED\_BY} edges. Buildings inherit source-prefixed \texttt{osm\_*} attributes from their primary match, while points of interest (green) attach via \texttt{HAS\_POI}.
    %OSM enrichment: CityGML building nodes (pink) fused to OSM building nodes (blue) by \texttt{ENRICHED\_BY} edges, each recording its overlap ratio and Jaccard with the \emph{primary} (largest-overlap) match flagged. A building copies the source-prefixed \texttt{osm\_*} values of its primary match only; values from non-primary matches stay on their own OSM node. The lower CityGML building overlaps both OSM buildings but takes \texttt{building}=\emph{residential} from its primary match, while the OSM \emph{commercial} node is primary for the upper building. A point of interest (green) attaches by a \texttt{HAS\_POI} edge from the building whose footprint contains it, recording its distance to the centroid. Enrichment is additive and never overwrites an authoritative value; gray nodes are the attached geometry (CityGML solids and surfaces right, OSM footprints left).
  }
  \Description{Graph diagram with two pink CityGML building nodes on the right connected by ENRICHED_BY edges to two blue OSM building nodes in the center; the lower pink node has an additional non-primary edge to the upper blue node; a green OSM point-of-interest node labeled restaurant hangs below the upper pink node, reached by a HAS_POI edge carrying a distance_m attribute; gray nodes for geometry and spatial index surround both sides.}
  \label{fig:osm_enrich_example}
\end{figure*}

\section{Per-City Dataset and Cross-Source Analysis}\label{app:dataset_stats}

This section gives
an overview comparison of all five datasets~\Cref{fig:city_stats}, 
the full per-city statistics summarized in \Cref{sec:dataset}, 
dataset provenance and scale (\Cref{tab:city_sources}), 
graph size and content (\Cref{tab:city_graph_content}), 
%storage and build cost (\Cref{tab:city_storage}), 
%the OSM ingestion completeness census (\Cref{tab:osm_census}), 
footprint correspondence between CityGML and OSM (\Cref{tab:footprint_cases}), 
and dual-sourced attribute agreement (\Cref{tab:dual_attrs}). 
% \Cref{fig:lod3_facades} shows the reconstructed LoD3 layer.

\begingroup
\definecolor{cham}{HTML}{2A78D6}\definecolor{chel}{HTML}{008300}%
\definecolor{czur}{HTML}{E87BA4}\definecolor{cnyc}{HTML}{EDA100}\definecolor{ctok}{HTML}{1BAF7A}%
\colorlet{cgray}{black!7}%
\tikzset{cbar/.style={rounded corners=0.5pt},
  rl/.style={left,font=\normalsize,inner sep=2.5pt},
  vl/.style={right,font=\small,text=black!55,inner sep=2.5pt}}%
\begin{figure}[h!]
  \centering
  % Each of the four quadrant charts is forced to an identical bounding box
  % (-2.1..6.5 wide) inside a 0.5\textwidth box, so all four have the SAME width,
  % exactly half of the full-width composition chart below.
  % ---- row 1: detail density | schema richness ----
  \begin{minipage}[t]{0.5\textwidth}\centering
    {\normalsize\bfseries Detail density: nodes per building}\par\vspace{2pt}
    \begin{tikzpicture}\def\MX{369.8}\def\AW{5.0}\def\bh{0.26}
      \useasboundingbox(-2.1,0.15)rectangle(6.5,-3.35);
      \foreach \lbl/\md/\mn/\c [count=\i] in {Hamburg/30/41/cham,Helsinki/54/107.4/chel,Zurich/253/369.8/czur,New York/31/36.9/cnyc,Tokyo/21/28.9/ctok}{
        \pgfmathsetmacro\yb{-\i*0.62}\pgfmathsetmacro\wm{\md/\MX*\AW}\pgfmathsetmacro\wa{\mn/\MX*\AW}
        \fill[cgray,cbar](0,\yb+0.32)rectangle(\AW,\yb+0.32+\bh);\fill[\c,cbar](0,\yb+0.32)rectangle(\wm,\yb+0.32+\bh);\node[vl]at(\wm,\yb+0.32+\bh/2){\md};
        \fill[cgray,cbar](0,\yb)rectangle(\AW,\yb+\bh);\fill[\c!42!white,cbar](0,\yb)rectangle(\wa,\yb+\bh);\node[vl]at(\wa,\yb+\bh/2){\mn};
        \node[rl]at(0,\yb+0.29){\lbl};}
    \end{tikzpicture}\par\vspace{-2pt}{\small\color{black!55}Darker = median \qquad lighter = mean}\par\vspace{5pt}
  \end{minipage}%
  \hfill 
  \begin{minipage}[t]{0.5\textwidth}\centering
    {\normalsize\bfseries Storage by source layer}\par\vspace{2pt}
    \begin{tikzpicture}\def\MX{99.18}\def\AW{5.0}\def\bh{0.38}
      \useasboundingbox(-2.1,0.15)rectangle(6.5,-3.35);
      \foreach \lbl/\cg/\tot/\tt/\c [count=\i] in {Hamburg/12.71/14.77/15/cham,Helsinki/0.22/0.28/0.28/chel,Zurich/27.06/30.21/30/czur,New York/27.51/32.35/32/cnyc,Tokyo/91.92/99.18/99/ctok}{
        \pgfmathsetmacro\yy{-\i*0.62}\pgfmathsetmacro\wcg{\cg/\MX*\AW}\pgfmathsetmacro\wt{\tot/\MX*\AW}
        \fill[cgray,cbar](0,\yy)rectangle(\AW,\yy+\bh);\fill[\c,cbar](0,\yy)rectangle(\wcg,\yy+\bh);\fill[\c!38!white,cbar](\wcg,\yy)rectangle(\wt,\yy+\bh);
        \node[rl]at(0,\yy+\bh/2){\lbl};\node[vl]at(\wt,\yy+\bh/2){\tt\,GiB};}
    \end{tikzpicture}\par\vspace{-2pt}{\small\color{black!55}Darker = CityGML \qquad lighter = $+$ OSM}\par\vspace{5pt}
  \end{minipage}
  \hfill 
  %\par\vspace{4pt}
  % ---- row 2: storage | enrichment ----
  \begin{minipage}[t]{0.5\textwidth}\centering
    {\normalsize\bfseries Source schema richness: thematic keys}\par\vspace{2pt}
    \begin{tikzpicture}\def\MX{49}\def\AW{5.0}\def\bh{0.38}
      \useasboundingbox(-2.1,0.15)rectangle(6.5,-3.35);
      \foreach \lbl/\v/\c [count=\i] in {Hamburg/19/cham,Helsinki/49/chel,Zurich/22/czur,New York/5/cnyc,Tokyo/20/ctok}{
        \pgfmathsetmacro\yy{-\i*0.62}\pgfmathsetmacro\bw{\v/\MX*\AW}
        \fill[cgray,cbar](0,\yy)rectangle(\AW,\yy+\bh);\fill[\c,cbar](0,\yy)rectangle(\bw,\yy+\bh);
        \node[rl]at(0,\yy+\bh/2){\lbl};\node[vl]at(\bw,\yy+\bh/2){\v};}
    \end{tikzpicture}%\par\vspace{0pt}{\small\strut}
  \end{minipage}%
  \hfill 
  \begin{minipage}[t]{0.5\textwidth}\centering
    {\normalsize\bfseries OSM enrichment rate}\par\vspace{2pt}
    \begin{tikzpicture}\def\MX{100}\def\AW{5.0}\def\bh{0.38}
      \useasboundingbox(-2.1,0.15)rectangle(6.5,-3.35);
      \draw[dashed,black!30](\AW,0.02)--(\AW,-3.15);\node[font=\small,text=black!45,below]at(\AW,-3.15){100\%};
      %\foreach \lbl/\v/\c [count=\i] in {Hamburg/85.3/cham,Helsinki/83.1/chel,Zurich/91.5/czur,New York/98.8/cnyc,Tokyo/58.9/ctok}{
      \foreach \lbl/\v/\c [count=\i] in {Hamburg/84.9/cham,Helsinki/81.1/chel,Zurich/91.2/czur,New York/98.7/cnyc,Tokyo/57.8/ctok}{
        \pgfmathsetmacro\yy{-\i*0.62}\pgfmathsetmacro\bw{\v/\MX*\AW}
        \fill[cgray,cbar](0,\yy)rectangle(\AW,\yy+\bh);\fill[\c,cbar](0,\yy)rectangle(\bw,\yy+\bh);
        \node[rl]at(0,\yy+\bh/2){\lbl};\node[vl]at(\bw,\yy+\bh/2){\v\%};}
    \end{tikzpicture}%\par\vspace{1pt}{\small\strut}
  \end{minipage}
  % \par\vspace{6pt}
  % ---- full-width: node composition (monochrome per city); width = 2x a quadrant ----
  % {\normalsize\bfseries Node composition: every graph is geometry-dominated}\par\vspace{2pt}
  % \begin{tikzpicture}\def\AW{15.0}\def\bh{0.40}
  %   \useasboundingbox(-2.1,0.78)rectangle(15.1,-3.32);
  %   \foreach \nm/\mix/\lx in {Geometry/100/0,Semantic/64/2.9,{Crowd-sourced OSM}/42/5.9,{Index / containers}/22/10.6}{
  %     \fill[cham!\mix!white,cbar](\lx,0.30)rectangle(\lx+0.32,0.62);\node[right,font=\small,text=black!55]at(\lx+0.38,0.46){\nm};}
  %   \foreach \lbl/\sem/\geo/\osm/\idx/\gt/\c [count=\i] in {Hamburg/2.2/90.8/6.8/0.1/91/cham,Helsinki/0.7/85.6/13.5/0.2/86/chel,Zurich/0.3/95.2/4.5/0.1/95/czur,New York/2.4/91.9/5.6/0.1/92/cnyc,Tokyo/3.6/91.1/5.1/0.1/91/ctok}{
  %     \pgfmathsetmacro\yy{-\i*0.60}\pgfmathsetmacro\tot{\sem+\geo+\osm+\idx}
  %     \pgfmathsetmacro\wg{\geo/\tot*\AW}\pgfmathsetmacro\ws{\sem/\tot*\AW}\pgfmathsetmacro\wo{\osm/\tot*\AW}
  %     \pgfmathsetmacro\xb{\wg+\ws}\pgfmathsetmacro\xc{\wg+\ws+\wo}
  %     \fill[\c,cbar](0,\yy)rectangle(\wg,\yy+\bh);\fill[\c!64!white,cbar](\wg,\yy)rectangle(\xb,\yy+\bh);\fill[\c!42!white,cbar](\xb,\yy)rectangle(\xc,\yy+\bh);\fill[\c!22!white,cbar](\xc,\yy)rectangle(\AW,\yy+\bh);
  %     \node[left,font=\normalsize,inner sep=3pt]at(0,\yy+\bh/2){\lbl};\node[font=\small,text=white]at(\wg/2,\yy+\bh/2){\gt\%};}
  % \end{tikzpicture}
  \caption{
    Comparison of the five \benchname{} graphs. Zurich is the densest, Tokyo the largest, Helsinki the most attribute-rich, and New York the most OSM-enriched. Geometry accounts for 85--95\,\% of graph content in every city. 
    % Full statistics in \Cref{app:dataset_stats}.
    % Per-city comparison of the five \benchname{} graphs (fixed color per city). \emph{Detail density}: median (solid) and mean (faded) nodes per building; level of detail, not city size, drives graph scale (a Zurich building costs $\sim$12$\times$ a Tokyo one). \emph{Schema richness}: distinct source thematic keys (Helsinki~49 vs New York~5), the multilingual axis. \emph{Storage}: the authoritative CityGML layer (darker) dominates, OSM fusion (lighter) adds a modest slice. \emph{OSM enrichment rate}: \num{58.9}--\num{98.8}\,\%. \emph{Node composition}: every graph is 85--95\,\% geometry. Full per-city figures in \Cref{app:dataset_stats}.
    }
  \Description{Five side-by-side bar-chart panels comparing Hamburg, Helsinki, Zurich, New York, and Tokyo across graph characteristics: detail density (median and mean nodes per building), schema richness (distinct thematic keys), storage split between CityGML and OSM, OSM enrichment rate percentages, and node composition dominated by geometry.}
  \label{fig:city_stats}
\end{figure}
\endgroup

% \begingroup
% \definecolor{cham}{HTML}{2A78D6}\definecolor{chel}{HTML}{008300}%
% \definecolor{czur}{HTML}{E87BA4}\definecolor{cnyc}{HTML}{EDA100}\definecolor{ctok}{HTML}{1BAF7A}%
% \colorlet{cgray}{black!7}%
% \tikzset{cbar/.style={rounded corners=0.5pt},
%   rl/.style={left,font=\normalsize,inner sep=2.5pt},
%   vl/.style={right,font=\small,text=black!55,inner sep=2.5pt}}%
% \begin{figure*}[ht!]
%   \centering
%   % Each of the four quadrant charts is forced to an identical bounding box
%   % (-2.1..6.5 wide) inside a 0.5\textwidth box, so all four have the SAME width,
%   % exactly half of the full-width composition chart below.
%   % ---- row 1: detail density | schema richness ----
%   \begin{minipage}[t]{0.5\textwidth}\centering
%     {\normalsize\bfseries Detail density: nodes per building}\par\vspace{2pt}
%     \begin{tikzpicture}\def\MX{369.8}\def\AW{5.0}\def\bh{0.26}
%       \useasboundingbox(-2.1,0.15)rectangle(6.5,-3.35);
%       \foreach \lbl/\md/\mn/\c [count=\i] in {Hamburg/30/41/cham,Helsinki/54/107.4/chel,Zurich/253/369.8/czur,New York/31/36.9/cnyc,Tokyo/21/28.9/ctok}{
%         \pgfmathsetmacro\yb{-\i*0.62}\pgfmathsetmacro\wm{\md/\MX*\AW}\pgfmathsetmacro\wa{\mn/\MX*\AW}
%         \fill[cgray,cbar](0,\yb+0.32)rectangle(\AW,\yb+0.32+\bh);\fill[\c,cbar](0,\yb+0.32)rectangle(\wm,\yb+0.32+\bh);\node[vl]at(\wm,\yb+0.32+\bh/2){\md};
%         \fill[cgray,cbar](0,\yb)rectangle(\AW,\yb+\bh);\fill[\c!42!white,cbar](0,\yb)rectangle(\wa,\yb+\bh);\node[vl]at(\wa,\yb+\bh/2){\mn};
%         \node[rl]at(0,\yb+0.29){\lbl};}
%     \end{tikzpicture}\par\vspace{-2pt}{\small\color{black!55}Darker = median \qquad lighter = mean}
%   \end{minipage}%
%   \begin{minipage}[t]{0.5\textwidth}\centering
%     {\normalsize\bfseries Storage by source layer}\par\vspace{2pt}
%     \begin{tikzpicture}\def\MX{99.18}\def\AW{5.0}\def\bh{0.38}
%       \useasboundingbox(-2.1,0.15)rectangle(6.5,-3.35);
%       \foreach \lbl/\cg/\tot/\tt/\c [count=\i] in {Hamburg/12.71/14.77/15/cham,Helsinki/0.22/0.28/0.28/chel,Zurich/27.06/30.21/30/czur,New York/27.51/32.35/32/cnyc,Tokyo/91.92/99.18/99/ctok}{
%         \pgfmathsetmacro\yy{-\i*0.62}\pgfmathsetmacro\wcg{\cg/\MX*\AW}\pgfmathsetmacro\wt{\tot/\MX*\AW}
%         \fill[cgray,cbar](0,\yy)rectangle(\AW,\yy+\bh);\fill[\c,cbar](0,\yy)rectangle(\wcg,\yy+\bh);\fill[\c!38!white,cbar](\wcg,\yy)rectangle(\wt,\yy+\bh);
%         \node[rl]at(0,\yy+\bh/2){\lbl};\node[vl]at(\wt,\yy+\bh/2){\tt\,GiB};}
%     \end{tikzpicture}\par\vspace{-2pt}{\small\color{black!55}Darker = CityGML \qquad lighter = $+$ OSM}
%   \end{minipage}
%   \par\vspace{4pt}
%   % ---- row 2: storage | enrichment ----
%   \begin{minipage}[t]{0.5\textwidth}\centering
%     {\normalsize\bfseries Source schema richness: thematic keys}\par\vspace{2pt}
%     \begin{tikzpicture}\def\MX{49}\def\AW{5.0}\def\bh{0.38}
%       \useasboundingbox(-2.1,0.15)rectangle(6.5,-3.35);
%       \foreach \lbl/\v/\c [count=\i] in {Hamburg/19/cham,Helsinki/49/chel,Zurich/22/czur,New York/5/cnyc,Tokyo/20/ctok}{
%         \pgfmathsetmacro\yy{-\i*0.62}\pgfmathsetmacro\bw{\v/\MX*\AW}
%         \fill[cgray,cbar](0,\yy)rectangle(\AW,\yy+\bh);\fill[\c,cbar](0,\yy)rectangle(\bw,\yy+\bh);
%         \node[rl]at(0,\yy+\bh/2){\lbl};\node[vl]at(\bw,\yy+\bh/2){\v};}
%     \end{tikzpicture}%\par\vspace{0pt}{\small\strut}
%   \end{minipage}%
%   \begin{minipage}[t]{0.5\textwidth}\centering
%     {\normalsize\bfseries OSM enrichment rate}\par\vspace{2pt}
%     \begin{tikzpicture}\def\MX{100}\def\AW{5.0}\def\bh{0.38}
%       \useasboundingbox(-2.1,0.15)rectangle(6.5,-3.35);
%       \draw[dashed,black!30](\AW,0.02)--(\AW,-3.15);\node[font=\small,text=black!45,below]at(\AW,-3.15){100\%};
%       \foreach \lbl/\v/\c [count=\i] in {Hamburg/85.3/cham,Helsinki/83.1/chel,Zurich/91.5/czur,New York/98.8/cnyc,Tokyo/58.9/ctok}{
%         \pgfmathsetmacro\yy{-\i*0.62}\pgfmathsetmacro\bw{\v/\MX*\AW}
%         \fill[cgray,cbar](0,\yy)rectangle(\AW,\yy+\bh);\fill[\c,cbar](0,\yy)rectangle(\bw,\yy+\bh);
%         \node[rl]at(0,\yy+\bh/2){\lbl};\node[vl]at(\bw,\yy+\bh/2){\v\%};}
%     \end{tikzpicture}%\par\vspace{1pt}{\small\strut}
%   \end{minipage}
%   % \par\vspace{6pt}
%   % ---- full-width: node composition (monochrome per city); width = 2x a quadrant ----
%   % {\normalsize\bfseries Node composition: every graph is geometry-dominated}\par\vspace{2pt}
%   % \begin{tikzpicture}\def\AW{15.0}\def\bh{0.40}
%   %   \useasboundingbox(-2.1,0.78)rectangle(15.1,-3.32);
%   %   \foreach \nm/\mix/\lx in {Geometry/100/0,Semantic/64/2.9,{Crowd-sourced OSM}/42/5.9,{Index / containers}/22/10.6}{
%   %     \fill[cham!\mix!white,cbar](\lx,0.30)rectangle(\lx+0.32,0.62);\node[right,font=\small,text=black!55]at(\lx+0.38,0.46){\nm};}
%   %   \foreach \lbl/\sem/\geo/\osm/\idx/\gt/\c [count=\i] in {Hamburg/2.2/90.8/6.8/0.1/91/cham,Helsinki/0.7/85.6/13.5/0.2/86/chel,Zurich/0.3/95.2/4.5/0.1/95/czur,New York/2.4/91.9/5.6/0.1/92/cnyc,Tokyo/3.6/91.1/5.1/0.1/91/ctok}{
%   %     \pgfmathsetmacro\yy{-\i*0.60}\pgfmathsetmacro\tot{\sem+\geo+\osm+\idx}
%   %     \pgfmathsetmacro\wg{\geo/\tot*\AW}\pgfmathsetmacro\ws{\sem/\tot*\AW}\pgfmathsetmacro\wo{\osm/\tot*\AW}
%   %     \pgfmathsetmacro\xb{\wg+\ws}\pgfmathsetmacro\xc{\wg+\ws+\wo}
%   %     \fill[\c,cbar](0,\yy)rectangle(\wg,\yy+\bh);\fill[\c!64!white,cbar](\wg,\yy)rectangle(\xb,\yy+\bh);\fill[\c!42!white,cbar](\xb,\yy)rectangle(\xc,\yy+\bh);\fill[\c!22!white,cbar](\xc,\yy)rectangle(\AW,\yy+\bh);
%   %     \node[left,font=\normalsize,inner sep=3pt]at(0,\yy+\bh/2){\lbl};\node[font=\small,text=white]at(\wg/2,\yy+\bh/2){\gt\%};}
%   % \end{tikzpicture}
%   \caption{
%     Comparison of the five \benchname{} graphs. Zurich is the densest, Tokyo the largest, Helsinki the most attribute-rich, and New York the most OSM-enriched. Geometry accounts for 85--95\,\% of graph content in every city. Full statistics in \Cref{app:dataset_stats}.
%     % Per-city comparison of the five \benchname{} graphs (fixed color per city). \emph{Detail density}: median (solid) and mean (faded) nodes per building; level of detail, not city size, drives graph scale (a Zurich building costs $\sim$12$\times$ a Tokyo one). \emph{Schema richness}: distinct source thematic keys (Helsinki~49 vs New York~5), the multilingual axis. \emph{Storage}: the authoritative CityGML layer (darker) dominates, OSM fusion (lighter) adds a modest slice. \emph{OSM enrichment rate}: \num{58.9}--\num{98.8}\,\%. \emph{Node composition}: every graph is 85--95\,\% geometry. Full per-city figures in \Cref{app:dataset_stats}.
%     }
%   \Description{Five side-by-side bar-chart panels comparing Hamburg, Helsinki, Zurich, New York, and Tokyo across graph characteristics: detail density (median and mean nodes per building), schema richness (distinct thematic keys), storage split between CityGML and OSM, OSM enrichment rate percentages, and node composition dominated by geometry.}
%   \label{fig:city_stats}
% \end{figure*}
% \endgroup

\input{tables/city_sources_table.tex}

% \input{tables/city_enrichment_table.tex}

\input{tables/city_graph_content_table.tex}

% \input{tables/city_storage_table.tex}

\begin{table}[h!]
  \caption{
  Footprint correspondence between CityGML and OSM buildings.
  % for the four fused city instances: connected components of the overlap graph by case, with the buildings involved.
  }
  \label{tab:footprint_cases}
  %\small
  \begin{tabular}{lrrrr}
    \toprule
    Case & Components & Share & CityGML bldgs & OSM bldgs \\
    \midrule
    \multicolumn{5}{l}{\emph{Hamburg}} \\
    1:1 & \num{256219} & 87.2\,\% & \num{255984} & \num{255984} \\
    1:n & \num{19167} & 6.5\,\% & \num{47936} & \num{19096} \\
    n:1 & \num{16146} & 5.5\,\% & \num{16142} & \num{47838} \\
    n:m & \num{2363} & 0.8\,\% & \num{9537} & \num{9151} \\
    \midrule
    \multicolumn{5}{l}{\emph{Helsinki}} \\
    1:1 & \num{1692} & 84.2\,\% & \num{1686} & \num{1686} \\
    1:n & \num{127} & 6.3\,\% & \num{285} & \num{125} \\
    n:1 & \num{129} & 6.4\,\% & \num{127} & \num{281} \\
    n:m & \num{62} & 3.1\,\% & \num{318} & \num{302} \\
    \midrule
    \multicolumn{5}{l}{\emph{Zurich}} \\
    1:1 & \num{75541} & 90.8\,\% & \num{75430} & \num{75430} \\
    1:n & \num{4858} & 5.8\,\% & \num{12442} & \num{4853} \\
    n:1 & \num{2003} & 2.4\,\% & \num{1999} & \num{4325} \\
    n:m & \num{787} & 0.9\,\% & \num{3812} & \num{3744} \\
    \midrule
    \multicolumn{5}{l}{\emph{New York}} \\
    1:1 & \num{1053850} & 99.3\,\% & \num{1053727} & \num{1053727} \\
    1:n & \num{3046} & 0.3\,\% & \num{8448} & \num{3015} \\
    n:1 & \num{3199} & 0.3\,\% & \num{3178} & \num{6149} \\
    n:m & \num{850} & 0.1\,\% & \num{4277} & \num{4351} \\
    \midrule
    \multicolumn{5}{l}{\emph{Tokyo}} \\
    1:1 & \num{823560} & 85.9\,\% & \num{821655} & \num{821655} \\
    1:n & \num{41392} & 4.3\,\% & \num{90648} & \num{40561} \\
    n:1 & \num{22288} & 2.3\,\% & \num{22273} & \num{46975} \\
    n:m & \num{71562} & 7.5\,\% & \num{224653} & \num{224226} \\
    \bottomrule
  \end{tabular}

  \smallskip\raggedright\footnotesize
  Reading: in 1:n one OSM polygon covers several cadastral buildings; in n:1 several OSM polygons tile one cadastral building. The case mix is itself a per-city property spanning more than an order of magnitude: the fragmented share (n:1 plus n:m) is 0.4\,\% for New York (whose crowd mapping tracks the cadastre almost one-to-one), 3.3\,\% for Zurich, 6.3\,\% for Hamburg, 9.5\,\% for Helsinki, and 9.8\,\% in Tokyo, whose crowd mapping splits more cadastral buildings into several tagged parts.
\end{table}

\begin{table}[h!]
  \caption{
  Dual-sourced building attributes when a CityGML, OSM or ML-predicted value describe the same property. 
  %Both values are retained on every building (\texttt{osm\_*} prefix beside the authoritative property). The dual coverage is a per-city variable, not a constant of the construction: Hamburg has all three, Helsinki exposes height and storey redundancy, Zurich only height (its source encodes neither storeys nor roof material), and New York has none of the three (its LoD1--2 source carries neither surveyed height nor storeys, and no ML layer was produced) so all its cross-source redundancy is geometric (\Cref{tab:footprint_cases}) and locational (POIs). Height agreement itself varies with how each national source defines a building's height.
  }
  \label{tab:dual_attrs}
  %\small %\footnotesize
  \begin{tabularx}{\linewidth}{llrX}
    \toprule
    Property & City & Dual & Agreement or $|\Delta|$\\
    \midrule
    Storeys & Hamburg & \num{144364} & 87.2\,\% exact; 98.6\,\% $\pm$1 \\
    & Helsinki & \num{48} & 70.8\,\% exact; 91.7\,\% $\pm$1 \\
    % & Zurich & \num{0} & - \\
    % & New York & \num{0} & - \\
    & Tokyo & \num{116022} & 85.9\,\% exact; 94.3\,\% $\pm$1 \\
    \addlinespace
    Height & Hamburg & \num{3404} & 85.0\,\% $\leq$ 2\,m; 6.4\,\% $>$ 5\,m \\
    & Helsinki & \num{76} & 61.8\,\% $\leq$ 2\,m; 28.9\,\% $>$ 5\,m \\
    & Zurich & \num{770} & 46.4\,\% $\leq$ 2\,m; 22.6\,\% $>$ 5\,m \\
    % & New York & \num{0} & - \\
    & Tokyo & \num{61373} & 94.4\,\% $\leq$ 2\,m; 2.5\,\% $>$ 5\,m \\
    \addlinespace
    Roof mat. & Hamburg & \num{3641} & 75.2\,\% over 5 classes \\
    % & Helsinki & \num{0} & - \\
    % & Zurich & \num{0} & - \\
    % & New York & \num{0} & - \\
    % & Tokyo & \num{0} & - \\
    \bottomrule
  \end{tabularx}
\end{table}

% \begin{figure}[h!]
%   \centering
%    \includegraphics[width=\linewidth]{figures/Facade_Buildings_New.png}
%   \caption{An excerpt of reconstructed LoD3 building models for the Hamburg HafenCity district~\cite{nguyen_pykci_2026}. In \benchname{}, each window and door becomes an \texttt{Opening} node on corresponding wall surface. The facades attach to their CityGML LoD2 anchors by explicit \texttt{HAS\_LOD3\_FACADE} edges.}
%   \label{fig:lod3_facades}
% \end{figure}

% \begin{figure}[h!]
%   \centering
%   \begin{subfigure}{0.48\linewidth}
%     \includegraphics[width=\linewidth]{figures/undurchsichtig_windows_doors_bbox_hafencity_1.png}
%     \caption{Reconstructed facades}
%   \end{subfigure}\hfill
%   \begin{subfigure}{0.48\linewidth}
%     \includegraphics[width=\linewidth]{figures/durchsichtig_windows_doors_bbox_hafencity_1.png}
%     \caption{Openings exposed}
%   \end{subfigure}
%   \caption{The reconstructed LoD3 layer for the Hamburg HafenCity district. (a) Opaque rendering of the reconstructed facades; (b) the same buildings rendered translucent to expose the per-opening bounding boxes. In \benchname{}, each window and door becomes an \texttt{Opening} node on its wall surface, and each opening cuts an interior ring (a hole) into the wall polygon. The facades attach to their authoritative LoD2 anchors by explicit \texttt{HAS\_LOD3\_FACADE} edges and never overwrite the measured geometry.}
%   \label{fig:lod3_facades}
% \end{figure}

% \section{Non-Unique Source Identifiers: Traceability}\label{app:dupids}

% This appendix documents every source building affected by a non-unique \texttt{gml:id} in the three Tier~1 cities where it occurs, so that the responsible data providers can reproduce the finding and, if they choose, correct their datasets at the source. \benchname{} does not edit the source: it resolves each collision at ingest (\Cref{subsec:fidelity}) and records the disposition, but the underlying identifiers remain the providers' to fix. The complete, machine-readable record (for \emph{every} affected building, with its source file, line range, native-CRS bounding box, centroid, footprint \texttt{gml:posList}, content classification, and resulting node id) ships with the release as \texttt{output/stats/duplicate\_ids\_traceability.\{md,json\}}; \Cref{tab:dupids} summarizes it and the paragraphs below give representative cases.

% \clearpage
% \begin{table}
%   \caption{Non-unique source identifiers per affected city and how identity resolution disposes of them. ``Buildings involved'' counts source building elements that share a \texttt{gml:id} with at least one other; ``recovered'' are distinct features a naive \texttt{MERGE} would have collapsed and lost, now kept as separate nodes; ``deduplicated'' are true copies a naive keep-all would have fabricated, now merged to one node. Coordinates and extents are recorded in each dataset's native CRS.}
%   \label{tab:dupids}
%   \small
%   \begin{tabularx}{\linewidth}{lccccX}
%     \toprule
%     City & Elements & Distinct ids & Involved & Disposition \\
%     \midrule
%     Helsinki & \num{2980} & \num{2919} & 122 & 61 recovered (split) \\
%     Zurich & \num{102673} & \num{102628} & 62 & 40 recovered (split), 5 dedup \\
%     Tokyo & \num{2980839} & \num{2005294} & \num{1782636} & \num{975077} dedup, 468 split \\
%     \bottomrule
%   \end{tabularx}
% \end{table}
%
% WITH CRS
% \begin{table}
%   \caption{Non-unique source identifiers per affected city and how identity resolution disposes of them. ``Buildings involved'' counts source building elements that share a \texttt{gml:id} with at least one other; ``recovered'' are distinct features a naive \texttt{MERGE} would have collapsed and lost, now kept as separate nodes; ``deduplicated'' are true copies a naive keep-all would have fabricated, now merged to one node. Coordinates and extents are recorded in each dataset's native CRS.}
%   \label{tab:dupids}
%   \small
%   \begin{tabularx}{\linewidth}{lccccX}
%     \toprule
%     City & CRS & Elements & Distinct ids & Involved & Disposition \\
%     \midrule
%     Helsinki & 3879 & \num{2980} & \num{2919} & 122 & 61 recovered (split) \\
%     Zurich & 2056 & \num{102673} & \num{102628} & 62 & 40 recovered (split), 5 dedup \\
%     Tokyo & 6697 & \num{2980839} & \num{2005294} & \num{1782636} & \num{975077} dedup, 468 split \\
%     \bottomrule
%   \end{tabularx}
% \end{table}

% \emph{Helsinki} contains \num{61} identifiers each carried by two \emph{distinct} buildings (verified by differing geometry), so all \num{61} are split and recovered. For example \texttt{BID\_0338b8c4-\ldots} appears twice in \texttt{helsinki\_citygml2\_lod2\_2019.gml}, at lines \num{31179}--\num{31822} and \num{58951}--\num{61965}: the two occurrences share a bounding box but not a geometry, so the second is stored under a synthesized id with \texttt{source\_gml\_id} set to the original, and both survive.

% \emph{Zurich} has \num{17} colliding identifiers over \num{62} buildings in \num{78} tiles, of two kinds. The placeholder \texttt{gml:id="ID\_"} is stamped on \num{30} buildings spread over \num{25} tiles, mostly disjoint and city-wide; the first keeps \texttt{ID\_}, \num{25} distinct buildings receive source-tagged synthesized ids, and \num{4} geometry-less placeholder copies are deduplicated as identical content. Separately, \num{16} genuine UUIDs are each duplicated \emph{within a single tile}: for instance \texttt{ID\_10C4883F-\ldots} at lines \num{2989} and \num{7530489} of tile \texttt{1091-23}, two near-coincident but distinct buildings, of which \num{15} are split and one geometry-less copy is deduplicated.

% \emph{Tokyo} duplicates at the file level rather than the identifier level: PLATEAU ships each ward as a separate package, so a mesh tile on a ward boundary recurs in every adjacent ward (\num{240} shared tiles; ward-span 2--4). \num{177} of these tiles are byte-identical across their wards and \num{63} differ only in ward-specific \texttt{uro:*} ADE data-quality metadata while the building \emph{geometry} is identical (verified: every shared building in the sampled differing tiles has coordinate-identical geometry). Both cases resolve to one node per building (byte-identical copies on the content hash, geometry-identical copies on the geometry signature; see \Cref{subsec:fidelity}), so Tokyo's \num{975545} duplicate building elements collapse without any phantom split. The released record lists all \num{288} redundant tile files with their wards and content digests.

\section{Source Thematic Property Inventory}\label{app:thematic}

Every property key that originates in each city's source CityGML (core attributes plus \texttt{gen:*} generic attributes), with the number of buildings carrying it, is inventoried in \Cref{tab:city_thematic_attrs}. These verbatim, multilingual source vocabularies (German ALKIS/AdV codes, Finnish, Swiss-German, a sparse English LoD1--2 schema, and Japanese PLATEAU keys) are the semantics a text-to-query model must bridge (\Cref{subsec:query_tasks}) and the attributes available for representation-learning tasks (\Cref{subsec:rl_tasks}).

\input{tables/city_thematic_attrs_table.tex}

\section{Question Suite Design and Statistics}\label{app:questions}

This section summarizes the question suite specification. The full document (all 84 template definitions with gold-query sketches, slot providers, sizing model, and freeze process) ships with the benchmark repository. \Cref{tab:categories}'s instance counts are the frozen, live-verified v1.0 numbers.
% , not design targets. Nineteen of the 84 templates (sixteen plus a further three on Helsinki) were added after cross-referencing the per-city property inventory of \Cref{app:thematic} against the initial 65-template set and finding properties with real, substantial coverage that no question touched; \Cref{app:questions_native} lists them.

\begin{table*}[h!]
  \caption{
    Question categories with template and instance counts, aggregated across all five cities.  
    % Question categories with template counts and live-verified instance counts (\Cref{subsec:construction}), pooled over all five released city graphs and all languages. Tier~2 categories exist on the Hamburg deep-fusion instance only. The multilingual row is a cross-cutting tag on templates already counted in the rows above (German/Japanese variants of ten cross-city templates plus three Tokyo schema-bridging templates), so it is not added into the total.
    }
  \label{tab:categories}
  % \footnotesize
  \begin{tabularx}{\textwidth}{lXcrr}
    \toprule
    Category & Capability probed & Tier & Tmpl. & Instances \\
    \midrule
    Aggregate / filter / top-k & schema grounding, retrieval, ranking, incl.\ native per-city attributes 
    %(floor area, storeys, volume, year built, basement storeys) 
              & 1 & 18 & 230 \\
    Multi-hop semantic & traversal over the compact schema, incl.\ native per-city attributes 
    % (building material, town code, lifecycle status) 
              & 1 & 11 & 105 \\
    Spatial & window, proximity, geometric multi-hop, 3D structure, density & 1 & 12 & 462 \\
    Cross-source & (dis)agreement, coverage gaps, address completeness, dual roof shapes, OSM-only
    % agreement, disagreement, gaps in both directions, address coverage, dual roof-shape reporting, and OSM-only supplementation where authoritative data is absent 
              & 1 & 13 & 237 \\
    Provenance-filtered & source and confidence constraints, provenance metadata
    %source and confidence constraints, incl.\ reconstruction confidence (Hamburg) and each source's own pre-fusion provenance metadata (Helsinki, Zurich) 
            & 1 & 12 & 162 \\
    Coverage-aware & partial-coverage reasoning, incl.\ multi-material predictions & 2 & 6 & 24 \\
    LoD3 showcase & facade/interior structure & 2 & 4 & 21 \\
    Infeasible & hallucination resistance (guard-verified) & 1 & 8 & 153 \\
    Multilingual & German/Japanese language variants; Japanese-key gold queries & 1 & 3 & 95 \\
    \midrule
    Total & & & 84 & \num{1394} \\
    \bottomrule
  \end{tabularx}

  \smallskip\raggedright\footnotesize
  Tier~2 categories are available only for Hamburg. \emph{Multilingual} is an annotation applied across existing templates (German, Japanese, etc.) and is therefore excluded from the total.
\end{table*}

\subsection{Template Inventory}\label{app:questions_inventory}

The categories of \Cref{tab:categories} are divided into capability subfamilies. The five distinctive categories are deliberately the deepest: spatial alone accounts for roughly a third of the feasible instances, and the five together for approximately 70\,\%. Templates marked cross-city use only the common attribute subset of \Cref{sec:dataset} and instantiate on all five cities. City-specific templates (ALKIS code lists, Japanese-keyed attributes, prediction and LoD3 layers) are tagged as such and instantiate only where their layers exist, a constraint the generator asserts at run time.

% \begin{table}
%   \caption{Template inventory by category and subfamily.}
%   \label{tab:app_inventory}
%   \footnotesize
%   \begin{tabularx}{\linewidth}{lXr}
%     \toprule
%     Category & Subfamilies and probes & Tmpl. \\
%     \midrule
%     Aggregate & counts, means, percentiles, dataset extent, native attributes (Helsinki floor area, construction year, storey count, volume) & 10 \\
%     Filter / top-k & thresholds, ranking, compound predicates, missing-value awareness, native attributes (Tokyo basement storeys, redevelopment districts) & 8 \\
%     Multi-hop & grouped traversal, part hierarchies, boundary-surface census, per-district top-1, native attributes (Helsinki material and lifecycle status, Tokyo town code) & 11 \\
%     Spatial & S1 window; S2 proximity and $k$-nearest; S3 geometric multi-hop (spatial $\times$ semantic); S4 3D structure; S5 pair density & 12 \\
%     Cross-source & value disagreement, gaps in both directions, floor-count conflicts, agreement rate, net-new content by kind, address coverage, dual roof-shape reporting, OSM-only height where authoritative is absent & 13 \\
%     Provenance & confidence thresholds, correspondence cases (1:1/1:n/n:1/n:m), source-restriction traps, source-qualified values, non-primary edges, LoD3 reconstruction confidence, native pre-fusion provenance (Helsinki integration method, Zurich data vintage and revision history) & 12 \\
%     Coverage & covered share, covered-subset distributions, city-wide count trap, per-district coverage uniformity, multi-material predictions & 6 \\
%     LoD3 showcase & windows, doors, openings per wall, rooms & 4 \\
%     Infeasible & absent attribute, absent layer, false premise, temporal, deliberately absent personal data & 8 \\
%     Multilingual & German/Japanese variants of ten cross-city templates; Japanese-key schema bridging (Tokyo) & 3 \\
%     \midrule
%     Total & & 84 \\
%     \bottomrule
%   \end{tabularx}
% \end{table}

% \subsection{Per-City Native-Content Extension}\label{app:questions_native}

% Cross-referencing the per-city property inventory (\Cref{app:thematic}) against the initial 65-template set surfaced properties with real coverage, from tens of thousands to millions of buildings, that no question touched: OSM-copied address fields; a second, roof-\emph{shape} cross-source signal alongside the existing roof-\emph{material} one; Hamburg's own reconstruction-confidence and -method scalars on the LoD3 layer; the multi-material tail of the roof-material predictions; several further Japanese-keyed Tokyo attributes beyond the two schema-bridging templates already in the suite; roughly thirty unused Finnish cadastral attributes on Helsinki, the attribute-richest city in the corpus (\Cref{tab:city_thematic_attrs}) and the least queried before this pass; and Zurich's native data-vintage and revision-history fields. One gap was structural rather than a missing property: every existing cross-source template required both an authoritative and a crowd-sourced value present, which silently excluded New York's defining fact (\num{97.5}\,\% OSM height coverage against zero authoritative height) from every cross-source question; the new template asks what the crowd-sourced layer alone reports precisely where the authoritative source is silent, verified to return a substantive answer (\num{1056741} of \num{1083437} buildings, mean \num{8.3}\,m). Two further additions surface provenance axes distinct from our own OSM fusion: Helsinki's pre-existing record-linkage metadata from its own prior data integration, and Zurich's native survey-vintage and revision fields, provenance about a source's own history, not only about our fusion of it. Per-city instantiation is governed by the declarative requirement gate (\Cref{subsec:construction}): each feasible template lists the graph tokens its gold query reads, and the generator instantiates it only where the city's measured schema provides them. This supersedes an earlier reliance on slot self-gating alone, which was safe for slotted templates but let a no-slot aggregate on a city lacking the property emit a degenerate null-or-zero row rather than being withheld, an unfair question that then had to be removed by hand and, being hand-applied per city, was applied inconsistently across cities; the gate removes both the inconsistency and the manual step, and is data-driven on token presence rather than a hardcoded per-city allow-list. A follow-up pass then addressed the corpus's most under-queried city relative to its richness: Helsinki carries the most source-unique attributes of any city (34, against 27 for Hamburg and 16 for Zurich) yet was among the least templated, so three further Helsinki templates were added, its native storey count (a field distinct from the pan-city storey attribute, which Helsinki populates on only 92 buildings), building volume, and building lifecycle status (active, planned, under construction, or demolished), bringing it to parity with Tokyo. A fourth candidate, the quantitative footprint match-quality logged by Helsinki's own prior record linkage, was authored first but dropped after live verification against the released graph showed the field, and the co-occurring fields it was meant to complement (matching method, footprint-area difference), constant across all \num{2846} buildings that carry them: a single recorded outcome, not a distribution, so an ``average'' framed as a quality measure would have been a degenerate restatement of an existing template's count rather than new signal; the building-status field, by contrast, exhibits a real four-way split (\num{1921}/\num{245}/\num{8}/\num{2}) and standard, unambiguous terminology, so it replaced the dropped candidate. Fields whose semantics require the data provider's dictionary, notably Helsinki's parallel roof-height datums (the plain versus \texttt{brec\_}-prefixed height variants), remain released in the dataset but deliberately excluded from v1.0 gold. All 84 templates, including the three Helsinki additions, the paraphrase expansion, and the German/Japanese language-variant mechanism (\Cref{subsec:construction}), have since passed live machine-verification (G1--G4) against the released graph of all five cities: \num{1394}/\num{1406} instantiated questions (\num{99.1}\,\%), every failure the deliberate feasibility flip (\Cref{subsec:query_tasks}), none attributable to the additions.

\subsection{Difficulty Rubric}\label{app:questions_difficulty}

Each question is assigned an authored difficulty level (\Cref{tab:app_difficulty}). 
% These labels remain fixed throughout the study. 
After completing the baseline evaluation matrix, we additionally report the observed difficulty of each question, measured by its failure rate across models, and quantify the correlation between authored and observed difficulty. Any discrepancy between the two is treated as an empirical finding rather than a reason to revise the original difficulty labels.
% Each question carries an authored design difficulty (\Cref{tab:app_difficulty}). The authored tags stay frozen; after the baseline matrix has run we additionally report the observed difficulty (per-question failure rates) and the correlation between the two, treating a divergence as a finding rather than an occasion to re-tag.
Gold answers exhibit a diverse set of result formats (\Cref{tab:app_answer_types}), determined by the structure of the corresponding gold query outputs. Single-value results are the most common, including aggregates, counts, and scalar lookups. These are followed by multi-column or grouped tables, ranked top-$k$ lists, and infeasible cases that require refusal. Notably, none of the benchmark templates produces a bare boolean answer.
% Gold answers span a range of result shapes (\Cref{tab:app_answer_types}), derived from each gold query's result structure: single scalars dominate (aggregates, counts, single-value lookups), followed by multi-column or grouped tables, ranked top-$k$ lists, and the infeasible refusals; no template returns a bare boolean.

\begin{table}
  \caption{Authored difficulty rubric and its distribution over the \num{1394} frozen question instances. Conditions of lower levels may also apply at higher ones.}
  \label{tab:app_difficulty}
  % \footnotesize
  \begin{tabularx}{\linewidth}{cXr}
    \toprule
    Level & Definition & Questions \\
    \midrule
    1 & one node label, one property, one aggregate; no filter & 47 (3.4\%) \\
    2 & one filter, one relationship hop, or ranking (\texttt{ORDER BY}/\texttt{LIMIT}) & 272 (19.5\%) \\
    3 & two composed elements: two hops, a grouped aggregate, one spatial predicate, one edge-property (provenance) filter, or a missing-value subtlety & 714 (51.2\%) \\
    4 & cross-concern composition: spatial $\times$ semantic, cross-source comparison with computation, coverage normalization, percentile or conditional share, schema bridging across languages & 247 (17.7\%) \\
    5 & three or more schema regions plus a geometric or set-level predicate: geometric multi-hop joins, per-entity grouped top-1 over two-hop patterns, n:m correspondence reasoning & 114 (8.2\%) \\
    \bottomrule
  \end{tabularx}
\end{table}

\begin{table}
  \caption{Gold answer-type distribution over the 1394 frozen question instances.}
  \label{tab:app_answer_types}
  % \footnotesize
  \begin{tabularx}{\linewidth}{Xr}
    \toprule
    Answer type & Questions \\
    \midrule
    Scalar, numeric & 556 (39.9\%) \\
    Table (multi-column or grouped) & 352 (25.3\%) \\
    Ranked list (top-$k$, ordered and limited) & 259 (18.6\%) \\
    Refusal (infeasible) & 153 (11.0\%) \\
    Scalar, string & 48 (3.4\%) \\
    List (single column) & 26 (1.9\%) \\
    \bottomrule
  \end{tabularx}
\end{table}

\subsection{Per-City Feasibility}\label{app:questions_matrix}

\Cref{tab:app_matrix} specifies which benchmark categories are instantiated for each city. Missing entries are intentional design choices.
% rather than coverage gaps. 
In particular, the coverage and LoD3 rows define the Tier~2 deep-fusion categories. Likewise, the absent-layer infeasible subfamily relies on the absence of prediction layers outside Hamburg, allowing the same question template to be feasible in one city and infeasible in another while keeping the question text unchanged.
% \Cref{tab:app_matrix} states which categories instantiate on which city. Absences are design features, not gaps: the coverage and LoD3 rows define the Tier~2 deep-fusion categories, and the absent-layer infeasible subfamily exists precisely because prediction layers are absent outside Hamburg, letting the same question text flip between feasible and infeasible across cities.

\begin{table}
  \caption{Category $\times$ city feasibility. HH Hamburg, ZH Zurich, HEL Helsinki, TYO Tokyo, NYC New York.}
  \label{tab:app_matrix}
  %\footnotesize
  \begin{tabular}{lccccc}
    \toprule
    Category & HH & ZH & HEL & TYO & NYC \\
    \midrule
    Aggregate / filter / multi-hop & \CIRCLE & \CIRCLE & \CIRCLE & \CIRCLE & \CIRCLE \\
    Spatial\textsuperscript{1} & \CIRCLE & \CIRCLE & \CIRCLE & \CIRCLE & \CIRCLE \\
    Cross-source / provenance\textsuperscript{2} & \CIRCLE & \CIRCLE & \CIRCLE & \CIRCLE & \CIRCLE \\
    Coverage-aware & \CIRCLE & -- & -- & -- & -- \\
    LoD3 showcase\textsuperscript{3} & \CIRCLE & -- & -- & -- & -- \\
    Infeasible (other subfamilies) & \CIRCLE & \CIRCLE & \CIRCLE & \CIRCLE & \CIRCLE \\
    Infeasible (absent layer) & -- & \CIRCLE & \CIRCLE & \CIRCLE & \CIRCLE \\
    Multilingual schema bridging & -- & -- & -- & \CIRCLE & -- \\
    Dual roof-shape reporting & \CIRCLE & \CIRCLE & -- & -- & -- \\
    LoD3 confidence scalars\textsuperscript{3} & \CIRCLE & -- & -- & -- & -- \\
    Native single-city content\textsuperscript{4} & -- & \CIRCLE & \CIRCLE & \CIRCLE & -- \\
    OSM-only height\textsuperscript{5} & \CIRCLE & \CIRCLE & \CIRCLE & \CIRCLE & \CIRCLE \\
    \bottomrule
  \end{tabular}
  \par\smallskip\raggedright\footnotesize
  \textsuperscript{1} After metric-CRS reprojection at ingest for Tokyo and New York (\Cref{subsec:construction}). 
  \textsuperscript{2} After whole-city OSM fusion. 
  \textsuperscript{3} Applicable only to selected buildings in Hamburg (\Cref{subsec:predicted_layers}). 
  \textsuperscript{4} Single-home-city templates, gated by \Cref{subsec:construction}'s requirement gate. Helsinki: floor area, storeys, volume, material, integration method, building status; Zurich: data vintage, revision history; Tokyo: basement storeys, town code, redevelopment district. Hamburg's own single-city content is the coverage and LoD3 rows above. 
  \textsuperscript{5} Cross-city by construction (no per-city gate needed), but only substantive where a real gap exists between the two sources' coverage, most pronounced for New York. 
  %(\Cref{app:questions_native}).
\end{table}

\subsection{Verification Gate}\label{app:questions_gate}

Every question passes five machine-checked validation gates before entering the benchmark suite. \textbf{G1} verifies that the gold query executes successfully without error. \textbf{G2} requires the canonicalized answer hash to be identical across three independent executions, ensuring result stability. \textbf{G3} checks that the result is non-empty unless the underlying template explicitly permits empty outputs. \textbf{G4}, applied to infeasible questions, executes a guard query that must demonstrate the absence of the requested information in the target city's graph by returning a count of zero. \textbf{G5} enforces benchmark versioning: whenever the dataset is rebuilt, all gold answers are re-materialized, and any change triggers a patch-version increment together with a changelog entry, preventing gold answers from silently drifting from the released graph. In addition to these automated checks, we perform human validation on a stratified sample comprising at least 20\% of questions, sampled across categories, difficulty levels, and cities. Each sampled question is independently reviewed 
%by two annotators 
for gold-answer correctness and natural language quality.
%, with disagreements resolved through adjudication. 
At benchmark freeze time, we document the rejection rates for each validation gate together with the human-review acceptance rates.
% Every question passes five machine-checked gates before entering the suite: G1 the gold query executes without error; G2 its canonicalized answer hash is identical across three runs; G3 the result is non-empty unless the template explicitly permits emptiness; G4 for infeasible questions, a guard query proves on the target city's graph that the asked-for information is absent (the guard must return a zero count); G5 on every dataset rebuild all gold answers are re-materialized, and any change bumps the suite patch version with a changelog entry, so gold answers can never silently drift from the released graph. On top of the machine gates, a stratified sample of at least 20\,\% (by category, difficulty, and city) is reviewed by two humans for gold correctness and natural phrasing, with disagreements adjudicated; per-gate rejection rates and review acceptance rates are published here at freeze time.

\subsection{Example Questions with Gold Queries}\label{app:questions_examples}

One representative example is provided for each distinctive capability. Placeholders enclosed in angle brackets are instantiated at generation time using values retrieved directly from the released graph. 
% The complete set of all \num{84} template gold queries is listed in \Cref{app:questions_all}.

\emph{Spatial (S1 window).} ``Which is the tallest building inside the window $\langle$wkt$\rangle$?''
\begin{lstlisting}
CALL spatial.intersects('features', '<wkt>') YIELD node
WITH node WHERE node:Building
  AND node.measured_height IS NOT NULL
RETURN node.id AS id, node.measured_height AS height_m
ORDER BY node.measured_height DESC, node.id ASC LIMIT 1
\end{lstlisting}

\emph{Cross-source (value disagreement).} ``Which buildings have an OSM height that differs from the surveyed height by more than $\langle d \rangle$ meters?''
\begin{lstlisting}
MATCH (b:Building)-[r:ENRICHED_BY {is_primary: true}]
      ->(o:OsmFeature)
WHERE o.osm_height IS NOT NULL
  AND b.measured_height > -999
  AND abs(toFloat(o.osm_height) - b.measured_height) > <d>
RETURN b.id AS id, b.measured_height AS surveyed_m,
       toFloat(o.osm_height) AS osm_m
ORDER BY abs(toFloat(o.osm_height) - b.measured_height)
  DESC, id ASC
\end{lstlisting}

\emph{Provenance (source-restriction trap).} ``Using only authoritative cadastral data, what is the average building height?'' The gold query reads the surveyed property only. A model that also averages the crowd-sourced \texttt{osm\_height} produces a plausibly close but wrong number.
%, which the structural match score exposes.
\begin{lstlisting}
MATCH (b:Building)
WHERE b.measured_height > -999
RETURN round(avg(b.measured_height) * 100) / 100.0
       AS avg_height_m
\end{lstlisting}

\emph{Coverage-aware (city-wide count trap).} ``How many buildings in Hamburg have a $\langle$material$\rangle$ roof?'' The gold answer couples the count with its coverage context. A bare count is scored as wrong.
\begin{lstlisting}
MATCH (b:Building)
WITH count(b) AS total,
  count(CASE WHEN b.predictedroofmaterial IS NOT NULL
        THEN 1 END) AS covered,
  count(CASE WHEN b.predictedroofmaterial = '<material>'
        THEN 1 END) AS matching
RETURN matching, covered, total,
  round(1000.0 * covered / total) / 10.0 AS coverage_pct
\end{lstlisting}

\emph{Infeasible (absent attribute).} ``In which year was each building last renovated?'' Gold behavior is an explicit refusal. The guard proves the absence on the target city:
\begin{lstlisting}
MATCH (b:Building)
WHERE b.renovation_year IS NOT NULL
RETURN count(b) AS n  // gate G4: must return n = 0
\end{lstlisting}

\emph{Multilingual (Japanese-key schema bridging, Tokyo).} ``Which buildings belong to the district plan `Ichigaya-Yanagich\=o'?'', also posed in Japanese (\texttt{lang: ja}) as \jp{「地区計画『市谷柳町地区』に属する建物はどれですか。」} The property key itself is Japanese, so the gold query must bridge the schema regardless of the question language:
\begin{lstlisting}[escapeinside={(*@}{@*)}]
MATCH (b:Building)
WHERE b.`(*@\jp{地区計画}@*)` = '(*@\jp{市谷柳町地区}@*)'
RETURN b.id AS id ORDER BY id
\end{lstlisting}

% \subsection{Complete Gold Query Listing}\label{app:questions_all}

% All \num{84} template gold queries, one per template, grouped by category and ordered as in \Cref{tab:categories}. \Cref{tab:app_query_overview} summarizes the count per query type. Slots shown as \lstinline|<slot>| are inlined with values read from the released graph at generation time (\Cref{subsec:construction}); a leading comment gives each template's identifier and a natural-language gloss. Infeasible templates list the guard query that must return zero for the question to be genuinely unanswerable. This listing is machine-generated from \texttt{bench/templates.py} by \texttt{bench/gen\_gold\_listing.py}.

% \begin{table}
%   \caption{Task~A gold-query overview: distinct gold-query templates and instantiated questions per query type (the categories of \Cref{tab:categories}), with a concise descriptor of what each type probes. The queries themselves follow in \Cref{app:questions_all}.}
%   \label{tab:app_query_overview}
%   \footnotesize
%   \begin{tabularx}{\linewidth}{Xrr}
%     \toprule
%     Query type and probes & Templates & Questions \\
%     \midrule
%     \textbf{Aggregate} counts, means, percentiles, dataset extent & 10 & \num{90} \\
%     \textbf{Filter / top-k} thresholds, top-$k$ ranking, compound filters & 8 & \num{140} \\
%     \textbf{Multi-hop} traversal over the compact schema & 11 & \num{105} \\
%     \textbf{Spatial} bounding box, proximity, geometric predicates & 12 & \num{462} \\
%     \textbf{Cross-source} OSM vs.\ authoritative agreement and disagreement & 13 & \num{237} \\
%     \textbf{Provenance} source, confidence, and primary-match constraints & 12 & \num{162} \\
%     \textbf{Coverage-aware} partial-coverage-aware aggregation (Hamburg) & 6 & \num{24} \\
%     \textbf{LoD3 showcase} facade and interior detail (Hamburg) & 4 & \num{21} \\
%     \textbf{Infeasible} unanswerable; gold behavior is to refuse & 8 & \num{153} \\
%     \midrule
%     \textbf{Total} & 84 & \num{1394} \\
%     \bottomrule
%   \end{tabularx}
% \end{table}

% \paragraph{Aggregate.}
% \begin{lstlisting}[basicstyle=\ttfamily\footnotesize,escapeinside={(*@}{@*)}]
% // agg_avg_height : What is the average measured building height in meters?
% MATCH (b:Building)
% WHERE b.measured_height > -999
% RETURN round(avg(b.measured_height) * 100) / 100.0 AS avg_height_m

% // agg_count_buildings : How many buildings does the dataset contain?
% MATCH (b:Building)
% RETURN count(b) AS n

% // agg_dataset_extent : What rectangular extent does the dataset cover, in kilometers?
% MATCH (ds:Dataset)
% RETURN round((ds.bbox_max_x - ds.bbox_min_x) / 100.0) / 10.0 AS width_km, round((ds.bbox_max_y - ds.bbox_min_y) / 100.0) / 10.0 AS height_km

% // agg_helsinki_avg_floor_area : What is the average total floor area (kerrosala) of the buildings, in square m...
% MATCH (b:Building)
% WHERE b.kerrosala IS NOT NULL
% RETURN round(avg(toFloat(b.kerrosala)) * 100) / 100.0 AS avg_floor_area_m2, count(b) AS n

% // agg_helsinki_avg_storeys : What is the average number of storeys (kerroksia) of the buildings, and how ma...
% MATCH (b:Building)
% WHERE b.kerroksia IS NOT NULL
% RETURN round(avg(toFloat(b.kerroksia)) * 100) / 100.0 AS avg_storeys, max(toInteger(b.kerroksia)) AS max_storeys, count(b) AS n

% // agg_helsinki_avg_volume : What is the average building volume (tilavuus) in cubic meters, and how many b...
% MATCH (b:Building)
% WHERE b.tilavuus IS NOT NULL
% RETURN round(avg(toFloat(b.tilavuus)) * 100) / 100.0 AS avg_volume_m3, count(b) AS n

% // agg_helsinki_avg_year_built : What is the average construction year (valmistunut) of the buildings?
% MATCH (b:Building)
% WHERE b.valmistunut IS NOT NULL
% RETURN round(avg(toInteger(b.valmistunut))) AS avg_year_built, count(b) AS n

% // agg_max_storeys : What is the highest number of storeys above ground of any building?
% MATCH (b:Building)
% WHERE b.storeys_above_ground IS NOT NULL
% RETURN max(b.storeys_above_ground) AS max_storeys

% // agg_median_height : What is the median measured building height?
% MATCH (b:Building)
% WHERE b.measured_height > -999
% RETURN round(percentileCont(b.measured_height, 0.5) * 100) / 100.0 AS median_height_m

% // tokyo_jp_numeric_attr : What is the average floor-area conversion factor across buildings?
% MATCH (b:Building)
% WHERE b.`(*@\jp{延べ面積換算係数}@*)` IS NOT NULL
% RETURN round(avg(toFloat(b.`(*@\jp{延べ面積換算係数}@*)`)) * 1000) / 1000.0 AS avg_factor
% \end{lstlisting}

% \paragraph{Filter / top-k.}
% \begin{lstlisting}[basicstyle=\ttfamily\footnotesize,escapeinside={(*@}{@*)}]
% // filter_by_function : How many buildings have the function '<function_name>'?
% MATCH (b:Building)-[:HAS_FUNCTION]->(f:BuildingFunction {name: '<function_name>'})
% RETURN count(b) AS n

% // filter_compound : How many buildings taller than <height_threshold> m have roof type '<roof_type...
% MATCH (b:Building)-[:HAS_ROOF_TYPE]->(r:RoofType {name: '<roof_type_name>'})
% WHERE b.measured_height > <height_threshold>
% RETURN count(b) AS n

% // filter_has_basement : How many buildings have at least one storey below ground?
% MATCH (b:Building)
% WHERE b.storeys_below_ground IS NOT NULL AND b.storeys_below_ground >= 1
% RETURN count(b) AS n

% // filter_missing_height : How many buildings have no recorded height?
% MATCH (b:Building)
% WHERE b.measured_height IS NULL
% RETURN count(b) AS n

% // filter_redevelopment_district : How many buildings belong to a legally-designated redevelopment-promotion dist...
% MATCH (b:Building)
% WHERE b.`(*@\jp{再開発等促進区を定める地区計画}@*)` IS NOT NULL
% RETURN count(b) AS n

% // filter_taller_than : How many buildings are taller than <height_threshold> meters?
% MATCH (b:Building)
% WHERE b.measured_height > <height_threshold>
% RETURN count(b) AS n

% // topk_most_storeys : List the <topk> buildings with the most storeys above ground.
% MATCH (b:Building)
% WHERE b.storeys_above_ground IS NOT NULL
% RETURN b.id AS id, b.storeys_above_ground AS storeys
% ORDER BY storeys DESC, id ASC LIMIT <topk>

% // topk_tallest : List the <topk> tallest buildings with their heights.
% MATCH (b:Building)
% WHERE b.measured_height IS NOT NULL
% RETURN b.id AS id, b.measured_height AS height_m
% ORDER BY b.measured_height DESC, b.id ASC LIMIT <topk>
% \end{lstlisting}

% \paragraph{Multi-hop.}
% \begin{lstlisting}[basicstyle=\ttfamily\footnotesize,escapeinside={(*@}{@*)}]
% // multihop_boundary_census : How many boundary surfaces of each type does building <any_building> have?
% MATCH (:Building {id: '<any_building>'})-[:HAS_BOUNDARY]->(s:BoundarySurface)
% RETURN s.surface_type AS surface_type, count(s) AS n
% ORDER BY n DESC, surface_type ASC

% // multihop_building_parts : Which buildings consist of more than one building part?
% MATCH (b:Building)-[:HAS_BUILDING_PART]->(:BuildingPart)
% WITH b, count(*) AS n_parts
% WHERE n_parts >= 2
% RETURN b.id AS id, n_parts
% ORDER BY n_parts DESC, id ASC

% // multihop_district_counts : How many buildings are located in each district?
% MATCH (b:Building)-[:LOCATED_IN]->(d:District)
% RETURN d.id AS district, count(b) AS n
% ORDER BY n DESC, district ASC

% // multihop_district_top_function : What is the most common building function in district <district>?
% MATCH (b:Building)-[:LOCATED_IN]->(:District {id: '<district>'})
% MATCH (b)-[:HAS_FUNCTION]->(f:BuildingFunction)
% RETURN f.name AS function, count(b) AS n
% ORDER BY n DESC, function ASC LIMIT 1

% // multihop_function_avg_height : What is the average building height per building function?
% MATCH (b:Building)-[:HAS_FUNCTION]->(f:BuildingFunction)
% WHERE b.measured_height > -999
% RETURN f.name AS function, round(avg(b.measured_height) * 100) / 100.0 AS avg_height_m, count(b) AS n
% ORDER BY n DESC, function ASC

% // multihop_helsinki_building_status : What is the distribution of building lifecycle status (rakennuksen_tila) acros...
% MATCH (b:Building)
% WHERE b.rakennuksen_tila IS NOT NULL
% RETURN b.rakennuksen_tila AS status, count(b) AS n
% ORDER BY n DESC, status ASC

% // multihop_helsinki_material_distribution : What is the distribution of building materials (rakennusaine) across buildings...
% MATCH (b:Building)
% WHERE b.rakennusaine IS NOT NULL
% RETURN b.rakennusaine AS material, count(b) AS n
% ORDER BY n DESC, material ASC

% // multihop_jp_town_code_top : Which 10 town/block codes ((*@\jp{町丁目コード}@*)) have the most buildings?
% MATCH (b:Building)
% WHERE b.`(*@\jp{町}@*)_(*@\jp{丁目コード}@*)` IS NOT NULL
% RETURN b.`(*@\jp{町}@*)_(*@\jp{丁目コード}@*)` AS town_code, count(b) AS n
% ORDER BY n DESC, town_code ASC LIMIT 10

% // multihop_rooftype_distribution : Show the distribution of roof types across all buildings.
% MATCH (b:Building)-[:HAS_ROOF_TYPE]->(r:RoofType)
% RETURN r.name AS roof_type, count(b) AS n
% ORDER BY n DESC, roof_type ASC

% // tokyo_jp_key_distribution : How many buildings are there per district plan?
% MATCH (b:Building)
% WHERE b.`(*@\jp{地区計画}@*)` IS NOT NULL
% RETURN b.`(*@\jp{地区計画}@*)` AS district_plan, count(b) AS n
% ORDER BY n DESC, district_plan ASC

% // tokyo_jp_key_lookup : Which buildings belong to the district plan '<jp_district_plan>'?
% MATCH (b:Building)
% WHERE b.`(*@\jp{地区計画}@*)` = '<jp_district_plan>'
% RETURN b.id AS id
% ORDER BY b.id
% \end{lstlisting}

% \paragraph{Spatial.}
% \begin{lstlisting}[basicstyle=\ttfamily\footnotesize,escapeinside={(*@}{@*)}]
% // spatial_bbox_agg : What is the average building height inside the window <bbox_wkt>?
% CALL spatial.intersects('features', '<bbox_wkt>') YIELD node
% WITH node
% WHERE node:Building AND node.measured_height > -999
% RETURN round(avg(node.measured_height) * 100) / 100.0 AS avg_height_m, count(node) AS n

% // spatial_bbox_by_rooftype : Within the window <bbox_wkt>, how many buildings are there per roof type?
% CALL spatial.intersects('features', '<bbox_wkt>') YIELD node
% WITH node
% WHERE node:Building
% MATCH (node)-[:HAS_ROOF_TYPE]->(r:RoofType)
% RETURN r.name AS roof_type, count(node) AS n
% ORDER BY n DESC, roof_type ASC

% // spatial_bbox_count : How many buildings lie inside the bounding box <bbox_wkt>?
% CALL spatial.intersects('features', '<bbox_wkt>') YIELD node
% WITH node
% WHERE node:Building
% RETURN count(node) AS n

% // spatial_bbox_tallest : Which is the tallest building inside the bounding box <bbox_wkt>?
% CALL spatial.intersects('features', '<bbox_wkt>') YIELD node
% WITH node
% WHERE node:Building AND node.measured_height IS NOT NULL
% RETURN node.id AS id, node.measured_height AS height_m
% ORDER BY node.measured_height DESC, node.id ASC LIMIT 1

% // spatial_dense_pairs : Within the window <small_bbox_wkt>, how many pairs of buildings stand less tha...
% CALL spatial.intersects('features', '<small_bbox_wkt>') YIELD node
% WITH node
% WHERE node:Building
% WITH collect(node) AS bldgs
% UNWIND bldgs AS a
% UNWIND bldgs AS b
% WITH a, b
% WHERE a.id < b.id AND point.distance(a.location, b.location) < 10
% RETURN count(*) AS n

% // spatial_isolated : Within the window <small_bbox_wkt>, how many buildings have no other building ...
% CALL spatial.intersects('features', '<small_bbox_wkt>') YIELD node
% WITH node
% WHERE node:Building
% WITH collect(node) AS bldgs
% UNWIND bldgs AS a
% WITH a, size([x IN bldgs
% WHERE x.id <> a.id AND point.distance(a.location, x.location) < 100]) AS n_near
% WHERE n_near = 0
% RETURN count(a) AS n

% // spatial_knearest : Which <topk> buildings are closest to building <nearby_building>?
% MATCH (b0:Building {id: '<nearby_building>'})
% CALL spatial.withinDistance('features', b0.location, 0.5) YIELD node
% WITH b0, node
% WHERE node:Building AND node.id <> b0.id
% RETURN node.id AS id, round(point.distance(b0.location, node.location) * 100) / 100.0 AS distance_m
% ORDER BY distance_m ASC, id ASC LIMIT <topk>

% // spatial_near_amenity_tallest : What is the tallest building within <radius_m> m of a <amenity> point of inter...
% MATCH (o:OsmFeature {osm_amenity: '<amenity>'})
% CALL spatial.withinDistance('features', o.location, <radius_m> / 1000.0) YIELD node
% WITH node
% WHERE node:Building AND node.measured_height IS NOT NULL
% RETURN node.id AS id, node.measured_height AS height_m
% ORDER BY node.measured_height DESC, node.id ASC LIMIT 1

% // spatial_near_water : How many buildings stand within <radius_m> m of the water feature at (<water_c...
% WITH split('<water_centroid_xy>', ',') AS xy
% WITH point({x: toFloat(xy[0]), y: toFloat(xy[1])}) AS p
% CALL spatial.withinDistance('features', p, <radius_m> / 1000.0) YIELD node
% WITH node
% WHERE node:Building
% RETURN count(DISTINCT node) AS n

% // spatial_neighbors : How many buildings stand within 100 meters of building <nearby_building>?
% MATCH (b0:Building {id: '<nearby_building>'})
% CALL spatial.withinDistance('features', b0.location, 0.1) YIELD node
% WITH node
% WHERE node:Building AND node.id <> b0.id
% RETURN count(node) AS n

% // spatial_roof_complexity : Which <topk> buildings have the most roof polygons?
% MATCH (b:Building)-[:HAS_BOUNDARY]->(:BoundarySurface {surface_type: 'RoofSurface'})-[:HAS_POLYGON]->(p:GeometryPolygon)
% WITH b, count(p) AS n_polygons
% RETURN b.id AS id, n_polygons
% ORDER BY n_polygons DESC, id ASC LIMIT <topk>

% // spatial_within_distance_point : How many buildings lie within <radius_m> m of the point (<building_centroid_xy...
% WITH split('<building_centroid_xy>', ',') AS xy
% WITH point({x: toFloat(xy[0]), y: toFloat(xy[1])}) AS p
% CALL spatial.withinDistance('features', p, <radius_m> / 1000.0) YIELD node
% WITH node
% WHERE node:Building
% RETURN count(node) AS n
% \end{lstlisting}

% \paragraph{Cross-source.}
% \begin{lstlisting}[basicstyle=\ttfamily\footnotesize,escapeinside={(*@}{@*)}]
% // xsrc_addr_coverage : For how many buildings is a street address known?
% MATCH (b:Building)
% WITH count(b) AS total, count(b.addr_street) AS with_addr
% RETURN with_addr, total, round(1000.0 * with_addr / total) / 10.0 AS coverage_pct

% // xsrc_agreement_rate : For buildings where both sources report a height, what share agree within 1 me...
% MATCH (b:Building)-[:ENRICHED_BY {is_primary: true}]->(o:OsmFeature)
% WHERE b.measured_height > -999 AND o.osm_height IS NOT NULL
% WITH count(*) AS total, count(CASE WHEN abs(toFloat(o.osm_height) - b.measured_height) <= 1.0 THEN 1 END) AS agree
% RETURN agree, total, CASE WHEN total = 0 THEN null ELSE round(1000.0 * agree / total) / 10.0 END AS agreement_pct

% // xsrc_buildings_by_street : Which buildings are on street '<street_name>'?
% MATCH (b:Building)
% WHERE b.addr_street = '<street_name>'
% RETURN b.id AS id
% ORDER BY b.id

% // xsrc_dual_roof_reporting : How many buildings have both an authoritative roof type and a crowd-sourced OS...
% MATCH (b:Building)
% WHERE b.roof_type_code IS NOT NULL AND b.osm_roof_shape IS NOT NULL
% RETURN count(b) AS n

% // xsrc_height_disagreement : Which buildings have an OSM height that differs from the surveyed height by mo...
% MATCH (b:Building)-[r:ENRICHED_BY {is_primary: true}]->(o:OsmFeature)
% WHERE o.osm_height IS NOT NULL AND b.measured_height > -999 AND abs(toFloat(o.osm_height) - b.measured_height) > <deviation_m>
% RETURN b.id AS id, b.measured_height AS surveyed_m, toFloat(o.osm_height) AS osm_m
% ORDER BY abs(toFloat(o.osm_height) - b.measured_height) DESC, id ASC

% // xsrc_levels_vs_storeys : Which buildings report a different floor count in OSM than in the cadastre?
% MATCH (b:Building)
% WHERE b.osm_building_levels IS NOT NULL AND b.storeys_above_ground IS NOT NULL AND toInteger(b.osm_building_levels) <> b.storeys_above_ground
% RETURN b.id AS id, b.storeys_above_ground AS cadastre_storeys, toInteger(b.osm_building_levels) AS osm_levels
% ORDER BY id ASC

% // xsrc_mean_abs_deviation : What is the mean absolute difference between OSM and surveyed building heights...
% MATCH (b:Building)-[:ENRICHED_BY {is_primary: true}]->(o:OsmFeature)
% WHERE o.osm_height IS NOT NULL AND b.measured_height > -999
% RETURN round(avg(abs(toFloat(o.osm_height) - b.measured_height)) * 100) / 100.0 AS mean_abs_dev_m, count(*) AS n_pairs

% // xsrc_name_lookup : What is the OSM name of building <named_osm_building>?
% MATCH (:Building {id: '<named_osm_building>'})-[:ENRICHED_BY {is_primary: true}]->(o:OsmFeature)
% RETURN o.name AS osm_name

% // xsrc_net_new_by_kind : What does OpenStreetMap add that the cadastre lacks? Count unmatched OSM featu...
% MATCH (o:OsmFeature)
% WHERE o.matched = false
% RETURN o.osm_kind AS kind, count(o) AS n
% ORDER BY n DESC, kind ASC

% // xsrc_osm_only_height_stats : For buildings with no authoritative height on record, what does OpenStreetMap ...
% MATCH (b:Building)
% WHERE (b.measured_height IS NULL OR b.measured_height <= -999) AND b.osm_height IS NOT NULL
% RETURN round(avg(toFloat(b.osm_height)) * 100) / 100.0 AS avg_osm_height_m, count(b) AS n

% // xsrc_pois_in_building : Which points of interest are located inside building <poi_building>?
% MATCH (:Building {id: '<poi_building>'})-[:HAS_POI]->(p:OsmFeature)
% RETURN p.osm_id AS osm_id, p.name AS name, p.osm_amenity AS amenity
% ORDER BY osm_id ASC

% // xsrc_reverse_gap : How many OSM buildings have no cadastral counterpart?
% MATCH (o:OsmFeature {osm_kind: 'building'})
% WHERE o.matched = false
% RETURN count(o) AS n

% // xsrc_unenriched_gap : How many buildings have no OpenStreetMap counterpart at all?
% MATCH (b:Building)
% WHERE NOT (b)-[:ENRICHED_BY]->(:OsmFeature)
% RETURN count(b) AS n
% \end{lstlisting}

% \paragraph{Provenance.}
% \begin{lstlisting}[basicstyle=\ttfamily\footnotesize,escapeinside={(*@}{@*)}]
% // prov_helsinki_integration_method : How are the buildings distributed across the data-matching methods (matching_m...
% MATCH (b:Building)
% WHERE b.matching_mode IS NOT NULL
% RETURN b.matching_mode AS matching_mode, count(b) AS n
% ORDER BY n DESC, matching_mode ASC

% // prov_high_confidence : How many buildings have a primary OSM match with a Jaccard overlap of at least...
% MATCH (b:Building)-[r:ENRICHED_BY {is_primary: true}]->(:OsmFeature)
% WHERE r.jaccard >= <jaccard_min>
% RETURN count(DISTINCT b) AS n

% // prov_lod3_confidence : What is the reconstruction confidence and method for building <lod3_building>'...
% MATCH (b:Building {id: '<lod3_building>'})
% WHERE b.lod3_confidence IS NOT NULL
% RETURN b.lod3_confidence AS confidence, b.lod3_method AS method, b.lod3_source AS source

% // prov_lod3_high_confidence_count : How many LoD3-reconstructed buildings have a reconstruction confidence of at l...
% MATCH (b:Building)
% WHERE b.lod3_confidence >= <confidence_threshold>
% RETURN count(b) AS n

% // prov_match_cases : How are the OSM-to-cadastre building matches distributed over the corresponden...
% MATCH (b:Building)
% WHERE b.osm_match_type IS NOT NULL
% RETURN b.osm_match_type AS match_case, count(b) AS n
% ORDER BY n DESC, match_case ASC

% // prov_multi_matched : How many buildings are matched by two or more OSM features?
% MATCH (b:Building)
% WHERE b.osm_match_count >= 2
% RETURN count(b) AS n

% // prov_only_authoritative : Using only authoritative cadastral data, what is the average building height?
% MATCH (b:Building)
% WHERE b.measured_height > -999
% RETURN round(avg(b.measured_height) * 100) / 100.0 AS avg_height_m

% // prov_secondary_edges : How many non-primary correspondence edges does the fusion graph keep?
% MATCH (:Building)-[r:ENRICHED_BY {is_primary: false}]->(:OsmFeature)
% RETURN count(r) AS n

% // prov_source_qualified_value : What is the height of building <dual_height_building> according to each availa...
% MATCH (b:Building {id: '<dual_height_building>'})-[r:ENRICHED_BY {is_primary: true}]->(o:OsmFeature)
% RETURN b.measured_height AS authoritative_height_m, toFloat(o.osm_height) AS osm_height_m, r.jaccard AS match_jaccard

% // prov_weakest_matches : List the <topk> least confident primary OSM matches.
% MATCH (b:Building)-[r:ENRICHED_BY {is_primary: true}]->(:OsmFeature)
% WHERE r.jaccard IS NOT NULL
% RETURN b.id AS id, round(r.jaccard * 1000) / 1000.0 AS jaccard
% ORDER BY r.jaccard ASC, b.id ASC LIMIT <topk>

% // prov_zurich_data_vintage : What is the distribution of source survey years (herkunft_jahr) across the bui...
% MATCH (b:Building)
% WHERE b.herkunft_jahr IS NOT NULL
% RETURN b.herkunft_jahr AS source_year, count(b) AS n
% ORDER BY n DESC, source_year ASC

% // prov_zurich_revised_count : How many buildings have a recorded reason for revision (grund_aenderung), i.e....
% MATCH (b:Building)
% WHERE b.grund_aenderung IS NOT NULL
% RETURN count(b) AS n
% \end{lstlisting}

% \paragraph{Coverage-aware.}
% \begin{lstlisting}[basicstyle=\ttfamily\footnotesize,escapeinside={(*@}{@*)}]
% // cov_by_district : Report the roof-material prediction coverage per district.
% MATCH (b:Building)-[:LOCATED_IN]->(d:District)
% WITH d, count(b) AS n_total, count(CASE WHEN b.predictedroofmaterial IS NOT NULL THEN 1 END) AS n_covered
% RETURN d.id AS district, n_covered, n_total, round(1000.0 * n_covered / n_total) / 10.0 AS coverage_pct
% ORDER BY coverage_pct DESC, district ASC

% // cov_material_avg_height : What is the average height of buildings with predicted roof material '<materia...
% MATCH (b:Building)
% WHERE b.predictedroofmaterial = '<material>' AND b.measured_height > -999
% RETURN round(avg(b.measured_height) * 100) / 100.0 AS avg_height_m, count(b) AS n

% // cov_material_distribution : Show the distribution of predicted roof materials over the buildings that have...
% MATCH (b:Building)
% WHERE b.predictedroofmaterial IS NOT NULL
% RETURN b.predictedroofmaterial AS material, count(b) AS n
% ORDER BY n DESC, material ASC

% // cov_multi_material_count : How many buildings have more than one predicted roof material?
% MATCH (b:Building)
% WHERE b.predicted_roof_material_count >= 2
% RETURN count(b) AS n

% // cov_share : What share of buildings has a predicted roof material at all?
% MATCH (b:Building)
% WITH count(b) AS total, count(CASE WHEN b.predictedroofmaterial IS NOT NULL THEN 1 END) AS covered
% RETURN covered, total, round(1000.0 * covered / total) / 10.0 AS coverage_pct

% // cov_trap_citywide : How many buildings have a predicted '<material>' roof?
% MATCH (b:Building)
% WITH count(b) AS n_total, count(CASE WHEN b.predictedroofmaterial IS NOT NULL THEN 1 END) AS n_covered, count(CASE WHEN b.predictedroofmaterial = '<material>' THEN 1 END) AS n_matching
% RETURN n_matching, n_covered, n_total, round(1000.0 * n_covered / n_total) / 10.0 AS coverage_pct
% \end{lstlisting}

% \paragraph{LoD3 showcase.}
% \begin{lstlisting}[basicstyle=\ttfamily\footnotesize,escapeinside={(*@}{@*)}]
% // lod3_door_count : How many doors does building <lod3_building> have?
% MATCH (:Building {id: '<lod3_building>'})-[:HAS_LOD3_FACADE]->(:BoundarySurface)-[:HAS_OPENING]->(o:Opening {opening_type: 'Door'})
% RETURN count(o) AS n

% // lod3_openings_per_wall : Which wall of building <lod3_building> has the most openings?
% MATCH (:Building {id: '<lod3_building>'})-[:HAS_LOD3_FACADE]->(w:BoundarySurface {surface_type: 'WallSurface'})-[:HAS_OPENING]->(o:Opening)
% WITH w, count(o) AS n_openings
% RETURN w.id AS wall_id, n_openings
% ORDER BY n_openings DESC, wall_id ASC LIMIT 1

% // lod3_room_count : How many rooms does building <room_building> contain?
% MATCH (:Building {id: '<room_building>'})-[:HAS_INTERIOR_ROOM]->(r:Room)
% RETURN count(r) AS n

% // lod3_window_count : How many windows does building <lod3_building> have?
% MATCH (:Building {id: '<lod3_building>'})-[:HAS_LOD3_FACADE]->(:BoundarySurface)-[:HAS_OPENING]->(o:Opening {opening_type: 'Window'})
% RETURN count(o) AS n
% \end{lstlisting}

% \paragraph{Infeasible.}
% \begin{lstlisting}[basicstyle=\ttfamily\footnotesize,escapeinside={(*@}{@*)}]
% // inf_energy_class : How many buildings have energy efficiency class A?
% // REFUSAL required; guard query must return 0:
% MATCH (b:Building)
% WHERE b.energy_class IS NOT NULL
% RETURN count(b) AS n

% // inf_false_premise : How many buildings are in the district '<fake_district>'?
% // REFUSAL required; guard query must return 0:
% MATCH (d:District {id: '<fake_district>'})
% RETURN count(d) AS n

% // inf_interior_lod : How many rooms does building <any_building> have?
% // REFUSAL required; guard query must return 0:
% MATCH (:Building)-[:HAS_INTERIOR_ROOM]->(:Room)
% RETURN count(*) AS n

% // inf_layer_absent_city : What roof material is predicted for building <any_building>?
% // REFUSAL required; guard query must return 0:
% MATCH (b:Building)
% WHERE b.predictedroofmaterial IS NOT NULL
% RETURN count(b) AS n

% // inf_owner : Who owns building <any_building>?
% // REFUSAL required; guard query must return 0:
% MATCH (b:Building)
% WHERE b.owner IS NOT NULL
% RETURN count(b) AS n

% // inf_renovation_year : In which year was each building last renovated?
% // REFUSAL required; guard query must return 0:
% MATCH (b:Building)
% WHERE b.renovation_year IS NOT NULL
% RETURN count(b) AS n

% // inf_residents : How many people live in each building?
% // REFUSAL required; guard query must return 0:
% MATCH (b:Building)
% WHERE b.resident_count IS NOT NULL
% RETURN count(b) AS n

% // inf_temporal : How much did the average building height change between 2020 and 2024?
% // REFUSAL required; guard query must return 0:
% MATCH (ds:Dataset)
% WHERE ds.snapshot_year IS NOT NULL OR ds.version_year IS NOT NULL
% RETURN count(ds) AS n
% \end{lstlisting}

% \section{Per-Query and Per-Task Results}\label{app:full_results_a}
\section{LLM Setup for Text-to-Query}\label{app:full_results_a}

Both query-generation models are evaluated under identical, deliberately \emph{closed-book} conditions. Each question prompt contains only the natural-language question, the project's Cypher and spatial-query rules, and a fixed per-city schema string consisting of the cached APOC-sampled graph schema together with the documented schema-gap patches (\Cref{subsec:experiments_setup}). No live database connection or external tool access is available. The schema string is byte-identical for both models within a given city, ensuring that neither model is exposed to information unavailable to the other. Consequently, a model is expected to refuse requests for attributes that are absent from the provided schema text. Outputs are cached and evaluated offline. Generated Cypher queries are executed against the released graph and compared with gold answers after result-set canonicalization: row multisets are compared with column names removed and floating-point values rounded to $10^{-6}$. As a result, adding an additional column or renaming an existing column alters the row signature and is therefore treated as incorrect rather than silently accepted. 

The model qwen2.5-coder:7b is served locally through Ollama using the \texttt{Q4\_K\_M} quantization, a pinned model tag, greedy decoding (temperature $=0$), the default context window, and a single generation per question. Claude Sonnet~5 is evaluated as a closed-book Claude Code subagent, with the model itself acting as the query generator rather than through the \texttt{ChatAnthropic} API endpoint. No \texttt{ANTHROPIC\_API\_KEY}, live database access, or external tools are available beyond the supplied schema text. The evaluated model snapshot is \texttt{claude-sonnet-5}, tested on 2026-07-19 using the \emph{high} reasoning-effort setting. 

To improve token efficiency, questions within a split are submitted in batches of approximately 25 rather than as individual calls. Each batch prompt explicitly instructs the model to answer every question independently, relying only on the provided schema and rules and without influence from other questions in the batch. This design trades strict per-call isolation for a substantial reduction in the number of model calls. While answering many related questions in a single context may plausibly increase internal consistency relative to fully independent calls, this effect was not measured and no claim is made that the setup is equivalent to one-call-per-question evaluation. 
%For reproducibility, the exact byte-for-byte batch prompts (including the complete schema string) and the raw model outputs for every city and split are released with the benchmark suite under \texttt{bench/out/eval\_prompts\_*\_subagent/}. 

% Task~A results are reported in aggregate by city, with development and held-out test performance shown side by side in \Cref{tab:results_query}. The complete per-category breakdown underlying these aggregates is presented in \Cref{tab:app_results_models}, which provides one development-split table per model. In each case, the overall feasible EX (last row) reproduces the corresponding EX value reported in \Cref{tab:results_query}.

% Both tables are machine-generated from the stored per-question result files by \texttt{bench/gen\_result\_matrix.py}, so they cannot drift from the released results.

% \begin{table}
%   \centering
%   \caption{Task~A per-category execution accuracy for Claude Sonnet 5 (closed-book), all five Tier-1 city dev splits (the per-category breakdown of \Cref{tab:results_query}'s aggregate EX). Each cell is EX\% with the category's question count for that city in parentheses; \textemdash{} marks a category with no admitted question on that city. Category order and definitions follow \Cref{tab:categories}. Per-city question counts differ because a city omits native-content and prediction-layer templates its source schema cannot answer (\Cref{subsec:experiments_setup}). The \emph{Infeasible} row scores a correct refusal; its per-city F1 against over-refusals is in \Cref{tab:results_query}.}
%   \label{tab:app_results_claude}
%   % \footnotesize
%   \begin{tabular}{lrrrrr}
%     \toprule
%     Query category & Hamburg & Helsinki & Zurich & NYC & Tokyo \\
%     \midrule
%     Aggregate & 83.3~(6) & 50.0~(8) & 83.3~(6) & 66.7~(3) & 83.3~(6) \\
%     Filter / top-k & 66.7~(15) & 66.7~(15) & 50.0~(8) & \textemdash & 60.0~(10) \\
%     Multi-hop & 66.7~(6) & 33.3~(6) & 100.0~(3) & 66.7~(3) & 50.0~(8) \\
%     Spatial & 75.0~(24) & 75.0~(24) & 72.7~(22) & 75.0~(16) & 54.5~(22) \\
%     Cross-source & 40.0~(15) & 40.0~(15) & 35.7~(14) & 62.5~(8) & 45.5~(11) \\
%     Provenance & 30.8~(13) & 50.0~(12) & 58.3~(12) & 66.7~(9) & 44.4~(9) \\
%     Coverage-aware & 50.0~(6) & 0.0~(2) & \textemdash & \textemdash & \textemdash \\
%     LoD3 showcase & 100.0~(4) & \textemdash & \textemdash & \textemdash & 0.0~(1) \\
%     Infeasible & 90.0~(10) & 90.0~(10) & 90.0~(10) & 90.0~(10) & 90.0~(10) \\
%     \midrule
%     \textbf{Overall (feasible EX)} & 60.7~(89) & 56.1~(82) & 61.5~(65) & 69.2~(39) & 53.7~(67) \\
%     \bottomrule
%   \end{tabular}
% \end{table}

% \begin{table}
%   \centering
%   \caption{Task~A per-category execution accuracy for qwen2.5-coder:7b (local, Q4\_K\_M), all five Tier-1 city dev splits (the per-category breakdown of \Cref{tab:results_query}'s aggregate EX). Each cell is EX\% with the category's question count for that city in parentheses; \textemdash{} marks a category with no admitted question on that city. Category order and definitions follow \Cref{tab:categories}. Per-city question counts differ because a city omits native-content and prediction-layer templates its source schema cannot answer (\Cref{subsec:experiments_setup}). The \emph{Infeasible} row scores a correct refusal; its per-city F1 against over-refusals is in \Cref{tab:results_query}.}
%   \label{tab:app_results_qwen}
%   % \footnotesize
%   \begin{tabular}{lrrrrr}
%     \toprule
%     Query category & Hamburg & Helsinki & Zurich & NYC & Tokyo \\
%     \midrule
%     Aggregate & 33.3~(6) & 25.0~(8) & 33.3~(6) & 66.7~(3) & 33.3~(6) \\
%     Filter / top-k & 6.7~(15) & 20.0~(15) & 0.0~(8) & \textemdash & 0.0~(10) \\
%     Multi-hop & 33.3~(6) & 50.0~(6) & 0.0~(3) & 0.0~(3) & 0.0~(8) \\
%     Spatial & 16.7~(24) & 0.0~(24) & 0.0~(22) & 0.0~(16) & 0.0~(22) \\
%     Cross-source & 26.7~(15) & 26.7~(15) & 28.6~(14) & 25.0~(8) & 18.2~(11) \\
%     Provenance & 0.0~(13) & 0.0~(12) & 0.0~(12) & 0.0~(9) & 0.0~(9) \\
%     Coverage-aware & 0.0~(6) & 0.0~(2) & \textemdash & \textemdash & \textemdash \\
%     LoD3 showcase & 100.0~(4) & \textemdash & \textemdash & \textemdash & 0.0~(1) \\
%     Infeasible & 0.0~(10) & 0.0~(10) & 0.0~(10) & 0.0~(10) & 0.0~(10) \\
%     \midrule
%     \textbf{Overall (feasible EX)} & 19.1~(89) & 14.6~(82) & 9.2~(65) & 10.3~(39) & 6.0~(67) \\
%     \bottomrule
%   \end{tabular}
% \end{table}

% \begin{table*} 
%     \centering 
%     \caption{Task~A per-category execution accuracy (EX\%) across all five city development splits for Claude Sonnet 5 (closed-book) and qwen2.5-coder:7b (local, Q4\_K\_M). Each cell reports EX\% with the category's question count in parentheses.} 
%     \label{tab:app_results_models} 
%     % \small 
%     \begin{tabular}{lccccc ccccc} 
%         \toprule & \multicolumn{5}{c}{Claude Sonnet 5} & \multicolumn{5}{c}{qwen2.5-coder:7b} \\ \cmidrule(lr){2-6}\cmidrule(lr){7-11} 
%         Query category & Hamburg & Helsinki & Zurich & NYC & Tokyo & Hamburg & Helsinki & Zurich & NYC & Tokyo \\ \midrule 
%         Aggregate & 83.3 (6) & 50.0 (8) & 83.3 (6) & 66.7 (3) & 83.3 (6) & 33.3 (6) & 25.0 (8) & 33.3 (6) & 66.7 (3) & 33.3 (6) \\ 
%         Filter / top-k & 66.7 (15) & 66.7 (15) & 50.0 (8) & \textemdash & 60.0 (10) & 6.7 (15) & 20.0 (15) & 0.0 (8) & \textemdash & 0.0 (10) \\ 
%         Multi-hop & 66.7 (6) & 33.3 (6) & 100.0 (3) & 66.7 (3) & 50.0 (8) & 33.3 (6) & 50.0 (6) & 0.0 (3) & 0.0 (3) & 0.0 (8) \\ Spatial & 75.0 (24) & 75.0 (24) & 72.7 (22) & 75.0 (16) & 54.5 (22) & 16.7 (24) & 0.0 (24) & 0.0 (22) & 0.0 (16) & 0.0 (22) \\ 
%         Cross-source & 40.0 (15) & 40.0 (15) & 35.7 (14) & 62.5 (8) & 45.5 (11) & 26.7 (15) & 26.7 (15) & 28.6 (14) & 25.0 (8) & 18.2 (11) \\ 
%         Provenance & 30.8 (13) & 50.0 (12) & 58.3 (12) & 66.7 (9) & 44.4 (9) & 0.0 (13) & 0.0 (12) & 0.0 (12) & 0.0 (9) & 0.0 (9) \\ 
%         Coverage-aware & 50.0 (6) & 0.0 (2) & \textemdash & \textemdash & \textemdash & 0.0 (6) & 0.0 (2) & \textemdash & \textemdash & \textemdash \\ 
%         LoD3 showcase & 100.0 (4) & \textemdash & \textemdash & \textemdash & 0.0 (1) & 100.0 (4) & \textemdash & \textemdash & \textemdash & 0.0 (1) \\ 
%         Infeasible & 90.0 (10) & 90.0 (10) & 90.0 (10) & 90.0 (10) & 90.0 (10) & 0.0 (10) & 0.0 (10) & 0.0 (10) & 0.0 (10) & 0.0 (10) \\ \midrule 
%         \textbf{Feasible EX} & \textbf{60.7 (89)} & \textbf{56.1 (82)} & \textbf{61.5 (65)} & \textbf{69.2 (39)} & \textbf{53.7 (67)} & \textbf{19.1 (89)} & \textbf{14.6 (82)} & \textbf{9.2 (65)} & \textbf{10.3 (39)} & \textbf{6.0 (67)} \\ \bottomrule 
%     \end{tabular}
%     \par\smallskip\raggedright\footnotesize
%     \textemdash{} denotes that no admitted question exists for that city and category. Category definitions follow \Cref{tab:categories}. The \emph{Infeasible} row scores correct refusals.
% \end{table*}

% \section{Representation Learning}

% \begin{center}
%   \includegraphics[width=\linewidth]{figures/fig_prov_tsne.pdf}
%   \captionof{figure}{2D t-SNE of Hamburg building embeddings from the T3 encoders. Columns show roof type, OSM coverage, and roof-material coverage, where gray marks buildings the ML layer never covered. Compared with the agnostic encoder, the aware encoder more clearly isolates OSM-covered buildings and yields a more structured space.}
%   \label{fig:prov-tsne}
% \end{center}

% \section{Representation-Learning Benchmark: Supplementary}\label{app:repL}

% \Cref{fig:prov-tsne} shows the 2D t-SNE of Hamburg building embeddings from the provenance-agnostic and provenance-aware T3 matching encoders, the qualitative view referenced in \Cref{sec:repL:results}. Table~\ref{tab:repL-crosscity-roof} reports the T2 roof-type cross-city transfer, limited to Hamburg and Helsinki, the only two cities that carry the label.

% \begin{table}[ht!]
%   \centering
%   \caption{T2 roof-type cross-city transfer, leave-one-city-out on the only valid two-city pair (macro-F1 on the merged 3-class label, mean over \num{3} seeds, best per row in bold).} 
%   \label{tab:repL-crosscity-roof}
%   \setlength{\tabcolsep}{4pt}
%   \begin{tabular}{@{}lcc@{}}
%     \toprule
%     Held-out (roof type) & agnostic GNN & aware GNN\\
%     \midrule
%     Hamburg  & 0.203 & \textbf{0.208}\\
%     Helsinki & \textbf{0.336} & \textbf{0.337}\\
%     \bottomrule
%   \end{tabular}
% \end{table}

% \begin{figure*}[t]
%   \centering
%   \includegraphics[width=\linewidth]{figures/fig_prov_tsne.pdf}
%     \caption{2D t-SNE of Hamburg building embeddings from the T3 matching encoders. Columns show roof type, OSM coverage, and roof-material coverage, where gray marks buildings the ML layer never covered. Compared with the agnostic encoder, the aware encoder more clearly isolates OSM-covered buildings and yields a more structured space.}
%   \label{fig:prov-tsne}
% \end{figure*}

\section{Representation-Learning Benchmark: Supplementary}\label{app:repL}

\Cref{fig:prov-tsne} shows the 2D t-SNE of Hamburg building embeddings from the provenance-agnostic and provenance-aware T3 matching encoders, the qualitative view referenced in \Cref{sec:repL:results}. Table~\ref{tab:repL-crosscity-roof} reports the T2 roof-type cross-city transfer, limited to Hamburg and Helsinki, the only two cities that carry the label.

\begin{table}[ht!]
  \centering
  \caption{T2 roof-type cross-city transfer, leave-one-city-out on the only valid two-city pair (macro-F1 on the merged 3-class label, mean over \num{3} seeds, best per row in bold).}
  \label{tab:repL-crosscity-roof}
  \setlength{\tabcolsep}{4pt}
  \begin{tabular}{@{}lcc@{}}
    \toprule
    Held-out (roof type) & agnostic GNN & aware GNN\\
    \midrule
    Hamburg  & 0.203 & \textbf{0.208}\\
    Helsinki & \textbf{0.336} & \textbf{0.337}\\
    \bottomrule
  \end{tabular}
\end{table}

\begin{figure*}[t]
  \centering
  \includegraphics[width=\linewidth]{figures/fig_prov_tsne.pdf}
    \caption{2D t-SNE of Hamburg building embeddings from the T3 matching encoders. Columns show roof type, OSM coverage, and roof-material coverage, where gray marks buildings the ML layer never covered. Compared with the agnostic encoder, the aware encoder more clearly isolates OSM-covered buildings and yields a more structured space.}
  \label{fig:prov-tsne}
\end{figure*}

\end{document}
